% !TEX program = xelatex
\documentclass[11pt,a4paper]{ctexart}

\usepackage[margin=2.3cm]{geometry}
\usepackage{amsmath,amssymb,amsthm}
\usepackage{booktabs}
\usepackage{longtable}
\usepackage{array}
\usepackage{multirow}
\usepackage{graphicx}
\usepackage{caption}
\usepackage{enumitem}
\usepackage[hidelinks]{hyperref}
\usepackage{xcolor}
\usepackage{titlesec}

\hypersetup{colorlinks=true, linkcolor=black, urlcolor=blue, citecolor=black}
\captionsetup{font=small,labelfont=bf}
\setlength{\parskip}{0.35em}
\setlength{\parindent}{2em}
\setlength{\emergencystretch}{3em}
\setlength{\tabcolsep}{4pt}
\setlist{nosep}
\renewcommand{\arraystretch}{1.15}

\newcommand{\confhigh}{\textbf{高}}
\newcommand{\confmid}{中}
\newcommand{\conflow}{\textit{低}}

\title{\bfseries 趋势跟踪与均值回归：\\
\large 同一自相关结构的两侧剪影\\[2pt]
\large ——本质异同、执行分野与统一口径实证的全景比较}
\author{叨客宁宁}
\date{}

\begin{document}

\maketitle

\begin{abstract}
\noindent 趋势跟踪与均值回归被普遍表述为两种相反的策略：前者追涨杀跌，后者买跌卖涨。
本综述主张一个可能令人不适的论断——这种"相反"是表层修辞，二者在数学上共享同一份原料，
即收益序列的自相关结构；它们不是两套世界观，而是同一预测信号的两种符号选择。
为支撑该论断，本文做了三件事。
第一，把两条各自独立生长近百年的文献脉络并置编年：从 Bachelier(1900) 与
Uhlenbeck--Ornstein(1930) 的数学源头，经 Cowles(1933)、Kendall(1953)、Fama(1970) 的效率市场检验，
到 De Bondt--Thaler(1985) 的过度反应、Jegadeesh--Titman(1993) 的截面动量、
Gatev 等(2006) 的配对交易、Moskowitz--Ooi--Pedersen(2012) 的时间序列动量、
Hurst--Ooi--Pedersen(2017) 的百年回溯，直至 2020 年代的机器学习再挖掘。
第二，给出"统一—对偶"的分析框架：用一个含符号的回归强度参数、
方差比 $VR(q)$ 与半衰期/趋势持续期把两策略装入同一数学容器，
再逐条论证七组本质差异（平稳性假设、收益偏度、时间尺度、风险来源、
容量与换手、失效模式、入场门槛），并解释二者为何长期共存。
第三，用一个统一口径的对照回测检验上述论断：在中国期货主力连续与 A 股两个资产池上，
令趋势与反转共享同一资产池、同一调仓频率、同一成本模型，
唯一区别是信号的符号；结果支持"时间尺度决定归属"的核心预测。
全文收录 $80+$ 条逐条核验的来源（DOI/arXiv/官方仓库），附证据表、统一比较矩阵、
结论置信度分级、评测数据污染与不可比设定审计、可复现实验路线图。
本文对争议结论如实标注反方证据与边界，未获全文复核的具体数字一律标明，不作臆测。
\end{abstract}

\tableofcontents
\newpage

\section{引言：一个被修辞掩盖的问题}\label{sec:intro}

\subsection{问题的提法}

交易策略的通俗分类把趋势跟踪与均值回归摆在对立的两端：前者顺势加仓、追涨杀跌，信奉“趋势会延续”；后者逆势布局、买跌卖涨，信奉“偏离终将修复”。坊间叙述往往把这种对立上升为世界观之争——一种相信市场会走远，另一种相信市场会回头。

这类叙述有一个不易察觉的问题：它把\emph{交易动作}的差异误当成了\emph{信息来源}的差异。追涨与买跌确实是相反的动手方式，但两者据以决策的原料完全相同——都是资产过去一段时间的价格变化。趋势策略把这段历史当作延续的证据，反转策略把同一段历史当作过度的证据。二者从未使用不同的信息集，只是对同一份信息给出了相反的读数。

由此可以给出本文的核心论断：

\begin{quote}
\itshape
趋势跟踪与均值回归不是两套独立的世界观，而是同一个预测信号的两种符号选择。
把它们区分开的，既不是信息集、也不是数学工具，而是关于\emph{收益自相关符号}的一个假设。
\end{quote}

这个论断如果成立，会带来一连串可检验的推论：两条文献脉络应当在历史上共享源头；同一份数据的自相关函数应当在不同时间尺度上同时呈现正负两种符号；两个策略的收益分布应当在偏度上呈现系统性的镜像；而最直接的一条是——用同一资产池、同一成本模型、仅仅翻转信号符号，应当能复现出两者的绩效分野。

本文的任务就是把这个论断拆开、论证、并放到数据上检验。

\subsection{为什么这件事值得认真做}

否认“两策略本质对立”的说法，常被误解为一种和解式的修辞。恰恰相反：把它讲清楚是有实际后果的。

其一，\textbf{它改变了参数选择的方式}。如果两策略共享同一信号，那么决定用哪一个的关键变量就不是“你相信哪种市场观”，而是\emph{你打算在哪个时间尺度上持仓}。Bachelier 以来的收益序列自相关研究反复指向同一个经验规律：短周期与长周期呈现负自相关，中周期呈现正自相关。这意味着趋势与反转的归属由持有期决定，而不是由信念决定。把选策略的问题转化为选时间尺度的问题，是一个更接近可操作层面的提法。

其二，\textbf{它改写了风险预算的对象}。两者的收益分布并非对称的镜像。趋势策略的收益右偏、伴随罕见的深度左尾（动量崩溃），本质上是持有一个隐含的看涨期权结构；反转策略的收益左偏、依赖价格收敛，本质上是承担了卖方期权式的尾部风险\cite{funghsieh2001,danielmoskowitz2016}。如果只比较夏普比率而忽略偏度，就会把一个“平时赚钱、偶尔跳崖”的策略与一个“常态波动、极端时刻被资助”的策略当成同类。执行层面的止损、杠杆、再加仓规则，必须建立在对这种偏度差异的认识之上。

其三，\textbf{它暴露了共同的脆弱点}。两策略的发现在方法论上面对同一类质疑：多重检验与数据窥探。Sullivan 等(1999)对 7846 条技术规则的检验、White(2000)的 Reality Check、Harvey 等(2016)提出的更高 $t$ 值门槛，指向的是同一个问题——在一个被反复搜索的数据集上，“显著”是廉价的\cite{sullivan1999,white2000,hlz2016}。趋势与反转都难逃这一审查；把二者放入同一审计框架，比分别辩护更有说服力。

\subsection{比较的判据}

要比较两个策略，先要定义“比什么”。若不加限定地比较，结论会轻易退化为不可比的数字竞赛：一个是月频、低换手、跨资产的组合，另一个是日频、高换手、单资产的价差交易，说“前者年化更高”没有意义。

本文确立六个维度的比较判据，后文各章围绕它们展开：

\begin{enumerate}[leftmargin=2.2em]
  \item \textbf{信息与数学结构}：用什么变量预测什么量？预测函数的形状与符号如何？这是“本质”层面。
  \item \textbf{统计对象}：策略押注的究竟是哪一个可检验量——自相关系数、方差比、协整关系、还是半衰期？这决定了策略是否“先有入场券”。
  \item \textbf{收益分布}：偏度、峰度、以及极端市场中的表现（凸性 vs 凹性）。
  \item \textbf{执行代价}：换手量级、对交易成本的敏感度、容量与拥挤。
  \item \textbf{失效模式}：什么情况下会亏钱？亏钱是渐进的还是突发的？失败是可止损的还是“越止损越糟”的？
  \item \textbf{证据质量}：支持各自的核心证据在数据窥探、样本外衰减、发表后衰减上表现如何。
\end{enumerate}

\subsection{研究方法与证据纪律}

本文属于证据型综述，方法上遵循三条纪律。

\textbf{一是全部引用可回溯至原始来源}。文献的题名、作者、年份、期刊与 DOI 均通过 OpenAlex、Crossref、arXiv 官方接口或出版社页面逐一核对，不以记忆填写。凡接口返回与记忆冲突者，以接口为准。少数无法通过公开接口确证的经典（如 Alexander 1961 的早期工作），明确标注为待复核，并在正文中避免依赖其具体数值。

\textbf{二是区分“机制解释”与“实证事实”}。行为学解释（反应不足、过度反应）与风险补偿解释（时变风险溢价）在文献中并存且各有支持者，本文不做仲裁，而是并列呈现，并在结论置信度一章逐条给出分级\cite{cochrane2011,bsv1998,dhs1998,hongstein1999}。

\textbf{三是自备一组统一口径的对照实证}。文献中的绩效数字来自不同样本、不同成本假设、不同资产池，横向比较的可靠性有限。为检验“符号翻转”这一核心预测，本文在中国期货主力连续合约与 A 股两个资产池上重建了一套对照回测：两策略共享同一资产池、同一调仓频率、同一成本模型与同一评估指标，唯一区别是信号符号（详见第~\ref{sec:empirical}~节）。这组实证不宣称发现新的alpha，它的作用是\emph{把文献论断放到一个可控的变量对照下}。

\subsection{边界与结构}

本文的比较限于\emph{可交易的收益自相关策略}。宏观计量意义上的均值回归（如购买力平价的长期收敛检验）与非交易性的价格发现研究，仅在涉及机制解释时引用，不作为策略比较的对象。此外，本文讨论的市场证据以发达市场文献为主，中国市场部分主要依赖本文自建的对照实证与国内可得数据，样本长度与制度差异构成外推边界，将在第~\ref{sec:audit}~节与结论中明确。

全文结构如下。第~\ref{sec:genealogy}~节并置两条脉络的编年史；第~\ref{sec:unity}~节给出统一框架；第~\ref{sec:duality}~节论证七组本质差异；第~\ref{sec:matrix}~节用统一比较矩阵汇总证据；第~\ref{sec:mechanism}~节对照机制解释；第~\ref{sec:execution}~节落到执行差异；第~\ref{sec:empirical}~节是本文自建的对照实证；第~\ref{sec:audit}~节做污染与不可比审计；此后依次是结论置信度分级、可复现实验路线图、以及研究现状与遗留课题。

\section{双线谱系：两条脉络的编年并置}\label{sec:genealogy}

\subsection{并置的理由}

把两条文献脉络放在同一张时间轴上，是检验引言论断的第一个动作。若“两策略共享同一源头”成立，那么在早期（1920--1980）的文献中，我们应当看到同一批作者用同一套统计工具同时讨论两个方向，而不是两条互不相干的溪流。事实正是如此。

\subsection{共同的数学源头（1900--1930）}

现代金融时间序列的数学起点通常追溯到 Bachelier(1900) 的博士论文，他把价格变动建模为布朗运动\cite{bachelier1900}。这一建模选择本身就规定了后续争论的形态：一旦价格服从带独立增量的随机游走，收益就既无自相关、也无均值回复，任何基于历史价格的策略在理论上都无法获利。此后近一个世纪的实证研究，本质上是不断检验这个零假设在多大程度上被违反。

Uhlenbeck 与 Ornstein(1930) 提出的过程给出了“违反”的最简便形式：一个带线性回复力的扩散过程，偏离均值越远，被拉回的倾向越强\cite{uo1930}。这个后来被称为 OU 过程的对象，在数学上是均值回归策略的母体——它同时规定了回归强度（半衰期）与平稳分布，把“价格会回来”这句话变成了可估参数。它与随机游走的差别只有一个漂移项：$dX_t=\theta(\mu-X_t)dt+\sigma dW_t$ 中的 $\theta$ 是正、是零、还是负。\textbf{这一个符号，就是本文全篇要反复回到的地方。}

\subsection{效率市场的检验期（1930--1970）：同一位作者，两个方向}

1930 至 1970 年是两条脉络共用工具的时期。

Cowles(1933) 与 Cowles \& Jones(1937) 检验专业预测者和技术交易规则能否跑赢市场，方法是对收益序列做自相关与游程检验，结论偏负面\cite{cowles1933,cowlesjones1937}。这项工作的重要性不在于结论，而在于它确立了“用序列相关检验可预测性”的方法论范式——而这正是后来动量与反转研究共同使用的工具。同一路径上，Fama 与 Blume(1966) 对“过滤规则”（filter rules）的检验给出了更直接的交易含义：在扣除交易成本后，机械的顺势规则无法战胜买入持有\cite{famablume1966}。这一结论在三十年后被 Brock 等(1992) 的移动平均与突破规则检验所动摇\cite{brock1992}，而后者又被 Sullivan 等(1999) 用数据窥探校正部分推翻\cite{sullivan1999}——这条“发现—质疑—再质疑”的链条是趋势研究中最典型的方法论演进样本。

Kendall(1953) 对英国商品与股票价格的时间序列分析，进一步支持了价格近似随机游走的判断\cite{kendall1953}。Fama(1970) 的效率市场综述把上述结果收拢为一个可检验的框架：弱式有效意味着历史价格对预测未来收益无增量信息\cite{fama1970}。Levy(1967) 的相对强弱选股准则则提供了一个早期的经验反例：按过去涨幅排序的组合在随后一段时间表现更优\cite{levy1967}——这是被后来的截面动量文献反复引用的先声。

这段历史的微妙之处在于：\emph{它同时是动量与反转的先验否定}。如果市场是弱式有效的，“追涨”与“买跌”都无效；两者的合法性必须等到效率市场的经验反例足够多之后才被重新讨论。

\subsection{分岔点：1985--1993}

真正的分岔发生在 1980 年代中期到 1990 年代初，而且两条脉络几乎同时破土。

\textbf{反转一侧}。Shiller(1981) 与 Summers(1986) 从价格波动过大这一现象出发，论证价格可以长期偏离基本面\cite{shiller1981,summers1986}。De Bondt \& Thaler(1985) 给出了最直接的交易含义：以 1926--1982 年数据构造形成期 36 个月的输家与赢家组合，发现输家随后跑赢赢家，被解释为过度反应后的回归\cite{debondt1985}。Fama \& French(1988) 用回归分离股价的永久与暂时成分，发现长周期负自相关，为“均值回归在长尺度上存在”提供了更严格的计量证据\cite{famafrench1988}。Poterba \& Summers(1988) 以方差比检验在多个国家与更长的历史样本上得到类似结论\cite{poterbsummers1988}。这一路径随后在国际样本上被反复检验：Richards(1997) 检验各国股指的赢家——输家反转能否用风险解释\cite{richards1997}，Balvers 等(2000) 把参数化反向策略推广到 18 个国家指数并给出可交易的截面对冲方案\cite{balvers2000}。

\textbf{动量一侧}。几乎同时，Jegadeesh \& Titman(1993) 报告了方向相反的结果：以 3--12 个月为形成期、3--12 个月为持有期，买入赢家卖出输家的组合在随后一年获得显著超额收益\cite{jegadeesh1993}。这与长周期反转并不矛盾——两者作用在不同的时间尺度上。Lehmann(1990) 与 Jegadeesh(1990) 则报告了更短尺度（周频、月频）上的负自相关，即短期反转\cite{lehmann1990,jegadeesh1990}。Chan 等(1996) 进一步检验动量的收益来源，指出盈利惯性与信息扩散无法完全解释，价格动量本身独立有效\cite{chan1996}。

至此，经验事实已经显露出它的真实形状：\textbf{收益自相关随尺度变化符号}。短尺度为负、中尺度为正、长尺度为负。趋势与反转不是两个事实，而是同一个自相关函数在不同横坐标上的取值。

\subsection{两条脉络的独立生长（1990--2010）}

意识到“尺度决定归属”之后，两条脉络并未合流，反而各自向纵深发展，形成两套几乎不交叉的技术语言。

\textbf{反转一侧走向统计套利}。Lo \& MacKinlay(1990) 澄清了反转利润的来源问题：许多被报告的反转收益来自股票之间的交叉自相关（大市值滞后带动小市值），而非个股自身的过度反应\cite{lomackinlay1990}。这一批评把反转研究推向更精细的技术路线。Engle \& Granger(1987) 的协整框架为此提供了计量基础\cite{englegranger1987}：只有当价差平稳（协整）时，“价差会回归”才是可检验的命题。此后配对交易成为主力范式——Gatev 等(2006) 用 1962--2002 年 CRSP 日数据，以最小化标准化价差平方和的距离法选对，报告年化超额收益约 11\%\cite{gatev2006}；Avellaneda \& Lee(2010) 用因子与主成分剥离残差构建统计套利组合\cite{avellaneda2010}；Elliott 等(2005) 把 Kalman 滤波引入配对价差的在线估计\cite{elliott2005}。Do \& Faff(2012) 随后发现配对交易利润随时间显著衰减，早期高收益不能外推\cite{dofaff2012}。

\textbf{动量一侧走向跨资产因子}。Rouwenhorst(1998) 在 12 个欧洲市场确认动量的普遍性\cite{rouwenhorst1998}；Moskowitz \& Grinblatt(1999) 发现行业动量解释了个股动量的大部分\cite{moskowitzgrinblatt1999}；Jegadeesh \& Titman(2001) 用 1990 年代的数据确认动量持续存在，并倾向于行为解释\cite{jegadeesh2001}。Fung \& Hsieh(1997,2001) 把 CTA 的收益刻画为“回溯跨式”式的凸性暴露，使趋势跟踪第一次获得了可因子化的表述\cite{funghsieh1997,funghsieh2001}。Moskowitz、Ooi \& Pedersen(2012) 提出的时间序列动量把趋势跟踪彻底学术化：在 58 个期货合约上，过去 12 个月超额收益正向预测未来收益，且几乎每个合约都盈利\cite{mop2012}。Asness 等(2013) 进一步证明价值与动量在八类资产上普遍存在且负相关\cite{asness2013}。

两套语言的差异此时已经固化：反转研究的关键词是\emph{协整、平稳性、残差、阈值}；动量研究的关键词是\emph{形成期、持有期、波动率缩放、横截面排序}。这一点将直接导致后文讨论的“入场门槛”差异。

跨资产层面，这条脉络还向外汇与商品延伸：Menkhoff 等(2012) 系统检验外汇动量策略，发现收益高但换手与交易成本显著\cite{menkhoff2012}；Erb \& Harvey(2006) 在商品期货上检验展期收益与动量择时的作用\cite{erbharvey2006}。技术分析的学术化则以 Lo 等(2000) 为代表，该文用计算算法与统计推断为形态识别提供了可检验的实现\cite{lo2000}。Baltas 与 Kosowski(2020) 后来把时间序列动量的收益拆解为波动率估计、交易规则与两两相关性三部分，是该领域最系统的技术性梳理\cite{baltaskosowski2020}。

\subsection{共同的批判期（2015--2020）}

数据窥探的批评其实早有预备：Sullivan 等(1999) 对 7846 条技术交易规则施加 bootstrap 校正，发现样本外超额收益不再显著\cite{sullivan1999}；White(2000) 提出 Reality Check\cite{white2000}。但真正的转折发生在 2015 年之后：Harvey 等(2016) 统计 316 个已发表因子，论证常规 $t>2$ 门槛过低，新因子应达到 $t>3$ 左右\cite{hlz2016}；McLean \& Pontiff(2016) 量化了发表后衰减——组合收益在样本外下降 26\%，发表后累计下降约 58\%\cite{mcleanpontiff2016}。

在这一审查下，两条脉络各自被点名。动量一侧，Huang 等(2020) 逐资产回归后质疑时间序列动量的稳健性，指出其表现与“按历史样本均值投资”的简单策略几乎无差别\cite{huang2020}；Daniel \& Moskowitz(2016) 则揭示了动量的深度左尾——在恐慌状态下的动量崩溃\cite{danielmoskowitz2016}。反转一侧，Do \& Faff(2012) 的利润衰减证据\cite{dofaff2012}、以及 2007 年 8 月量化危机的经验研究\cite{khandanilo2011}，指向反转策略在拥挤与去杠杆环境下的集体失败。

\subsection{百年尺度的证据与当代回归（2017--至今）}

Hurst、Ooi \& Pedersen(2017) 手工收集 1880 年以来的数据，把趋势跟踪的回溯拉到 138 年，发现每个十年均录得正收益，且在 60/40 组合的十次最大回撤中有八次为正\cite{hop2017}。这是“趋势=稳健因子”叙事最常被引用的证据，也构成对“样本期依赖”批评的直接回应。

与此同时，两策略在方法上的重新接近开始出现。Wood 等(2021) 的“慢动量加快反转”用深度学习与变点检测把两个方向融合进同一个交易规则\cite{wood2021}。这一尝试的动机恰恰是本文的核心论断：既然二者共享信号，那么把它们当作\emph{同一信号在不同尺度上的两个分量}来建模，在工程上是自然的。

\subsection{编年并置说明了什么}

把两条脉络放在一起看，有三个事实浮现。

\textbf{其一，源头是共同的}。OU 过程同时孕育了两者；Bachelier 的随机游走同时否定了两者。
\textbf{其二，分岔点在时间尺度}。1985 年前后同时出现的过度反应证据与动量证据，其差异并不在方法或数据，而在形成期与持有期的长度。
\textbf{其三，批判的降临是同步的}。数据窥探、发表后衰减、拥挤与尾部风险，对两者构成同构的威胁。

这三个事实合起来，支撑了引言中的论断：讲“两策略对立”，是把同一现象的两个切片读成了两种现象。下一节把这个直觉写成数学。

\subsection{一条常被忽略的支线：波动率的均值回归}

在价格方向的争论之外，还有一条几乎平行的支线，它同样古老，却很少被放进“趋势 vs 反转”的比较里。

Engle(1982) 提出的 ARCH 族与 Bollerslev(1986) 的 GARCH 把“波动率自身会回归其长期水平”变成可估计的命题\cite{engle1982,bollerslev1986}。这条支线的经验规律同样清晰：波动率具有显著的自相关（聚类），且在足够长的尺度上回归均值——这正是“波动率均值回归”策略族的基础。此后 Carr 与 Wu(2009) 用指数期权定价出方差风险溢价\cite{carrwu2009}，Bollerslev 等(2009) 发现该溢价能显著预测短期股票收益\cite{bollerslev2009}。

它为什么重要？因为它揭示了一件事：\emph{“均值回归”这个标签下其实装着不止一种载体}。价格的均值回归（$\theta>0$ 的 OU）与波动率的均值回归（GARCH 的条件方差收敛）在数学上不同构——前者的状态变量是一阶矩，后者是二阶矩。这就带来一个操作上的分岔：用价格反转的逻辑去做波动率交易会得到错误结论（波动率不“偏离均值后必然回归”到同一速度），反之亦然。把两者混淆，是实务中一类常见错误的来源。

这条支线也提供了一个对照视角：价格方向上的“趋势 vs 反转”之争，在波动率维度上并不存在对称的争论——因为波动率均值回归的机制（风险溢价的期限结构）比价格方向的机制更容易被数据锁定。

\subsection{被后来文献反复重写的三个节点}

把这条编年史压缩，有三个节点经历了最多的重写，它们分别对应本文后面三章的核心争论。

\textbf{技术规则是否有效（1966--1999）}：Fama 与 Blume(1966) 判定过滤规则无效\cite{famablume1966}，Brock 等(1992) 推翻这一判定\cite{brock1992}，Sullivan 等(1999) 再用数据窥探校正把结论拉回“样本外不显著”\cite{sullivan1999}。同一问题被改了三次，留下的教训是：\emph{“技术规则有效吗”这个问法本身不完整，必须附带“在多大的搜索空间内、用哪种多重检验校正”}。

\textbf{长周期反转是否真实（1985--1993）}：De Bondt--Thaler(1985)、Fama--French(1988)、Poterba--Summers(1988) 给出肯定\cite{debondt1985,famafrench1988,poterbsummers1988}，Richardson(1993) 系统指出这些证据对小样本与检验设定高度敏感\cite{richardson1993}。争论至今没有闭合，但它把“检验功效”概念引入了金融时间序列研究。

\textbf{时间序列动量是否稳健（2012--2020）}：Moskowitz 等(2012) 的 58 合约证据\cite{mop2012}与 Hurst 等(2017) 的百年样本\cite{hop2017}，对上 Huang 等(2020) 的逐资产再检验\cite{huang2020}与 Pitkäjärvi 等(2020) 的跨资产质疑\cite{pitkajarvi2020}。分歧的实质见第~\ref{sec:confidence}~节：它不是“有效或无效”的二分，而是“收益来自预测还是来自缩放”的分解问题。

三个节点的共同形态是同一件事：\emph{一个看似干净的实证结论，在被追问“搜索空间多大、样本多长、成本多少”之后，会逐步失去它的确定性}。这正是本文设立可信度分级的直接理由。

\section{本质的统一：一个信号，两个符号}\label{sec:unity}

\subsection{本质问题}

要判断两个策略是否“本质上相同”，需要先回答一个更基本的问题：\emph{一个交易策略究竟在赌什么}？

把任何策略的逐期损益写开，就是持仓权重与随后实现收益的内积。取期望，策略的期望收益等于
\begin{equation}
\mathbb{E}\!\left[\sum_t w_t\, r_{t+1}\right]
=\sum_t \mathrm{Cov}\!\left(w_t,\; r_{t+1}\right),
\end{equation}
其中 $r_{t+1}$ 是下一期收益，$w_t$ 是此刻的持仓（在无风险利率为零、且权重由可用信息决定的设定下，均值项可配对抵消）。这个分解给出了策略赚钱的唯一合法途径：\textbf{持仓必须与随后收益正相关}。所谓“有 alpha”，在数学上就是这个协方差为正。

于是不同策略族的区别，只能落在 $w_t$ 与 $r_{t+1}$ 各自“看什么”上：

\begin{itemize}[leftmargin=2.2em]
  \item 若 $w_t$ 取决于价格的\emph{绝对水平}或\emph{相对估值}（如市净率、展期收益率），那是价值/利差类策略——它们赌的是一阶量（水平）。
  \item 若 $w_t$ 取决于价格的\emph{历史变化}（过去一段时间的收益），那是本文讨论的对象——它们赌的是二阶量（自相关）。
\end{itemize}

这一划分是本文全部论证的起点：\textbf{趋势跟踪与均值回归同属后者，它们赌的都是符号，而不是水平。}两者都不问“这个资产现在贵不贵”，只问“它最近在往哪儿走，以及这个走势会不会持续”。

\subsection{第一性原理：唯一的自由度是符号}

设资产在 $t$ 期的收益为 $r_t$。定义\emph{信号}为过去 $L$ 期的累计收益
\begin{equation}
s_t(L)=\sum_{k=0}^{L-1} r_{t-k}.
\end{equation}
两类策略的持仓都可以写成同一个响应函数
\begin{equation}
w_t \;=\; \pm\, g\!\left(\frac{s_t(L)}{\sigma_t}\right),
\qquad g \text{ 单调不减},\ \sigma_t \text{ 为波动率尺度},
\end{equation}
其中取 $+$、$g$ 递增即趋势跟踪；取 $-$、$g$ 递减即均值回归。信号 $s_t(L)$、尺度 $\sigma_t$、窗口 $L$ 三个要素被完全共享；\textbf{唯一的自由度是那个符号}。

这个提法并不抽象，它精确对应了两类文献里的标准构造。趋势一侧，Moskowitz、Ooi \& Pedersen(2012) 的时间序列动量取
$w_t \propto \mathrm{sign}(s_t(L))/\sigma_t$，$L=12$ 个月\cite{mop2012}；截面动量则把 $s_t$ 在资产间排序后做多前列、做空后列\cite{jegadeesh1993}。反转一侧，配对交易取价差 $X_t=P^A_t-\beta P^B_t$ 对均值的\emph{负}偏离作为建仓依据，即 $w_t=-\mathrm{sign}(X_t-\mu)/\sigma$ 在阈值带外的版本\cite{gatev2006,elliott2005}；短期反转则直接用 $s_t(L)$ 的负号，$L$ 取一到四周\cite{lehmann1990,jegadeesh1990}。两组构造在符号之外逐项对应。

\subsection{同构论证：符号与回归强度是同一件事}

上面的统一仍停留在“构造形式相同”的层面。更强的同构在于：那个符号，可以严格地解释为均值回归速度的符号。

考虑对数价格 $X_t=\log P_t$ 服从 OU 过程
\begin{equation}
dX_t=\theta(\mu-X_t)\,dt+\sigma\,dW_t .
\end{equation}
$\theta>0$ 时 $X_t$ 围绕 $\mu$ 往复，偏离被拉回——这正是均值回归的连续时间载体\cite{uo1930,vasicek1977}。离散化后得到 AR(1)：
\begin{equation}
X_{t+\Delta}-\mu=e^{-\theta\Delta}\left(X_t-\mu\right)+\varepsilon_t ,
\end{equation}
自回归系数 $\phi=e^{-\theta\Delta}\in(0,1)$，自相关 $\rho_k=\phi^{k}$ 恒为正、随时间衰减，半衰期 $H=\ln 2/\theta$。

现在把 $\theta$ 的符号翻过来。若 $\theta<0$，漂移项变为 $-|\theta|(X_t-\mu)$，偏离均值不是被拉回而是被放大，价格呈指数发散——这是“趋势”的连续时间载体。此时离散系数 $\phi=e^{-|\theta|\Delta}>1$，序列非平稳，长期方差随持有期以超线性速度增长。用方差比的语言说，$\theta>0$ 对应 $VR(q)<1$，$\theta<0$ 对应 $VR(q)>1$。

\textbf{这就是本文主张的数学内核：趋势与均值回归是同一个 OU 族在两个参数半平面上的两个实例。}把 $\theta$ 从正推到负，策略从“买跌卖涨”连续地变成“追涨杀跌”。存在 $\theta=0$ 的分界，那里是随机游走，两种策略都无从获利。

这个同构带来一个立刻可用的推论：\emph{任何为均值回归设计的工具，把符号翻转后原则上都可用于趋势}。半衰期 $H$ 对应趋势的持续期 $1/|\theta|$；OU 的最优阈值带对应趋势的入场/离场带；协整关系对应“跨资产的同步趋势”。文献中“慢动量加快反转”的融合策略之所以可能，正是因为这个同构\cite{wood2021}。

\subsection{统一的统计量：方差比与自相关函数}

若符号统一了“怎么下注”，还需要一个量来刻画“事实站在哪一边”。这个量是收益的自相关结构。

设 $r_t$ 平稳，$\rho_k=\mathrm{Corr}(r_t,r_{t-k})$。$q$ 期收益的方差与 $\rho_k$ 的关系由方差比给出\cite{lomackinlay1988,campbell1997}：
\begin{equation}
VR(q)=\frac{\mathrm{Var}\!\left(\sum_{k=0}^{q-1}r_{t-k}\right)}{q\,\mathrm{Var}(r_t)}
=1+2\sum_{k=1}^{q-1}\left(1-\frac{k}{q}\right)\rho_k .
\end{equation}
$VR(q)>1$ 意味着多期收益的方差超过线性增长，即正自相关占优——这是趋势策略的生存条件；$VR(q)<1$ 意味着负自相关占优——这是反转策略的生存条件。

于是“选哪个策略”被改写为一个可以用数据回答的问题：\textbf{估出 $\rho_k$（或 $VR(q)$），看它在目标持有期上取什么符号。}

这一改写清除了很多无谓的争论。当有人问“趋势有效还是反转有效”，正确的问题是“在哪个 $q$ 上”。文献给出的答案出奇地一致——见下节。

\subsection{时间尺度不可分离：同一个自相关函数的三个区间}

把 $\rho_k$ 作为 $k$ 的函数画出来，几十年文献积累的形状是一条起伏的曲线：

\begin{itemize}[leftmargin=2.2em]
  \item \textbf{短尺度（日到周，$k\approx1$--$5$）}：$\rho_k<0$。Lehmann(1990) 的周频反转、Jegadeesh(1990) 的月频负自相关都在此区间\cite{lehmann1990,jegadeesh1990}。
  \item \textbf{中尺度（一到十二个月，$k\approx20$--$250$）}：$\rho_k>0$。Jegadeesh \& Titman(1993) 的截面动量与 Moskowitz 等(2012) 的时间序列动量在此区间\cite{jegadeesh1993,mop2012}。
  \item \textbf{长尺度（三到五年，$k\approx750$--$1250$）}：$\rho_k<0$。De Bondt \& Thaler(1985)、Fama \& French(1988)、Poterba \& Summers(1988) 的长期反转证据在此区间\cite{debondt1985,famafrench1988,poterbsummers1988}。
\end{itemize}

Cutler、Poterba \& Summers(1991) 的“投机动力学”把这种多尺度并存讲得最清楚：\emph{短期的动量证据与长期的反转证据同时成立，是因为它们瞄准了同一序列的不同频率成分}\cite{cutler1991}。这带来本文第一个可检验推论：

\begin{quote}
\itshape
若两策略共享同一信号，则在同一资产池上，趋势与反转的相对强弱应当随持有期 $L$ 单调反转——短 $L$ 上反转占优，中 $L$ 上趋势占优，长 $L$ 上又回到反转。
\end{quote}

第~\ref{sec:empirical}~节的对照回测正是为了检验这条推论（见本章末的预告）。

\subsection{统一的信息集与响应函数：为什么不能靠“减少信息”来区分}

还有一个常被忽略的统一性：两策略使用\emph{完全相同的信息集}，且都可以做到对信息量的“无损利用”。

考虑信息集 $\mathcal{F}_t$ 包含至 $t$ 为止的全部历史价格。最优的线性预测子由投影给出：$\mathbb{E}[r_{t+1}\mid\mathcal{F}_t]=\sum_k\beta_k r_{t-k+1}$。趋势策略取 $\sum_k\beta_k>0$ 的一阶近似，反转策略取 $\sum_k\beta_k<0$ 的一阶近似。二者都在用同一个 $\mathcal{F}_t$，差别只在估计出的 $\beta$ 的符号。

这一点有直接的操作含义：\textbf{不存在“因为信息不同所以结论不同”的可能性}。如果有人声称趋势策略看到的是“基本面变化”而反转策略看到的是“情绪过度”，那属于\emph{解释层}的分歧，不改变\emph{信号层}的同构；而解释层的分歧完全可以在同一个 $\rho_k$ 上并存——例如 Hong \& Stein(1999) 的信息渐进扩散模型同时给出了短期反应不足（动量）与长期过度反应（反转）\cite{hongstein1999}，Barberis 等(1998) 的投资者情绪模型也是如此\cite{bsv1998}。一个机制同时生成两个符号，这是本文在机制章要展开的内容。

\subsection{统一框架的边界}

必须诚实指出这个统一框架在哪里失效——否则它就成了过度概括。

\textbf{边界一：协整型反转不完全是时间序列反转。}配对交易的信号是\emph{跨资产的价差}偏离，而非单资产自身的历史收益。二者的可交易条件不同：协整要求两个序列存在长期均衡关系，这引入了趋势策略所没有的\emph{结构约束}。虽然形式上都可以写进响应函数，但“价差平稳”是一个比“自相关为负”更强的假设，也更易被破坏。这一差异将在第~\ref{sec:duality}~节作为第一组对偶展开。

\textbf{边界二：非线性与状态依赖。}符号统一的是\emph{一阶}描述。现实中的策略收益往往含状态依赖的非线性（趋势在低波动趋势市有效、反转在高波动区间市有效），线性自相关框架无法完全覆盖。这类现象需要 regime 切换模型处理\cite{hamilton1989,angtimmermann2012}。

\textbf{边界三：容量与成本改变了“有效”的含义。}$\rho_k$ 显著为负\emph{不等于}反转策略能赚钱。Martin \& Sch\"oneborn(2011) 的结论标题就是答案：\emph{Mean Reversion Pays, but Costs}\cite{martin2011}。统计上的正期望可能在扣除换手成本后归零。统一框架描述的是\emph{毛收益}的来源，能否落地是执行章的内容。

\textbf{边界四：波动率本身是均值回归的，且与价格方向不同构}。GARCH 类模型刻画的条件方差向长期均值回归\cite{bollerslev1986}、方差风险溢价的存在\cite{carrwu2009}与它对短期收益的预测能力\cite{bollerslev2009}，构成“波动率均值回归”这一独立家族。它的载体不是价格而是二阶矩，因此不能简单地用价格的符号翻转来刻画——这也是本文把波动率类策略列为统一框架特例而非同构对象的原因。

修正后的判断是：统一框架的适用范围是“信号层”——它精确地说明了两个策略共享什么、区别在哪；它\emph{不}声称两策略在执行、风险、容量上等价。恰恰相反，正是因为信号层如此接近，执行层的差异才成为决定成败的关键。这正是本文后半部分的重心。

\section{本质的差异：七组对偶}\label{sec:duality}

上一章论证了两策略在\emph{信号层}的同构。但同构不等于同质。本章给出七组对偶，逐条说明在共享同一信号的前提下，两个方向为什么会在假设、分布、机制与约束上分道扬镳。每组对偶按“命题—论证—证据—工程映射—边界”的顺序展开，并在末尾给出可检验的判断。

\subsection{对偶一：平稳性与非平稳性}

\textbf{命题}：均值回归策略要求被交易的对象具备平稳性（或至少存在协整关系）；趋势策略反而依赖非平稳性。两者对“世界是什么样”的假设正好相反。

\textbf{论证}。均值回归的可交易基础是价差 $X_t$ 的平稳性：只有当 $X_t$ 存在一个吸引子（长期均值）时，“偏离终将修复”才是可检验的命题。若 $X_t$ 含单位根，则偏离不会回归，反转策略面临无限累积的亏损。因此配对交易的现代形式把协整检验当作入场券：Engle \& Granger(1987) 的两步法——对两个价格序列做回归、再对残差做单位根检验——是这一门槛的标准实现\cite{englegranger1987}；Dickey--Fuller 检验则是单位根判断的入口\cite{dickeyfuller1979}。

趋势策略的处境完全不同。若价格服从带正漂移的随机游走，趋势信号仍然有效——因为漂移本身提供了方向。Hurst 等(2017) 的百年样本显示趋势跟踪在极端宏观环境下表现最好，本质上是因为那些时期价格序列呈强非平稳（危机后的持续单向运动）\cite{hop2017}。换言之，\textbf{趋势策略恰恰在反转策略最不该交易的时刻获利}。

\textbf{工程映射}。这一差异直接决定了策略研发的第一步动作不同。反转策略的研发流程必须包含\emph{可行性筛选}：先做 ADF/协整检验，通过者才进入交易规则设计；筛选失败的样本被丢弃。趋势策略没有对应的筛选步骤，它的第一步是估计动量统计量（$t$ 值、胜率、夏普），信号弱但仍在。

\textbf{边界}。这个对偶有一个重要的反例：\emph{短期反转}（日频、周频）并不依赖严格的协整关系，它更接近微观结构效应——买卖价差反弹、流动性提供者的补偿。这类反转的“平稳性”来自订单流机制而非基本面的长期均衡，其半衰期以小时或天计。因此，把“反转必先协整”绝对化会漏掉整个高频反转族。准确的表述是：\emph{中长尺度反转依赖协整，短尺度反转依赖微观结构}。

\textbf{可检验判断}：在同一资产池上，反转策略的绩效对“是否通过平稳性筛选”应当高度敏感；趋势策略的绩效则不应因剔除不平稳样本而系统性改善。

\subsection{对偶二：收益偏度——凸性与凹性}

\textbf{命题}：趋势策略的收益右偏、本质上是持有看涨期权式的凸性暴露；反转策略的收益左偏、本质上是卖出期权式的凹性暴露。二者在小概率极端事件上表现相反。

\textbf{论证}。Fung \& Hsieh(2001) 用“回溯跨式”（lookback straddle）期权组合复制趋势跟踪者的收益，指出其收益结构具有明确的凸性：在市场发生大幅单边运动时收益放大，在震荡市中缓慢失血\cite{funghsieh2001}。这一刻画解释了业界常说的“crisis alpha”：趋势策略在股灾、危机中往往正收益，因为危机会制造持续的单向运动。

反转策略的偏度方向相反。它的收益来自价格向均值收敛，在绝大多数时间里稳定收租，但在价格\emph{不}收敛时遭受集中损失。这在数学上等价于卖出价外期权：收益分布是“多次小赢、偶尔巨亏”。Harvey \& Siddique(2000) 的条件偏度定价文献为这类暴露提供了资产定价层面的框架\cite{harveysiddique2000}。

\textbf{证据}。趋势一侧有两个直接的实证：Daniel \& Moskowitz(2016) 记录的动量崩溃——在恐慌状态下（市场下跌后、波动率高企）伴随市场反弹，动量组合遭受罕见但持久的亏损，1927--2013 样本中最剧烈的两段是 1932 年 7--8 月与 2009 年 3--5 月\cite{danielmoskowitz2016}。注意：\emph{动量崩溃是趋势策略的“负偏”证据}——这提醒我们，趋势策略的偏度是“总体右偏、但尾部有深度左尾”，不是纯粹的右偏。

反转一侧的两段历史更直白。1998 年 LTCM 的收敛交易在俄罗斯违约后崩盘，其杠杆与流动性结构把一次“价差暂时扩大”放大为系统性失败\cite{lowenstein2000,jorion2000}；2007 年 8 月的“量化危机”中，因子型反转/统计套利组合被迫去杠杆，Khandani \& Lo(2011) 用因子模拟与 TAQ 交易数据识别出渐进的平仓过程，而非单日事件\cite{khandanilo2011}。两次事件的共同点是：反转策略的亏损不是渐进的，而是由\emph{不收敛+资金约束}共同触发的跳变。

\textbf{工程映射}。偏度差异决定了风险控制的手段不同。趋势策略的风险控制核心是\emph{波动率目标化}：Barroso \& Santa-Clara(2015) 用自身已实现方差缩放动量仓位，把夏普从 0.53 提升到 0.97、超额峰度从 18.2 降到 2.7\cite{barroso2015}；Harvey 等(2018) 在 60 多个资产上系统验证波动率目标改善了风险资产的夏普\cite{harvey2018}。反转策略的风险控制则不能依赖波动率缩放——理由见对偶四——它更依赖\emph{止损与资金冗余}：Leung \& Li(2015) 把止损退出纳入 OU 最优交易框架，正是因为不收敛的尾部风险无法靠仓位缩放消除\cite{leungli2015}。

\textbf{边界}。把趋势简单等同于“正偏”是危险的。动量崩溃的存在说明趋势策略也有深度左尾；而短期反转（提供流动性）的偏度未必是负的，因为流动性补偿本身可能呈正偏。偏度对偶严格成立于\emph{中低频、跨资产的组合口径}，在单资产高频繁口径上需要重新检验。

\subsection{对偶三：时间尺度——同一个自相关函数的三个区间}

\textbf{命题}：趋势与反转的归属由持有期决定，而非由信念决定。

\textbf{论证}。第~\ref{sec:unity}~节已经给出方差比 $VR(q)$ 的统一表述。这里补两点关键细节。

第一，美国市场的经验形状是\emph{波动}而非单调的。Lo \& MacKinlay(1988) 对 1962--1985 年 CRSP 数据做方差比检验，发现等权指数在周频显著偏离随机游走，而单只个股更接近随机游走\cite{lomackinlay1988}——这意味着反转信号在\emph{组合层面}比在\emph{个股层面}更容易被检验到，因为个股的特质噪声稀释了自相关。这是配对交易与截面反转选择“篮子”而非单资产的技术原因。

第二，长周期反转的证据本身存在争议。Richardson(1993) 明确质疑 Fama--French 的暂时成分结论，指出其证据对小样本与检验设定高度敏感\cite{richardson1993}。因此“长周期反转”的置信度明显低于“中周期动量”，这一点将在结论置信度章（第~\ref{sec:confidence}~节）体现。

\textbf{证据}。见第~\ref{sec:matrix}~节的统一比较矩阵中“持有期”一行，以及本文第~\ref{sec:empirical}~节的对照回测：同一资产池上，趋势与反转的相对夏普应当随 $L$ 单调反转。

\textbf{工程映射}。这一对偶有一个常被忽视的推论：\emph{“策略选择”其实可以转化为“持仓周期选择”}。组合管理的自由度因此从“选哪个策略”变成了“在哪几个尺度上分配风险预算”。Wood 等(2021) 的“慢动量加快反转”正是这一思路的实现：用同一信号的高频分量做反转、低频分量做趋势\cite{wood2021}。这不是把两个策略拼起来，而是承认它们是同一信号的两个频段。

\textbf{边界}。尺度依赖不是普适常数。交易成本、市场微观结构、投资者构成都会改变自相关函数的位置与幅度；同一 $L$ 在不同市场可能落在不同区间。因此“$L=12$ 个月是动量”这类说法是特定样本的经验概括，不是结构性定律。

\subsection{对偶四：风险来源——承担收敛风险 vs 收获延迟反应}

\textbf{命题}：反转策略的收益主要来自\emph{提供流动性、承担收敛失败风险}；趋势策略的收益主要来自\emph{风险溢价与信息延迟扩散}。前者是卖方（被资助方），后者是买方（资助方）。

\textbf{论证}。反转策略的经济角色是“当别人急着卖时买入、急着买时卖出”。这个动作本质上是做市：交易对手的紧迫需求把价格推离合理水平，反转交易者提供即时性，赚取补偿。其代价是承担两件事——价差可能继续扩大（不收敛），以及自己可能在需要资金时被抽资。Shleifer \& Vishny(1997) 的套利限制理论把第二点讲得最清楚：由于业绩与资金流的正反馈，套利者在错误定价加深时反而被迫减仓，从而无法等到收敛\cite{shleifervishny1997}。Brunnermeier \& Pedersen(2009) 的融资流动性模型给出同构的机制链条：价格下跌→保证金上升→被迫平仓→价格进一步下跌\cite{brunnermeier2009}。

趋势策略的经济角色是“在价格上涨后继续买入”。它的收益来源在文献中有两条解释。一条是风险补偿：趋势在极端市场中获利，说明它承担了某种在危机中才兑现的、被市场厌恶的风险\cite{hop2017}；另一条是行为：投资者对新信息的反应存在延迟，价格的调整不是一次完成，而是分步扩散。Hong \& Stein(1999) 的信息渐进扩散模型为后者提供了严格表述\cite{hongstein1999}，Barberis 等(1998) 与 Daniel 等(1998) 则从保守主义与过度自信两个角度给出互补解释\cite{bsv1998,dhs1998}。

\textbf{工程映射}。风险来源的差异解释了为什么两策略的“加仓纪律”方向相反。反转策略在亏损时应\emph{减少}而非增加敞口——因为亏损意味着不收敛的概率在上升，而资金约束会紧随其后。趋势策略在盈利时加仓是合理的（凸性），但在亏损时的加仓需要严格的波动率边界，否则会滑向动量崩溃。

\textbf{边界}。风险来源的二分不是互斥的。Ehsani \& Linnainmaa(2022) 的因子动量证据表明动量可以在因子层面被重新解释\cite{ehsani2022}；而反转收益中相当一部分被证明是流动性成本而非独立 alpha（Lo \& MacKinlay(1990) 的交叉自相关批评正是此意）\cite{lomackinlay1990}。因此更稳妥的表述是：\emph{两者的风险来源有主导倾向，而非排他归属}。

\subsection{对偶五：容量与换手——成本敏感度相差一个量级}

\textbf{命题}：反转策略的换手量级显著高于趋势策略，因而对交易成本敏感得多；趋势策略虽换手较低，却面临更严重的容量与拥挤问题。

\textbf{论证}。反转信号的半衰期短。以 OU 框架为例，最优交易带宽度与成本呈三次方根关系（Martin \& Sch\"oneborn 2011\cite{martin2011}），成本越高，带宽越宽、交易越稀——这意味着在给定成本下，反转策略被迫放弃一部分信号。Do \& Faff(2012) 的实证直接验证了这一点：把 Gatev 等的样本延伸并引入现实成本后，配对交易利润大幅衰减\cite{dofaff2012}。Korajczyk \& Sadka(2004) 用日内数据估计动量策略的成本，给出了盈亏平衡的基金规模\cite{korajczyksadka2004}；Frazzini 等(2018) 用实盘交易数据估计异象的真实成本，指出换手高的因子净收益衰减最重\cite{frazzini2018}；Novy-Marx \& Velikov(2016) 建立的异象成本分类法也指向同一结论\cite{novymarxvelikov2016}。

趋势一侧，King 等(2016)\cite{kim2016}与 Baltas \& Kosowski(2013) 的研究显示，时间序列动量的换手可以通过信号平滑（如线性趋势拟合、Yang--Zhang 区间估计量）显著降低，且样本外表现更好。换手低意味着成本敏感度低——这是趋势策略在长期样本中“看似稳健”的一个技术原因。

但趋势的代价是拥挤。Stein(2009) 的理论指出，对于缺乏锚定的策略（动量是典型代表），套利者无法知道有多少同行在同时进场，协调失败会放大价格偏离；同时杠杆外溢制造火售外部性\cite{stein2009}。Lou \& Polk(2022) 提出的 comomentum 指标从相关性角度度量动量拥挤度，发现高拥挤之后动量收益下降\cite{loupolk2022}；Baltas(2019) 系统研究了另类风险溢价策略的拥挤与容量\cite{baltas2019}。

\textbf{工程映射}。这一对偶决定了两个策略的规模天花板不同。反转策略的约束来自\emph{成本}：单笔交易的期望收益在扣除成本后可能所剩无几，因此它天然适合规模较小、速度较快的资金。趋势策略的约束来自\emph{拥挤}：信号本身在低换手下可容纳大资金，但同质化导致信号失效，因此它天然适合规模较大、但需要持续分散与容量监控的资金。

\textbf{边界}。这个二分有例外：跨资产的趋势策略在流动性差的品种上换手成本同样很高；而某些低频反转（如多年尺度的价值反转）换手极低。因此准确的表述应绑定\emph{尺度}：短尺度反转成本敏感，中尺度趋势拥挤敏感。

\subsection{对偶六：失效模式——“慢慢失血”与“突然跳崖”的错配}

\textbf{命题}：两个策略都会亏钱，但亏法不同。反转策略的失效是“不收敛+资金断裂”的联合事件，表现为突然的深亏；趋势策略的失效是“震荡市慢性出血”叠加“崩溃时跳崖”。

\textbf{论证}。反转策略的典型失败路径在两次历史事件中被完整记录。LTCM 的收敛交易在价差扩大时本应加仓等待，但杠杆与赎回约束让它无法等到收敛\cite{lowenstein2000,jorion2000}。2007 年 8 月的量化危机中，多个统计套利基金同时去杠杆，Khandani \& Lo(2011) 识别出的两次平仓事件集中在数小时到数日内完成，起始头寸集中于金融股、多账面市值比、空盈利动量\cite{khandanilo2011}。共同的结构是：\emph{亏损本身不可怕，可怕的是亏损触发了资金约束}。

趋势策略的失败有两副面孔。震荡市中是慢性出血：价格反复穿越信号阈值，策略不断止损，收益为负但幅度温和，波动率目标化会进一步降低敞口（Harvey 等 2018 指出这正是波动率目标的代价之一\cite{harvey2018}）。危机中是突然跳崖：动量崩溃。Daniel \& Moskowitz(2016) 强调崩溃的\emph{可预测性}——它发生在恐慌状态，且部分可由市场波动率与市场收益预测，据此构造的动态动量策略把 alpha 与夏普大约翻倍\cite{danielmoskowitz2016}。

\textbf{工程映射}。失效模式差异对应两套不同的“熔断”设计。反转策略需要的是\emph{资金冗余与止损线}：设定价差偏离的容忍上限，超过则减仓，且预留足够的保证金缓冲以防被迫平仓；不能用“越跌越买”的加仓逻辑。趋势策略需要的是\emph{状态过滤}：在恐慌状态下降低甚至反转敞口，避免在崩溃中继续追涨。

\textbf{边界}。两策略的失效在极端去杠杆事件中会同时发生——2007 年 8 月中，量化反转基金与动量基金都被冲击，因为冲击来自共同的资金端而非各自的信号端\cite{khandanilo2011,brunnermeier2009}。这说明\emph{分散化在资金冲击下会失效}，这是组合层面而非策略层面的问题。

\subsection{对偶七：入场门槛——先验证后交易 vs 边交易边估计}

\textbf{命题}：反转策略在结构上多一道“资格检验”；趋势策略没有这道门槛，但需要持续估计信号的稳定性。前者是“条件性”策略，后者是“无偏条件”策略。

\textbf{论证}。配对交易的标准流程是：选对（距离法或协整法）→ 检验价差平稳性 → 估计 OU 参数（回归速度、半衰期）→ 设定阈值带 → 交易。每一步都是筛选，不合格的样本被排除。Vidyamurthy(2004) 的专著系统化了这一流程\cite{vidyamurthy2004}；Elliott 等(2005) 用 Kalman 滤波在线估计 OU 状态，把“资格检验”从一次性改为滚动进行\cite{elliott2005}。

趋势策略的流程是：估计动量统计量（过去收益的符号、排序或回归斜率）→ 按波动率缩放建仓 → 定期调仓。没有独立的资格检验步骤，因为动量信号对每一个资产都可以计算，不需要先证明该资产“适合动量”。

这两条流程的差别带来可检验的推论：\emph{反转策略的有效样本数少于趋势策略}。反转需要先筛选出满足条件的配对/篮子，筛选失败的资产被浪费；趋势则可以用满整个资产池。在资产数量有限的市场（如单一期货品种），反转策略的样本基础更薄。

\textbf{工程映射}。这道门槛使反转策略的\emph{超参数}更多（选对规则、协整显著性、半衰期上下限、阈值带宽度），因而数据窥探的空间更大。这既是它容易过拟合的原因，也解释了为何配对交易文献中“利润衰减”证据如此普遍——Do \& Faff(2012) 的发现可以部分归因于早期研究的参数选择在样本外失效\cite{dofaff2012}。

\textbf{边界}。这个二分有反例：截面动量的“资格”隐含在排序中（只有流动性足够的资产才能进入组合），所以趋势策略并非完全没有筛选，只是筛选维度不同（流动性 vs 平稳性）。

\subsection{旁支落地：配对筛选的七个步骤}

对偶七断言反转多一道资格检验。把这道检验写成算法，就是配对交易文献的标准流程\cite{vidyamurthy2004}——七步，每步一个判据，不合格的样本在原地被拒绝，走完全程的才进入交易。

\begin{enumerate}[leftmargin=2.2em]
  \item \textbf{分组生成候选对}。只在同产业链或同板块内两两组合。跨组配对缺少经济意义上的长期均衡，协整检验在这种对上很容易给出伪通过。
  \item \textbf{相关性初筛}。形成期两序列对数价格的相关系数设下限（常见取 $0.6$）。这一步不是统计检验，是算力与稳健性的双重过滤——低相关配对即便偶然通过后续检验，价差波动也大到不可交易。
  \item \textbf{协整检验}。Engle--Granger 两步法\cite{englegranger1987}：对形成期做回归 $\ln P^A_t=\alpha+\beta\ln P^B_t+e_t$，再对残差 $e_t$ 做单位根检验\cite{dickeyfuller1979}，$p<0.05$ 才放行。这是全流程唯一的“资格检验”，对应对偶七命题里的平稳性要求。
  \item \textbf{回复速度估计}。对残差拟合稳定 AR(1)：$e_{t+1}=c+\phi e_t+u_t$，要求 $0<\phi<1$；换算连续时间参数 $\theta=-\ln\phi/\Delta$，半衰期 $H=\ln 2/\theta$。$\phi\ge1$ 意味着价差不回归，$\phi\le0$ 意味着高频振荡而非缓慢回复，两者都让半衰期失去意义。
  \item \textbf{半衰期区间筛选}。$H$ 设上下界：下界挡掉回复过快的对（交易频率高、成本吃光价差），上界挡掉占用资金过久的对（期间结构性断裂的概率随时间上升）。
  \item \textbf{滚动重估}。形成期与交易期严格分离、滚动推进，杜绝用交易期数据选对造成的 look-ahead——这一范式来自 Gatev 等(2006) 的 $6+6$ 个月设计\cite{gatev2006}；交易期的均值与标准差须用滚动统计量，不能沿用形成期常数，否则波动放大时 $|z|$ 虚高、开仓即触止损。
  \item \textbf{生成可交易对表}。输出对冲比率（静态回归系数，或 Kalman 滤波的在线估计\cite{elliott2005}）、价差均值与标准差、半衰期，据此设定开平仓阈值带。
\end{enumerate}

七步的顺序不可调换，且每一步的参数（相关性下限、显著性水平、半衰期区间、阈值带宽）都是自由度。这一点必须说透：\emph{筛选算法本身也在做多重检验}——一百个候选对过筛后剩五个“平稳”对，这五个的平稳性是选择的结果，不是证据。因此筛选参数必须在实验前锁定、并整体计入第~\ref{sec:audit}~节的试验族审计，否则筛选就从资格检验退化为变相的数据窥探。

\subsection{七组对偶的合并判断}

把七组对偶合起来，得到一个修正后的判断：

\begin{quote}
\itshape
趋势与均值回归共享同一个信号，却生活在两套约束里。
反转策略在一个“需要资格证明”的世界里做流动性生意，它的敌人是成本与资金约束，
它的失效方式是突然的；趋势策略在一个“无需资格证明”的世界里做延迟反应的生意，
它的敌人是拥挤与状态误判，它的失效方式是两面性的——慢性出血与尾部跳崖。
\end{quote}

这个判断是本文的第二个核心结论。它比“两策略对立”更精确，也比“两策略等价”更诚实。第~\ref{sec:execution}~节将把它翻译成具体的执行差异清单。

\section{统一比较矩阵}\label{sec:matrix}

第~\ref{sec:unity}~节与第~\ref{sec:duality}~节分别给出了“同”与“异”的论证。本章把二者压缩成两张表：一张是定性的对偶矩阵，一张是“论断—证据—置信度”的映射表。前者用于速查，后者用于审计——它的作用是把正文每一个可争议的判断落到具体来源上，使读者能够独立复核。

\subsection{对偶矩阵}

\begin{table}[htbp]
\centering\footnotesize
\caption{趋势跟踪与均值回归的逐维对偶（定性矩阵）}\label{tab:dual}
\begin{tabular}{@{}>{\raggedright\arraybackslash}p{2.5cm}>{\raggedright\arraybackslash}p{6.1cm}>{\raggedright\arraybackslash}p{6.1cm}@{}}
\toprule
\textbf{比较维度} & \textbf{趋势跟踪} & \textbf{均值回归} \\
\midrule
信号构造 & $w_t=+\,g\!\left(s_t(L)/\sigma_t\right)$：顺过去走势 &
$w_t=-\,g\!\left(s_t(L)/\sigma_t\right)$ 或 $w_t=-\,g\!\left((X_t-\mu)/\sigma\right)$：逆偏离 \\
数学载体 & OU 族 $\theta<0$ 分支；带漂移随机游走 &
OU 族 $\theta>0$ 分支；协整误差修正模型 \\
统计前提 & 无独立门槛，正自相关即可 & \textbf{先验证}平稳性/协整，再交易 \\
典型持有期 & 一至十二个月（中尺度） & 小时至周（短反转）、三至五年（长反转） \\
换手量级 & 中低（信号平滑后可进一步压低） & 高，半衰期短者尤甚 \\
成本敏感度 & 中 & 高；成本与最优带宽呈三次方根关系 \\
收益偏度 & 总体右偏，但含深度左尾（动量崩溃） & 左偏，收敛失败时集中亏损 \\
极端市场表现 & crisis alpha：危机中往往正收益 & 去杠杆中与同类策略同向巨亏 \\
容量约束来源 & 拥挤与协调失败（信号同质化） & 交易成本与资金约束 \\
失效模式 & 震荡市慢性出血 + 崩溃时跳崖 & 不收敛 + 被迫平仓，跳变式亏损 \\
风险管理手段 & 波动率目标化有效 & 波动率目标可能有害；需止损与资金冗余 \\
主导机制解释 & 风险溢价 / 信息延迟扩散 & 流动性补偿 / 过度反应 \\
证据总体强度 & 中周期动量证据强；百年回测样本长 & 配对交易早期强、后期衰减；短反转证据强 \\
数据窥探暴露 & 高（技术规则已被系统审查） & 高（超参数多、配对选择自由度高） \\
\bottomrule
\end{tabular}
\end{table}

\subsection{矩阵读法：三行容易被误读的地方}

\textbf{“风险管理手段”一行最反直觉}。波动率目标化被反复证明能提升动量/趋势策略的夏普\cite{barroso2015,harvey2018,moreiramuir2017}，但机械照搬到反转策略上是错的：趋势在波动率上升时信号质量下降（趋势被噪声打断），降敞口合理；反转的利润恰恰常来自高波动时段（偏离更大、回归空间更大），降敞口等于砍掉自己的收益来源。执行章（第~\ref{sec:execution}~节）展开这一点。

\textbf{“统计前提”一行的实操后果}是反转的研发流程多一道\emph{样本消耗}——先筛出满足协整/平稳条件的配对或篮子，不合格的样本被丢弃。资产数量有限的市场里，反转的有效样本基础因此比趋势薄，统计显著性也更难建立。这是“入场门槛”对偶的操作含义。

\textbf{“容量约束来源”一行最容易被忽略}。业界的直觉是“反转容量小、趋势容量大”；矩阵的画面更精细——反转的约束来自\emph{自己支付的成本}，确定、可测算；趋势的约束来自\emph{别人和你抢同一笔钱}，状态依赖、难以事前量化。容量研究（Stein 2009；Lou--Polk 2022）\cite{stein2009,loupolk2022}在过去十几年反复尝试却始终没有统一答案，根子就在这里。

\subsection{论断—证据映射}

下表把本文的核心论断逐条对上证据，并给出各自的置信度分级（分级口径见第~\ref{sec:confidence}~节）。

\begin{table}[htbp]
\centering\footnotesize
\caption{核心论断与证据映射}\label{tab:claim}
\begin{tabular}{@{}>{\raggedright\arraybackslash}p{3.5cm}>{\raggedright\arraybackslash}p{6.4cm}>{\raggedright\arraybackslash}p{1.5cm}>{\raggedright\arraybackslash}p{2.9cm}@{}}
\toprule
\textbf{论断} & \textbf{主要证据} & \textbf{置信度} & \textbf{主要反方/边界} \\
\midrule
两策略共享同一信号，仅符号不同 & 自相关函数多尺度符号反转\cite{lomackinlay1988,cutler1991,famafrench1988} & \confhigh & 协整型反转含结构约束\cite{englegranger1987} \\
收益自相关随持有期变号（U 形） & 短频反转\cite{lehmann1990,jegadeesh1990}、中频动量\cite{jegadeesh1993,mop2012}、长频反转\cite{debondt1985,poterbsummers1988} & \confhigh & 长周期反转证据受质疑\cite{richardson1993} \\
趋势策略呈凸性/危机中获利 & 回溯跨式复制\cite{funghsieh2001}、百年样本\cite{hop2017} & \confmid & 动量崩溃显示深度左尾\cite{danielmoskowitz2016} \\
反转策略呈凹性/收敛失败致灾 & LTCM\cite{lowenstein2000,jorion2000}、2007 量化危机\cite{khandanilo2011} & \confmid & 短期反转（流动性型）偏度未必为负 \\
反转换手高于趋势、成本敏感 & 配对利润衰减\cite{dofaff2012}、成本分类\cite{novymarxvelikov2016,frazzini2018} & \confhigh & 低频价值型反转换手很低 \\
趋势面临拥挤而非成本约束 & 协调失败\cite{stein2009}、comomentum\cite{loupolk2022}、拥挤度\cite{baltas2019} & \confmid & 容量约束的定量估计仍不统一 \\
两策略均受数据窥探威胁 & 7846 规则 bootstrap\cite{sullivan1999}、Reality Check\cite{white2000}、$t>3$~门槛\cite{hlz2016} & \confhigh & 方法学上可部分校正\cite{hansen2005,romanowolf2005} \\
发表后收益显著衰减 & 97 个因子实测\cite{mcleanpontiff2016} & \confmid & 衰减幅度因因子而异 \\
时间序列动量稳健性存疑 & 逐资产再检验\cite{huang2020}、跨资产信号\cite{pitkajarvi2020} & \confmid & 百年样本持相反结论\cite{hop2017} \\
\bottomrule
\end{tabular}
\end{table}

\section{机制解释的对照}\label{sec:mechanism}

信号层与统计层的比较完成之后，还有一个问题悬着：\emph{为什么会这样}。本章对照两条脉络对机制的争夺。这不是学理上的装饰——机制决定了策略在什么条件下应当被信任、在什么条件下应当被放弃。

\subsection{两套解释框架}

文献中对收益可预测性的解释可以粗分为两族。

\textbf{风险补偿族}认为，可预测性来自时变的风险溢价：某些状态下的预期收益更高，是因为那些状态下风险更大。动量在危机中获利，说明它承担了被市场厌恶的风险；反转策略的收益来自承担流动性风险与资金约束风险。

\textbf{行为偏差族}认为，可预测性来自投资者的系统性认知偏差：反应不足产生中短期动量，反应过度产生长期反转。

两族解释在文献中长期并存。Cochrane(2011) 的贴现率视角提供了一个把二者装进同一框架的努力：收益率可预测性的核心是贴现率的时变性，而贴现率的变化既可以是风险溢价的理性变化，也可以是情绪驱动的错价\cite{cochrane2011}。Cochrane(2005) 的综述进一步指出，区分二者的困难在于\emph{两者都能拟合同样的预测回归结果}\cite{cochrane2005}。

\subsection{反转一侧的机制}

反转策略的机制解释有三个层次。

\textbf{第一层是过度反应}。De Bondt \& Thaler(1985) 的原始解释：投资者对坏消息过度悲观，把价格压到低于合理水平，随后价格回升\cite{debondt1985}。Summers(1986) 与 Shiller(1981) 的价格波动过大证据为这一解释提供了宏观背景\cite{summers1986,shiller1981}。Lakonishok 等(1994) 把这一逻辑推广到价值投资：反向投资之所以有效，是因为市场外推了过去的增长\cite{lakonishok1994}。

\textbf{第二层是噪声交易者风险}。De Long 等(1990) 的模型给出了一个更精细的机制：即便理性套利者知道价格偏离，他们也不会立即纠正，因为噪声交易者可能让偏离进一步扩大，套利者承担了额外的风险\cite{delong1990}。这一模型解释了\emph{为什么偏离能持续}——而这正是反转策略需要“等待”的原因。

\textbf{第三层是流动性补偿}。这是最被后来的实证支持的一层。Lo \& MacKinlay(1990) 指出大量“反转利润”实为交叉自相关造成的聚合假象\cite{lomackinlay1990}；短周期的反转收益与买卖价差、流动性提供高度绑定。Avellaneda \& Lee(2010) 的因子剥离框架把反转定义为“残差的均值回复”，其实质是在剥离共同因子暴露后交易特质性错价\cite{avellaneda2010}。流动性作为定价因子的地位由 Pástor 与 Stambaugh(2003) 系统化：小盘股对市场流动性冲击更敏感，流动性风险的补偿构成收益的一部分\cite{pastorstambaugh2003}。这一层也解释了反转策略为什么容易在资金端出问题——Shleifer 与 Vishny(1997) 之后，Gromb 与 Vayanos(2010) 对该族理论做了系统综述\cite{shleifervishny1997,grombvayanos2010}。

三层机制并不互斥，而是作用在不同时间尺度上：短尺度以流动性补偿为主，中尺度混合过度反应与噪声交易者风险，长尺度以估值回归为主。这恰好对应第~\ref{sec:duality}~节“尺度决定归属”的判断。

\subsection{趋势一侧的机制}

趋势策略的机制解释同样分层。

\textbf{延迟反应}。Hong \& Stein(1999) 的信息渐进扩散模型是这一族最完整的表述：市场由“消息观察者”与“动量交易者”两类参与者构成，前者缓慢吸收私有信息，后者仅从价格推断信息，二者的互动先产生反应不足、后产生过度反应\cite{hongstein1999}。Barberis 等(1998) 用保守主义与代表性启发两个偏差给出平行解释\cite{bsv1998}；Daniel 等(1998) 用过度自信与自我归因给出第三个版本\cite{dhs1998}。三篇文献的构造不同，但都产出“先动量后反转”的同一预测——这与第~\ref{sec:unity}~节的自相关 U 形结构吻合。

\textbf{风险溢价}。Hurst 等(2017) 的百年证据是这一解释的主要支撑：趋势跟踪在每个十年、各种宏观环境下都为正，且与股债的长期相关性低\cite{hop2017}。若趋势收益纯属错价，它应当随市场成熟而消失；百年尺度上的稳定性更符合“承担了某种系统性风险”的解释。Erb \& Harvey(2006) 从商品展期收益的角度提供了一个具体的风险来源候选\cite{erbharvey2006}。Ilmanen(2011) 对各类风险溢价的系统梳理为这一族解释提供了统一框架：动量的收益可以按“承担了何种被厌恶的风险、在何种环境下兑现”来归类\cite{ilmanen2011}。一个有助于区分机制的线索是：若动量只是“对近期信息反应不足”的表象，那么它与基本面动量的相关性应当很高；Novy-Marx(2012) 检验后发现并非如此，动量并非主要由基本面动量驱动\cite{novymarx2012}。

\textbf{拥挤与协调失败}。Stein(2009) 的论证是第三种解释的雏形：动量属于“未锚定策略”（unhedged strategy），套利者无法知道同行的建仓规模，因而无法判断价格偏离是否会被纠正\cite{stein2009}。这一机制不解释趋势为何长期有效，但解释它为何会突然失效——这与动量崩溃的实证相符\cite{danielmoskowitz2016}。

\subsection{一个机制生成两个符号}

本章最重要的观察是：\textbf{不需要两套机制来解释两策略}。行为模型（Hong--Stein、BSV、DHS）在单一框架内同时产生短期动量的正自相关与长期反转的负自相关。这意味着第~\ref{sec:genealogy}~节观察到的“分岔”在机制层面根本不存在——两条脉络是在\emph{同一个机制的产物}上各自取了不同的一端。

这一观察有直接的操作含义。如果两策略来自同一机制，那么它们的失效也应当同时发生，因为机制被破坏时两个方向都会失效。经验证据支持这一点：2007 年 8 月，量化反转基金与趋势基金同时遭受冲击，因为冲击来自共同的资金端而非各自的信号端\cite{khandanilo2011,brunnermeier2009}。这解释了为什么简单地把两个策略等权组合起来并不能提供想象中的分散化——它们在某些状态下的相关性会趋近于一。

\subsection{机制之争对策略的实际意义}

机制解释的分歧不是纯粹的学术趣味，它决定了三件事。

\textbf{一是对“样本外失效”的解释}。若趋势收益是风险溢价，发表后不应衰减（风险仍需承担）；若是错价，则应当随套利而衰减。McLean \& Pontiff(2016) 发现的发表后衰减支持后者\cite{mcleanpontiff2016}；Hurst 等(2017) 的百年稳定性支持前者\cite{hop2017}。两种证据并存，说明趋势收益可能同时含两个成分，且比例随时间变化。

\textbf{二是对“何时应该收手”的判断}。若反转收益主要是流动性补偿，那么在流动性环境改善（如算法做市普及）后应当衰减——这与 Do \& Faff(2012) 报告的配对交易利润衰减时序吻合\cite{dofaff2012}。

\textbf{三是对“该不该加杠杆”的判断}。若趋势收益是风险溢价，加杠杆应当获得成比例的更高收益（在无摩擦条件下）；若是错价，加杠杆会放大拥挤风险。Stein(2009) 的杠杆外溢论证正是针对后者的警告\cite{stein2009}。

\subsection{边界：机制不可直接观测}

必须承认本章讨论的机制都无法被直接观测。我们观测到的只是收益序列与自相关结构，机制是\emph{对这些观测的解释}。Cochrane(2011) 的说法值得引述：几乎所有资产定价结果都可以同时用理性风险与行为偏差来解释，因为两者在价格数据上常难以区分\cite{cochrane2011}。

因此本章的结论应当保守表述：\emph{关于机制的争论应当影响我们对策略稳健性的先验，而不应当成为选择策略的依据}。选择依据应当是可以检验的东西——自相关结构在哪一段尺度上显著、成本与容量的实测约束、以及资金端的状态。这也是本文把重心放在第~\ref{sec:execution}~节与第~\ref{sec:empirical}~节的原因。

\subsection{能否设计一个可判决的检验}

机制不可直接观测，不等于无法设计区分性的检验。可行路径有三条，各自的可行性与困难值得写清，因为它们构成第~\ref{sec:future}~节遗留课题 Q9 的具体内容。

\textbf{路径一：看样本外是否衰减}。若趋势收益是风险溢价，它在文献发表后\emph{不应}系统性衰减（风险仍需有人承担，补偿仍在）；若是错价，则应当随套利而衰减。McLean 与 Pontiff(2016) 的实测支持后者\cite{mcleanpontiff2016}；Hurst 等(2017) 的百年稳定性支持前者\cite{hop2017}。困难在于两者可以同时成立——收益中同时含溢价成分与错价成分，比例还随时间变化，这使得“衰减多少”无法直接读成“溢价占多少”。

\textbf{路径二：看辅助变量的预测方向}。若反转收益来自流动性补偿，那么它与流动性指标（买卖价差、成交量、资金成本）应当同向变动；若来自过度反应，则应与情绪指标同向。Pástor 与 Stambaugh(2003) 的流动性因子框架提供了可用的代理\cite{pastorstambaugh2003}，Brunnermeier 与 Pedersen(2009) 的融资流动性模型给出了理论预测方向\cite{brunnermeier2009}。这条路径的困难在于代理变量的内生性：流动性差本身就是价格偏离的一部分。

\textbf{路径三：用状态的不可预测性做自然实验}。若收益是风险补偿，它在事前可识别的“高风险状态”下应当更高；若是错价，其高低应与是否被套利相关，而与事前风险状态关系较弱。Daniel 与 Moskowitz(2016) 对恐慌状态的刻画、以及 Ang 与 Timmermann(2012) 综述的 regime 方法提供了状态划分工具\cite{danielmoskowitz2016,angtimmermann2012}。困难是状态划分本身依赖模型，而不同模型给出的状态划分会导致不同的检验结论——这构成了近乎循环的依赖。

三条路径共同的障碍是同一个：\emph{任何用于区分机制的变量，都同时被机制所影响}。这解释了为什么这一问题在几十年后仍未闭合。本文的态度是把它记录为开放的困难，而不给出一个虚假的裁决。

\section{实际执行中的差异}\label{sec:execution}

前几章的比较都在信号层与统计层。但业界与文献反复给出同一个教训：\emph{统计上显著不等于能赚钱}。Martin \& Sch\"oneborn(2011) 把这个教训写进了标题——“均值回归能赚，但要付代价”\cite{martin2011}。本章把第~\ref{sec:duality}~节的对偶翻译成执行层面的具体差异，按从信号到资金的顺序展开。

\subsection{第一步就分岔：信号平滑与换手控制}

两个策略在离开信号层的第一步就走了不同的路。

\textbf{趋势一侧的核心工程问题是“少交易”}。动量信号在日频上噪声极大，若直接按日频符号调仓，换手会高到成本吞掉全部收益。文献给出的主流解法有三类：一是用线性趋势拟合的斜率替代离散符号（Baltas \& Kosowski 2013 报告线性趋势拟合信号样本外最优且换手最低\cite{baltaskosowski2013imp}）；二是用区间估计量替代收盘价波动率以降低噪声（Yang--Zhang 估计量在效率与偏差间取得平衡，同文献）；三是对头寸做指数平滑或部分调整（Gârleanu--Pedersen 的“目标前置、部分调整”给出了含成本的动态最优解\cite{garleanu2013}）。Kim 等(2016) 的系统检验进一步指出，时间序列动量的收益中有相当部分来自波动率缩放本身而非时序符号\cite{kim2016}——这意味着平滑与缩放的技术选择会显著改变策略的经济含义。

\textbf{反转一侧的核心工程问题是“带宽”}。反转策略不能频繁穿越阈值，否则会在价差尚未收敛时反复止损。理论给出的答案是：最优交易带（no-trade band）的宽度随交易成本以三次方根增长\cite{martin2011}。这一结论的实操含义是：\emph{成本越高的市场，反转策略应当越“迟钝”}，宁愿漏掉一部分信号，也不要在噪声中频繁进出。Boguslavsky 与 Boguslavskaya(2004) 在幂效用下给出的解析阈值\cite{boguslavsky2004}、Zhang 与 Zhang(2008) 的“低买高卖”最优规则\cite{zhangzhang2008}、以及 Leung 与 Li(2015) 引入止损退出的双停时解\cite{leungli2015}，构成了这一族的技术基础。

两者的对称性值得注意：趋势在“平滑信号”，反转在“拓宽带宽”。前者压的是信号的方差，后者压的是交易的频率。目标函数都是换手成本，手段却因信号的可加性不同而不同——趋势信号可以连续加权，反转信号必须通过阈值离散化。

\subsection{组合构建与中性化}

在两个策略上，组合层的任务不同。

\textbf{趋势策略需要“去相关”与“风险平价”}。跨资产的趋势信号在危机中高度同向，因此组合必须按资产波动率与相关性分配风险，否则单一资产会主导组合。行业/板块层面的中性化在股票动量中是必需的——Moskowitz \& Grinblatt(1999) 指出行业动量解释了大部分个股动量\cite{moskowitzgrinblatt1999}，若不中性化，组合实际上在押注行业轮动而非个股趋势。

\textbf{反转策略需要“彻底的非系统性”与“分散”}。配对交易的价差交易天然对冲了共同因子暴露，这是它的优势；但代价是\emph{特质风险无法被市场因子抵消}，因此必须靠配对数量分散。Avellaneda \& Lee(2010) 的 PCA 残差框架正是为这一目的设计的：先剥离前若干主成分（共同因子），再在残差上做均值回复\cite{avellaneda2010}。若不做因子剥离，反转组合会混入大量方向性风险，其“收敛”叙事就不成立。

一个容易被忽略的工程细节：反转策略的\emph{对冲比率}是估计量而非常数。Elliott 等(2005) 用 Kalman 滤波在线估计配对价差的 OU 状态\cite{elliott2005}，其工程含义是：对冲比率随时间变化，若沿用静态比率（如形成期的回归系数），价差序列会因参数漂移而失真。趋势策略没有对应的困扰，因为它的仓位由自身波动率决定，不依赖跨资产的稳定关系。

\subsection{波动率目标化的不对称作用}

这是本文认为执行层最重要的一条差异。

波动率目标化的做法是：按“目标波动率 / 已实现波动率”缩放仓位。它在趋势策略上被反复证明有效——Barroso \& Santa-Clara(2015) 把动量夏普从 0.53 提到 0.97、超额峰度从 18.2 降到 2.7\cite{barroso2015}；Harvey 等(2018) 在 60 多个资产上系统验证其改善风险资产夏普\cite{harvey2018}；Moreira \& Muir(2017) 表明波动率管理能提升多数因子的表现\cite{moreiramuir2017}。

但把这套方法搬到反转策略上会产生\emph{方向性错误}。理由是信号与波动率的关系在两策略中符号相反：趋势策略在波动率上升时信号质量下降（趋势被噪声打断），降敞口是保护；反转策略的利润经常\emph{恰好}来自高波动时段——偏离更大，回归的绝对空间也更大。因此在波动率上升时降低反转敞口，等于在信号最强时把仓位压到最小。

正确做法是把波动率信息用在\emph{风险预算}而非\emph{信号缩放}上：对反转策略，波动率上升应当触发的是止损线的收紧与资金冗余的检查，而不是机械的仓位缩放。Leung \& Li(2015) 把止损退出纳入最优框架，正是这一思路的理论化\cite{leungli2015}。

\subsection{成本与执行机制}

成本模型的选择会系统性影响两策略的比较结论。

\textbf{趋势策略的成本压力来自再平衡}。月频或周频调仓下，单次调仓的换手中等，但一年累积的换手仍可观；成本主要在“定期滚动”中流失，可以用信号平滑直接压低。

\textbf{反转策略的成本压力来自进出频率}。它的成本具有明显的非线性：价差偏离越大时，往往也是流动性越差的时刻（流动性螺旋），因此实际冲击成本高于按平均换手估算的值。Brunnermeier \& Pedersen(2009) 的模型给出了这一非线性\cite{brunnermeier2009}。工程上应当采用\emph{参与率约束的冲击成本模型}（冲击与参与率的平方根成正比），并对超过参与率上限的部分显式标记、不进入回测收益。

执行算法层面，Almgren \& Chriss(2000) 的框架给出冲击成本与择时风险的权衡\cite{almgrenchriss2000}；Perold(1988) 的“执行缺口”提供了一致的度量口径\cite{perold1988}。对反转策略，执行缺口的斜率更陡（因为它常在流动性差时交易）；对趋势策略，缺口更平坦但累积量更大。

\subsection{资金、杠杆与失效防线}

第~\ref{sec:duality}~节的对偶六指出两策略的失效模式不同，执行层的应对也必须不同。

反转策略的防线是\emph{资金冗余}。Shleifer \& Vishny(1997) 的套利限制理论说明，由于业绩与资金流正反馈，套利者在错价加深时反而被迫减仓\cite{shleifervishny1997}；2007 年 8 月的事件表明这一机制在大规模、多机构同时发生时会被放大\cite{khandanilo2011}。因此反转策略必须预留足以承受“偏离继续扩大”的保证金缓冲，并设定明确的偏离容忍上限——超过上限即减仓，而不是“越跌越买”。

趋势策略的防线是\emph{状态过滤}。动量崩溃发生在恐慌状态（市场下跌后、波动率高企并伴随反弹），且部分可预测\cite{danielmoskowitz2016}。因此趋势策略需要在极端状态下降低甚至反转敞口，而不是无条件追随短期走势。Chabot 等(2014) 的历史研究也指出动量与“追涨”行为在可预测的崩溃前会呈现特征性模式\cite{chabot2014}。

\subsection{Regime 依赖与切换}

两策略的适用状态不同，这一点在文献中有明确的结构。

趋势策略在\emph{持续单向运动}的环境中有效，这类环境通常伴随宏观变量的持续漂移；反转策略在\emph{区间震荡}环境中有效。Hamilton(1989) 的马尔可夫切换模型提供了建模基础\cite{hamilton1989}，Ang \& Bekaert(2002) 把它用于国际资产配置并显示相关性本身随状态变化\cite{angbekaert2002}，Ang \& Timmermann(2012) 的综述系统梳理了 regime 方法在金融中的应用\cite{angtimmermann2012}。

工程上，regime 的用途有两种。一是\emph{硬切换}：只在某状态下启用某策略；二是\emph{软加权}：按状态概率调整两策略的权重。Wood 等(2021) 的“慢动量加快反转”属于更彻底的方案——用变点检测识别状态转换点，让高频分量做反转、低频分量做趋势\cite{wood2021}。这一方案的动机与本文统一框架一致：既然两者是同一信号的不同频段，就不必在它们之间做二选一。

\subsection{组合融合：为什么简单叠加常常不分散}

一个常见的组合做法是把趋势与反转等权叠加，期望获得分散化。这在多数时候不起作用，原因在第~\ref{sec:mechanism}~节已经点出：两策略在某些状态下相关性趋近于一（共同的资金冲击），而在另一些状态下相关性为负（信号符号相反）。因此简单叠加的实际效果是\emph{在信号相反时互相抵消，在资金冲击时同时亏损}——两头不讨好。

更合理的做法有三条。一是\emph{按尺度分层}：短尺度配反转、中尺度配趋势，让不同频段各自发挥作用（对应 Wood 等的频段分解）。二是\emph{按状态加权}：用 regime 概率动态调整权重，而非固定比例。三是\emph{按风险预算而非资金分配}：先评估两策略在目标状态下的波动与尾部，再分配风险预算，避免某一策略在极端状态下主导组合。

\subsection{执行差异清单}

\begin{table}[htbp]
\centering\footnotesize
\caption{执行层面的差异清单}\label{tab:exec}
\begin{tabular}{@{}>{\raggedright\arraybackslash}p{3.0cm}>{\raggedright\arraybackslash}p{5.9cm}>{\raggedright\arraybackslash}p{5.9cm}@{}}
\toprule
\textbf{执行环节} & \textbf{趋势跟踪} & \textbf{均值回归} \\
\midrule
信号处理 & 平滑信号、压低换手；可用斜率替代符号 & 拓宽交易带（宽度 $\propto$ 成本立方根） \\
对冲设计 & 不依赖跨资产关系 & 依赖对冲比率，需在线估计（Kalman） \\
因子暴露 & 需行业/板块中性化，去同向相关 & 需彻底剥离共同因子，纯特质化 \\
风险缩放 & 波动率目标化有效 & 波动率目标化可能有害，改用止损+资金冗余 \\
成本结构 & 累积式，可由平滑压低 & 非线性式，流动性差时冲击更大 \\
资金防线 & 状态过滤，防动量崩溃 & 资金冗余 + 偏离上限，防强制平仓 \\
适用状态 & 持续单向运动、宏观漂移 & 区间震荡、偏离扩大 \\
组合融合 & 需按尺度或状态加权 & 简单叠加会互相抵消 \\
\bottomrule
\end{tabular}
\end{table}

这张清单是本文对实践最直接的一份交付。它给出的核心信息是一句反直觉的话：\textbf{两策略在信号上越接近，执行上的区别就越关键}——因为信号带来的边际 alpha 有限，而执行偏差（成本、对冲、资金）足以决定正负。

\section{统一口径对照实证}\label{sec:empirical}

前七章的比较建立在文献之上。文献的绩效数字彼此不可比（第~\ref{sec:audit}~节），因此需要一个变量受控的对照实验来检验核心论断。本章报告这样一组实验。

\subsection{实验设计}

\textbf{核心设计：唯一自由度是符号}。两策略共享以下全部设定，仅信号符号不同：

\begin{itemize}[leftmargin=2.2em]
  \item \textbf{资产池}：中国期货主力连续合约 45 个品种（经换月跳空守门前复权）；A 股宽基指数 7 个；A 股个股 5083 只（剔除上市不足与长期停牌）。
  \item \textbf{信号}：$s_t(L)=\sum_{k=0}^{L-1}\ln(1+r_{t-k})$，即过去 $L$ 日累计对数收益。趋势取 $+s$，反转取 $-s$。窗口 $L\in\{5,10,20,60,120,250\}$。
  \item \textbf{口径}：时序（TS）——按 $\mathrm{sign}(s_t)$ 乘波动率缩放后归一到总敞口 1；截面（XS）——按 $s_t$ 排序，多前列 20\%、空后列 20\%，两端等权，总敞口 1。
  \item \textbf{调仓}：每 5 个交易日调仓一次；信号使用至 $t$ 日收盘的信息，仓位自 $t+1$ 日生效；持有期内权重随收益漂移，换手按目标权重与实际权重之差计算。
  \item \textbf{成本}：期货单边 2 个基点；A 股指数 5 个基点（ETF/股指期货口径）；A 股个股单边 15 个基点（含佣金、印花税与滑点）。
  \item \textbf{样本}：2014-01-02 至 2026-09-30（3100 个交易日）。
\end{itemize}

\textbf{为什么这组设计能检验核心论断}。第~\ref{sec:unity}~节的论断是“两策略共享同一信号，唯一区别是符号”。若这个论断成立，那么在同一 $L$ 上，趋势与反转的换手应当\emph{完全相同}，成本拖累应当\emph{完全相同}，毛收益序列应当\emph{严格互为相反数}，净收益的差异应当完全由“信号符号 + 常数成本扣项”解释。这构成两条可证伪的验证门（对应第~\ref{sec:reproduce}~节的 G1 与 G2）。

\subsection{门禁检验：论断是否成立}

\begin{table}[htbp]
\centering\footnotesize
\caption{期货时序口径下两策略的对称性检验（G1/G2）}\label{tab:gate}
\begin{tabular}{@{}>{\raggedright\arraybackslash}p{1.6cm}>{\raggedright\arraybackslash}p{1.6cm}>{\raggedright\arraybackslash}p{2.0cm}>{\raggedright\arraybackslash}p{2.0cm}>{\raggedright\arraybackslash}p{2.0cm}>{\raggedright\arraybackslash}p{2.0cm}@{}}
\toprule
$L$ & 方向 & 年化换手 & 成本拖累 & 毛利夏普 & 净夏普 \\
\midrule
\multirow{2}{*}{5} & 趋势 & 52.41 & 1.06\% & $+0.40$ & $+0.20$ \\
                    & 反转 & 52.38 & 1.02\% & $-0.40$ & $-0.60$ \\
\midrule
\multirow{2}{*}{10} & 趋势 & 35.14 & 0.72\% & $+0.45$ & $+0.32$ \\
                     & 反转 & 35.12 & 0.68\% & $-0.44$ & $-0.58$ \\
\midrule
\multirow{2}{*}{60} & 趋势 & 16.14 & 0.33\% & $+0.26$ & $+0.21$ \\
                     & 反转 & 16.13 & 0.32\% & $-0.26$ & $-0.32$ \\
\midrule
\multirow{2}{*}{250} & 趋势 & 8.83 & 0.18\% & $+0.27$ & $+0.24$ \\
                      & 反转 & 8.85 & 0.17\% & $-0.27$ & $-0.30$ \\
\bottomrule
\end{tabular}
\end{table}

表~\ref{tab:gate}~给出三个结论。

\textbf{G1 与 G2 通过}。同一 $L$ 上，趋势与反转的年化换手几乎完全相等（$52.41$ 对 $52.38$，差异来自持有期内权重漂移的微小差异），成本拖累几乎完全相等（$1.06\%$ 对 $1.02\%$）；毛利夏普严格互为相反数（$+0.40$ 对 $-0.40$，$+0.45$ 对 $-0.44$，等等）。这确认了实现中没有引入非对称处理，“唯一区别是符号”在代码层面是成立的。

\textbf{成本的非对称作用是本实验最干净的发现}。以 $L=5$ 为例：毛利夏普为 $+0.40$ 与 $-0.40$（严格镜像），但净夏普为 $+0.20$ 与 $-0.60$——\emph{不是} $+0.20$ 与 $-0.20$。原因是成本是一个\emph{不区分方向的常数扣项}：它从两侧各减去约 $0.20$ 个夏普单位量级的收益（对应 1.04\% 的年化拖累 / 5.2\% 的年化波动 $\approx 0.20$）。好的一侧被削去一截，差的一侧被推向更差。\textbf{成本惩罚的是交易，而不是方向}；谁交易得更多，谁就被惩罚得更多——但当两策略交易量相同时，成本会放大两者的差距。

\textbf{短窗口的成本效率}。$L=5$ 时换手高达 52 倍/年，成本拖累 1.06\%，占毛利的一半以上；$L=250$ 时换手降至 8.83 倍/年，成本拖累 0.18\%，仅占毛利的 10\%。这直接印证了第~\ref{sec:duality}~节对偶五“换手与成本敏感度”的判断。

\subsection{期货：时序占优、截面呈 U 形}

\begin{table}[htbp]
\centering\footnotesize
\caption{中国期货：趋势与反转的夏普随窗口 $L$ 的变化}\label{tab:fut}
\begin{tabular}{@{}>{\raggedright\arraybackslash}p{1.0cm}>{\raggedright\arraybackslash}p{1.9cm}>{\raggedright\arraybackslash}p{1.9cm}>{\raggedright\arraybackslash}p{1.9cm}>{\raggedright\arraybackslash}p{1.9cm}>{\raggedright\arraybackslash}p{1.9cm}>{\raggedright\arraybackslash}p{1.9cm}@{}}
\toprule
\multirow{2}{*}{\textbf{$L$}} & \multicolumn{3}{c}{\textbf{时序口径（TS）}} & \multicolumn{3}{c}{\textbf{截面口径（XS）}} \\
\cmidrule(lr){2-4}\cmidrule(lr){5-7}
 & 趋势净 & 反转净 & 差值 & 趋势净 & 反转净 & 差值 \\
\midrule
5   & $+0.20$ & $-0.60$ & $+0.80$ & $-0.10$ & $-0.20$ & $+0.10$ \\
10  & $+0.32$ & $-0.58$ & $+0.90$ & $+0.11$ & $-0.31$ & $+0.42$ \\
20  & $+0.21$ & $-0.38$ & $+0.59$ & $-0.12$ & $-0.01$ & $-0.11$ \\
60  & $+0.21$ & $-0.32$ & $+0.53$ & $-0.82$ & $+0.72$ & $-1.54$ \\
120 & $+0.09$ & $-0.17$ & $+0.26$ & $-0.64$ & $+0.56$ & $-1.20$ \\
250 & $+0.24$ & $-0.30$ & $+0.54$ & $-0.32$ & $+0.25$ & $-0.57$ \\
\bottomrule
\end{tabular}
\end{table}

\textbf{时序口径：趋势全窗口占优，U 形未出现}。在 $L=5$ 到 $L=250$ 的全部六个窗口上，时间序列趋势策略的净夏普为正（$0.09\sim0.32$），时序反转策略的净夏普为负（$-0.17\sim-0.60$）。按美国市场的经验（短频反转、中频动量、长频反转），$L=5$ 与 $L=10$ 上本应出现反转占优，但中国期货市场并未出现。

这一结果值得单独强调：它\emph{不支持}“尺度决定归属”在中国期货时序口径上的适用性。可能的解释有三：一是中国商品期货的趋势持续性本身较强（产业资金与投机资金的行为特征）；二是 T+1 与涨跌停制度压制了短周期的反向交易；三是换手成本在短窗口上过高（$L=5$ 时 52 倍/年），足以淹没任何微弱的短期反转信号。本文无法区分这三种解释——它们是第~\ref{sec:future}~节的遗留课题 Q1。

\textbf{截面口径：U 形出现}。截面口径呈现了与直觉一致的结构：$L=10$ 时趋势占优（$+0.11$ 对 $-0.31$），$L=20$ 起反转追上，$L=60$ 与 $L=120$ 时反转显著占优（$+0.72$ 与 $+0.56$），$L=250$ 时优势收窄（$+0.25$）。趋势与反转的差值从 $+0.42$（$L=10$）单调变为 $-1.54$（$L=60$），再收窄至 $-0.57$（$L=250$）。

这是本文全部实验中\emph{唯一}出现 U 形结构的口径。它说明“尺度决定归属”不是普适定律，而是\emph{口径与市场联合决定}的结论：在有一定宽度横截面、且资产间存在共同因子的市场上，短窗口的横截面动量与中长窗口的横截面反转可以并存。

\textbf{反转侧在高 $L$ 上的经济含义}。$L=60$ 的截面反转净夏普 $0.72$、年化 $7.57\%$、最大回撤仅 $-18.74\%$，是本次实验中期货侧唯一兼具收益与稳健性的结果。它的含义是：在中国商品期货上买入过去三个月跌得最多的品种、卖出涨得最多的品种，在扣除 2 基点单边成本后仍有正的、显著的（NW $t=2.59$）收益。

\subsection{A 股：截面反转强，时序近零，短窗口被成本摧毁}

\begin{table}[htbp]
\centering\footnotesize
\caption{A 股：截面口径（5083 只个股）趋势与反转的对照}\label{tab:asu}
\begin{tabular}{@{}>{\raggedright\arraybackslash}p{1.0cm}>{\raggedright\arraybackslash}p{2.1cm}>{\raggedright\arraybackslash}p{2.1cm}>{\raggedright\arraybackslash}p{2.1cm}>{\raggedright\arraybackslash}p{2.1cm}>{\raggedright\arraybackslash}p{2.1cm}@{}}
\toprule
\textbf{$L$} & \textbf{趋势净夏普} & \textbf{反转毛利夏普} & \textbf{反转净夏普} & \textbf{反转换手} & \textbf{反转成本拖累} \\
\midrule
5   & $-2.49$ & $+1.08$ & $-0.35$ & 79.80 & 12.31\% \\
10  & $-2.11$ & $+1.10$ & $+0.10$ & 57.74 & 9.10\% \\
20  & $-2.15$ & $+1.46$ & $+0.76$ & 42.77 & 7.08\% \\
60  & $-1.93$ & $+1.46$ & $+0.99$ & 28.09 & 4.69\% \\
120 & $-1.63$ & $+1.24$ & $+0.86$ & 21.55 & 3.52\% \\
250 & $-1.17$ & $+0.86$ & $+0.56$ & 16.37 & 2.60\% \\
\bottomrule
\end{tabular}
\end{table}

\textbf{截面反转全面占优，且 U 形未出现}。A 股个股的截面动量策略在全部窗口上都是大幅负收益（净夏普 $-1.17\sim-2.49$，毛利夏普 $-0.88\sim-1.48$），其镜像的反转策略毛利夏普为 $+0.86\sim+1.46$。这与 A 股“追涨杀跌亏损、买跌卖涨获利”的市场共识一致，也是本次实验中最强的单一信号。

但注意：\emph{优势随 $L$ 的变化不是 U 形，而是倒 U 形或平台形}——反转的毛利夏普在 $L=20$ 与 $L=60$ 达到峰值（均为 $1.46$），在两端（$L=5$ 的 $1.08$、$L=250$ 的 $0.86$）下降。这与“短期反转最强”的微观结构直觉相反，说明 A 股的反转更可能是\emph{中期}现象——对应估值或情绪的回归，而非买卖价差反弹。

\textbf{短窗口被成本摧毁}。$L=5$ 时，反转的毛利夏普高达 $+1.08$，但换手 79.80 倍/年带来 12.31\% 的年化成本拖累，净夏普被压到 $-0.35$——\emph{信号完全被成本吞掉}。这是全文“成本敏感度”对偶最直接的证据：一个毛收益夏普大于 1 的信号，因换手过高而在扣除现实成本后变成负收益。随着 $L$ 增大，换手下降、成本拖累从 12.31\% 降到 2.60\%，反转的净夏普才转正并在 $L=60$ 达到 $0.99$。

这一结果对实践有直接含义：\textbf{在 A 股做统计套利型的截面反转，短窗口不可行，中期窗口（$L=20\sim120$）是唯一可行的区间}。这也解释了为什么 A 股的反转策略在实践中往往表现为“月度或季度调仓的低频价值/反转”，而不是高频统计套利。

\textbf{A 股指数时序：接近零收益，且不构成有效比较}。7 个宽基指数的时序结果（趋势净夏普 $-0.33\sim+0.17$，反转 $-0.27\sim+0.07$）全部在零附近，且最大回撤普遍在 $-50\%$ 至 $-79\%$ 之间。需要指出这组结果\emph{证据力很弱}：7 个 A 股宽基指数的相关系数高达 0.85 以上，其“组合”实质上是单一 A 股 beta 的择时，缺乏分散效应，且指数本身不可直接交易（需通过 ETF 或股指期货实现）。因此本文不把这组结果作为论据，仅作为对照列出。

\subsection{实证小结：论断的三个层次}

把四组结果合起来，可以对引言的论断给出分层的检验结论。

\textbf{第一层（信号统一）得到强支持}。期货时序口径下，趋势与反转的换手几乎完全相同（差异 $<0.1\%$）、成本拖累几乎完全相同、毛利夏普严格互为相反数。这从实证上确认了“两策略共享同一信号、唯一区别是符号”的构造性论断——它不是一个类比，而是一个可以逐项验证的事实。

\textbf{第二层（尺度决定归属）只在特定口径成立}。U 形结构只出现在期货截面口径上；在期货时序与 A 股截面上均未出现。因此这一论断应当修正为：\emph{“尺度决定归属”成立的条件是存在一个足够宽的横截面、且资产间存在共同因子}。在时序口径（单资产对自身历史）上，归属由市场特性决定，而非由尺度决定。

\textbf{第三层（成本与容量差异）得到强支持，且是实践中最重要的一层}。成本对短窗口策略的摧毁在两处得到直接印证：期货时序 $L=5$（成本占毛利 50\%）、A 股截面 $L=5$（毛利夏普 $1.08$ 被压到 $-0.35$）。而长窗口的成本占比降到 10\% 至 30\% 区间。这与第~\ref{sec:duality}~节对偶五以及第~\ref{sec:execution}~节的执行清单完全一致。

最后必须重复本章开头的定位：这组实证是\emph{变量受控的相对比较}，不是绝对收益水平的估计；由于未做多重检验校正（六个窗口 × 两方向 × 四组口径 = 48 个组合）、A 股侧做空不可实现、除权修复存在残余污染（第~\ref{sec:selfaudit}~节），本章的数字不应被当作“策略能赚多少”的答案。此外，本章全部指标基于日频收益的算术平均，未对自相关做夏普比率的标准误校正；Lo(2002) 指出在收益存在自相关时，夏普比率的样本估计方差会偏离独立同分布情形下的公式\cite{lo2002}，因此表~\ref{tab:gate}~中“成本使两侧差距不对称”的量化幅度应视为方向性判断，而非精确估计。

\section{评测污染与不可比设定审计}\label{sec:audit}

本文前半部分引用了大量文献结论。若不加审计地引用，就会把一个已知的问题带入结论：不同论文的绩效数字建立在不可比的设定上，且普遍存在数据窥探。本章按八个核验维度做审计，并\emph{同时审计本文自建回测的局限}——后者与前者同样重要，因为自建实证最容易被误当作无偏的裁判。

\subsection{审计框架}

\begin{table}[htbp]
\centering\footnotesize
\caption{八维核验口径在策略比较语境下的映射}\label{tab:audit}
\begin{tabular}{@{}>{\raggedright\arraybackslash}p{2.6cm}>{\raggedright\arraybackslash}p{11.2cm}@{}}
\toprule
\textbf{核验维度} & \textbf{本语境下的核验问题} \\
\midrule
样本数据 & 样本期、资产池、频率是否可比？是否含幸存者偏差与上市偏差？ \\
信号设定 & 形成期/持有期、阈值、参数是在样本内选的还是外生的？ \\
样本外稳健性 & 是否做 walk-forward 与跨市场泛化检验？ \\
概率/状态预测 & 是否给出参数估计与状态判断（OU 参数、半衰期、regime）？ \\
计算成本 & 数据与算力需求；重估频率 \\
许可 & 数据（CRSP/TAQ/交易所）与代码许可是否可得？ \\
公开基准 & 是否提供可复现基准（代码、复现报告）？ \\
数据污染与不可比 & look-ahead、数据窥探、多重检验、跨论文设定差异 \\
\bottomrule
\end{tabular}
\end{table}

\subsection{幸存者偏差与上市偏差}

两条脉络在这一项上的暴露程度不同。

趋势研究的资产池通常是\emph{存续至今}的期货合约或指数成分。Hurst 等(2017) 手工收集 1880 年以来的 CBOT 数据，所用 67 个市场是“可回溯到”的市场\cite{hop2017}；那些中途退市、停止交易的市场未被纳入。这构成典型的幸存者偏差：策略在“活下来的市场”上表现好，可能是选择的结果。

反转研究的偏差更隐蔽。配对交易需要在形成期就能配对，而“能配对”本身隐含了样本的存续要求。Gatev 等(2006) 用 1962--2002 年 CRSP 数据，其配对的股票池是存续股票\cite{gatev2006}；Do \& Faff(2012) 延伸样本至 2009 年后利润衰减，部分原因可能就是早期样本的存续偏差在后期被稀释\cite{dofaff2012}。

此外，反转策略还多一层偏差：\emph{它需要做空}。在多数市场中，可做空的股票集合远小于全部股票，且做空成本与可借量随时间变化。文献中的反转收益大多在“可做空”的假设下计算，实际可实现的收益更低。这一项在第~\ref{sec:empirical}~节的 A 股部分尤其致命，见后文审计。

\subsection{Look-ahead 偏差}

趋势研究的典型 look-ahead 来源是\emph{信号与执行的时间错配}：若在信号形成期的最后一个交易日收盘价上建仓，而实际执行发生在次日开盘，则回测高估收益。MOP12 的构造以月末信号、次月持有处理了这一项\cite{mop2012}。

反转研究的典型来源是\emph{配对选择与交易期重叠}。若配对是在包含交易期的数据上选出的，则价差平稳性被机械地保证。标准做法是形成期与交易期严格分离并滚动（Gatev 等的 6+6 个月滚动设计\cite{gatev2006}）。

本文自建回测的处理是：信号使用至 $t$ 日收盘的全部可得信息，仓位自 $t+1$ 日起生效，持有期内按收益漂移，调仓以目标权重与实际权重的偏离计换手。这一处理避免了最常见的 look-ahead，但仍有残余问题（见~\ref{sec:selfaudit}~小节）。

\subsection{数据窥探与多重检验}

这是两条脉络共同的、也是最严重的威胁。

\textbf{趋势一侧的历史更早被审查}。Sullivan 等(1999) 对 7846 条技术交易规则做 bootstrap 校正，发现样本外超额收益不再显著\cite{sullivan1999}；White(2000) 的 Reality Check 给出通用框架\cite{white2000}；Hansen(2005) 的 SPA 检验与 Romano--Wolf(2005) 的逐步多重检验进一步改进\cite{hansen2005,romanowolf2005}。Harvey 等(2016) 提出的更高门槛（新因子应达到 $t>3$ 左右）\cite{hlz2016}、以及 Bailey 与 López de Prado 的紧缩夏普比率与回测过拟合概率\cite{bailey2014,bailey2016}，构成当前的审计标准。

\textbf{反转一侧的暴露更隐蔽}。配对交易的超参数包括选对规则、协整显著性、半衰期上下限、阈值带宽度、形成期长度；每一个都是自由度。Krauss(2017) 的综述在梳理该领域时明确指出方法选择的空间巨大\cite{krauss2017}。这意味着反转文献中“发现显著收益”的比例可能被系统性高估，而 Do \& Faff(2012) 的利润衰减正是这一高估在样本外被修正的表征\cite{dofaff2012}。

\textbf{本文自建回测的暴露}。本文对 $L\in\{5,10,20,60,120,250\}$ 共六个窗口、两个方向、两种口径（时序/截面）、两个市场做了完整网格扫描，未做多重检验校正。因此本文报告的“某组合夏普为正”不应当被读作统计显著，只能读作\emph{在给定口径下的相对比较}。这是本文实证的定位：它是变量对照，不是显著性检验。

\subsection{样本外与发表后衰减}

McLean \& Pontiff(2016) 对 97 个已发表因子的实测是这一项最有力的证据：组合收益在样本外-发表前下降 26\%，发表后累计下降约 58\%\cite{mcleanpontiff2016}。这一结论对两条脉络同等适用——趋势与反转都是“被发表过的异象”。

对趋势策略的特别审查来自 Huang 等(2020)：逐资产回归显示时间序列动量的证据在样本内外均薄弱，池回归的大 $t$ 值不稳健，且策略表现与“按历史样本均值投资”的简单规则几乎相同\cite{huang2020}。Pitkäjärvi 等(2020) 从跨资产信号角度给出相关质疑\cite{pitkajarvi2020}。反方证据是 Hurst 等(2017) 的百年样本\cite{hop2017}——两方证据的张力在结论置信度章（第~\ref{sec:confidence}~节）中作为专门条目处理。

对反转策略的特别审查是利润的时间衰减：Gatev 等(2006) 报告的约 11\% 年化超额收益基于 1962--2002 样本\cite{gatev2006}，Do \& Faff(2012) 显示 1989 年之后明显下降\cite{dofaff2012}。这提示反转文献的高收益不能外推。

\subsection{跨论文不可比设定}

即便不考虑偏差，跨文献的绩效数字也常常不可比。差异来源至少四类：

\begin{enumerate}[leftmargin=2.2em]
  \item \textbf{样本期与资产池}：MOP12 用 1960s--2009 的 58 个合约\cite{mop2012}，HOP17 用 1880--2016 的 67 个市场\cite{hop2017}，Georgopoulou \& Wang(2017) 用 1969--2015 的股票与商品\cite{georgopoulou2017}；三者的“趋势跟踪年化”不是同一个量。
  \item \textbf{波动率缩放是否启用}：MOP12 的原始设定与 Barroso--Santa-Clara(2015)、Harvey 等(2018) 的缩放设定在风险暴露上不同\cite{mop2012,barroso2015,harvey2018}；Kim 等(2016) 指出收益的很大部分来自缩放本身\cite{kim2016}。
  \item \textbf{成本与费用假设}：HOP17 计入滚动成本与管理费\cite{hop2017}，部分复现研究忽略；反转一侧的成本假设差异更大（价差、融券费、冲击模型各不相同）\cite{dofaff2012,novymarxvelikov2016}。
  \item \textbf{频率}：月频、周频、日频策略的换手相差一个量级，成本敏感度随之变化\cite{baltaskosowski2011,korajczyksadka2004}。
\end{enumerate}

因此本文在引用任何绩效数字时均注明其设定条件，并在比较矩阵中拒绝把不同设定的数字直接排序。

\subsection{本文自建回测的局限}\label{sec:selfaudit}

自建实证同样需要审计，且这一小节是本文最应当被严格对待的部分。

\textbf{局限一：A 股个股采用不复权价的近似修复，存在残余污染}。通达信本地日线为不复权数据，本文的修复规则是——主板个股单日收益低于 $-10.5\%$、创业板/科创板低于 $-20.5\%$ 者判定为除权并置该日收益为零（等价于按当日比例对更早价格做前复权）。这一处理能捕获绝大多数送转除权（跳空幅度大于涨跌停），但\emph{漏掉跳空幅度小于涨跌停的除权}（如小额分红、10 送 1），后者构成残余污染。样本内共识别并修复 6121 处确定性除权；未识别的部分未做定量估计，属于本文已知的未修正偏差。

\textbf{局限二：A 股截面策略的做空假设不现实}。A 股个股的融券做空渠道有限、成本高、可借券池小。本文的 A 股截面多空组合因此在多数时段不可实现，其结果应被视为\emph{上限估计}而非可交易绩效。这一局限对反转策略的影响大于趋势策略，因为反转的多空结构更依赖空头端。

\textbf{局限三：未建模涨跌停与停牌的可交易性}。涨跌停板上无法成交，停牌期间无法调仓。本文的信号与收益计算未做这一过滤，因此在高波动时段（如 2015 年 A 股异常波动期间）会高估策略的可实现收益。期货侧的涨跌停同样未过滤。

\textbf{局限四：成本模型是简化的固定基点}。本文对期货采用单边 2 个基点、A 股采用单边 15 个基点（含佣金、印花税与滑点），未采用参与率约束的平方根冲击模型，也未考虑流动性时变。这一简化会低估反转策略（高换手）的成本，因为反转交易密集发生在流动性较差的时点。

\textbf{局限五：日线口径}。所有策略在日线频率上构建，无法刻画日内反转与执行细节；这与反转策略的实际高频属性存在口径落差。

\textbf{局限六：样本期较短且单一市场}。2014--2026 的十二年间，中国期货与 A 股的制度环境经历了显著变化（股指期货限制、融券政策调整、注册制改革）；结论不应外推到更早时期或其他市场。

这些局限不改变本文实证的定位——它是一次\emph{变量受控的相对比较}（同一资产池、同一成本、仅翻转符号），而不是对绝对收益水平的估计。任何把本文数字当作“策略年化收益”的读法都是误读。

\subsection{审计小结}

两条脉络的污染暴露在结构上同构：都有数据窥探与发表后衰减，都有不可比的样本设定，都有做空与流动性的实现障碍。差异在于侧重点——趋势研究更早、更彻底地接受了多重检验审查；反转研究的自由度（配对选择、参数筛选）更多，因此更依赖于“利润衰减”这类事后证据来暴露问题。

对本文自建实证的审计结论是：它可以支撑\emph{方向性比较}（两策略在同一口径下的相对表现），不能支撑\emph{水平性结论}（某策略能赚多少）。

\section{结论置信度分级}\label{sec:confidence}

本章对全文的核心结论给出置信度分级。分级不是修辞上的“我觉得”，而是按四条可操作的标准判定：\textbf{(i)} 证据是否来自多个独立来源；\textbf{(ii)} 是否经受过多重检验或样本外检验；\textbf{(iii)} 是否存在未解决的反方证据；\textbf{(iv)} 结论是机制判断还是可观测事实。四条标准中，“可观测事实”类的结论天然更容易达到高置信度，“机制解释”类的结论上限较低——因为机制不可直接观测（第~\ref{sec:mechanism}~节）。

\subsection{分级结果}

\newcommand{\lvA}{\textbf{高}}
\newcommand{\lvB}{\textbf{中}}
\newcommand{\lvC}{\textit{低}}

\begin{longtable}{@{}>{\raggedright\arraybackslash}p{3.6cm}>{\raggedright\arraybackslash}p{1.1cm}>{\raggedright\arraybackslash}p{4.3cm}>{\raggedright\arraybackslash}p{4.6cm}@{}}
\caption{核心结论的置信度分级}\label{tab:conf}\\
\toprule
\textbf{结论} & \textbf{置信度} & \textbf{支撑理由} & \textbf{限制与反方} \\
\midrule
\endfirsthead
\toprule
\textbf{结论} & \textbf{置信度} & \textbf{支撑理由} & \textbf{限制与反方} \\
\midrule
\endhead
\bottomrule
\endlastfoot
趋势与均值回归共享同一信号，唯一自由度是符号 & \lvA & 定义为可验证的构造事实；文献中的标准构造逐项对应\cite{mop2012,jegadeesh1993,gatev2006,lehmann1990} & 协整型反转含额外的结构约束，非纯粹符号翻转\cite{englegranger1987} \\
收益自相关随持有期变号 & \lvA & 多源独立证据：短频负\cite{lehmann1990,jegadeesh1990}、中频正\cite{jegadeesh1993,mop2012}、长频负\cite{debondt1985,poterbsummers1988}；方差比检验提供统一统计量\cite{lomackinlay1988} & 长周期负自相关受小样本质疑\cite{richardson1993}；本文 A 股实证未复现短频反转（见第~\ref{sec:empirical}~节） \\
趋势策略呈凸性、危机中获利 & \lvB & 期权复制\cite{funghsieh2001}、百年样本十年全正\cite{hop2017} & 动量崩溃显示深度左尾\cite{danielmoskowitz2016}；百年样本存在幸存者偏差 \\
反转策略呈凹性、收敛失败致灾 & \lvB & 两起独立的重大事件（LTCM 1998、量化危机 2007）\cite{lowenstein2000,jorion2000,khandanilo2011} & 事件样本少（n=2），难以做统计检验；样本选择可能偏于戏剧性案例 \\
反转策略换手高、成本敏感 & \lvA & 理论（带宽 $\propto$ 成本立方根\cite{martin2011}）与实证（利润衰减\cite{dofaff2012}、成本分类\cite{novymarxvelikov2016,frazzini2018}）一致 & 低频价值型反转换手极低，不适用本条 \\
趋势策略面临拥挤约束 & \lvB & 理论\cite{stein2009}与拥挤度量实证\cite{loupolk2022,baltas2019} & 拥挤度的定量估计尚无统一口径；容量上限的估计分歧较大 \\
波动率目标化对两策略作用不对称 & \lvB & 趋势侧多源验证\cite{barroso2015,harvey2018,moreiramuir2017} & 反转侧“波动率目标有害”缺少同等的系统实证，主要是机制推理 \\
两策略均受数据窥探威胁 & \lvA & 方法论文献高度一致\cite{sullivan1999,white2000,hlz2016,mcleanpontiff2016}；校正工具成熟\cite{hansen2005,romanowolf2005,bailey2014} & 校正后的结论仍依赖模型假设 \\
中周期时间序列动量真实存在 & \lvB & MOP12 的 58 合约证据\cite{mop2012}、百年样本\cite{hop2017} & 逐资产再检验质疑其稳健性\cite{huang2020,pitkajarvi2020}；矛盾未解决 \\
阈值带宽度随成本立方根变化 & \lvB & 解析结果\cite{martin2011}与数值验证\cite{boguslavsky2004,zhangzhang2008} & 假设较理想（线性成本、OU 动态），实盘偏离未知 \\
机制解释（行为 vs 风险）的归属 & \lvC & 两族解释均有模型支撑\cite{hongstein1999,bsv1998,dhs1998,hop2017} & 机制不可直接观测，无从判决\cite{cochrane2011,cochrane2005} \\
本文实证支持“持有期决定归属” & \lvC & 期货侧时序趋势优于时序反转，呈单调结构 & 仅为单一市场、单一时期、未做多重检验校正；不构成统计显著结论 \\
\end{longtable}

\subsection{分歧最深的一条}

表中“中周期时间序列动量真实存在”是整个比较中最分裂的结论。正方是 Moskowitz 等(2012) 的 58 个合约与 Hurst 等(2017) 的百年样本\cite{mop2012,hop2017}，反方是 Huang 等(2020) 的逐资产再检验与 Pitkäjärvi 等(2020) 的跨资产质疑\cite{huang2020,pitkajarvi2020}。本文不仲裁，只指出分歧的来源：Huang 等发现 TsMOM 的盈利与“按历史均值投资”几乎无差别，这一批评的实质是\emph{把“可预测性”与“策略盈利”分开}——策略可以盈利而预测性不显著，因为波动率缩放本身贡献收益（见 Kim 等 2016\cite{kim2016}）。因此这条结论应表述为：“中周期趋势策略在多数样本中盈利，但其中有多少来自时序预测、多少来自风险缩放，尚未定论。”

\subsection{一个整体的判断}

把分级结果汇总，可以看到一条清晰的分界：\emph{关于“是什么”的结论置信度高，关于“为什么”和“值多少”的结论置信度低}。

“两策略共享信号”“自相关随尺度变号”“反转换手更高”——这些是构造事实或稳健的经验规律，置信度高。
“趋势收益来自风险溢价”“反转收益来自流动性补偿”——这些是机制解释，置信度低。
“某策略年化收益 X\%”——跨文献不可比，且受发表后衰减影响，置信度低。

这条分界本身是本文最希望读者带走的东西。它意味着：\textbf{策略比较应当建立在高置信度的结构事实上（信号结构、尺度依赖、成本与容量约束），而对机制叙事与收益水平保持低信任。}

\section{可复现实验路线图}\label{sec:reproduce}

前文的比较依赖两类证据：文献（不可比、有污染）与本文自建回测（样本单一、口径简化）。要把比较推到更牢固的基础上，需要一组\emph{口径统一、可复现、带验证门}的实验。本章给出这样的路线图，按 Phase 0 至 Phase 4 划分，每阶段列明交付物与验证门。

设计原则有三条：\textbf{(i) 变量受控}——两策略共享除信号符号外的全部设定；\textbf{(ii) 口径前置}——成本、样本、指标在实验开始前锁定，不得事后调整；\textbf{(iii) 门禁明确}——每阶段设通过条件，未通过则不得进入下一阶段。

\subsection{Phase 0：基建与口径锁定}

\textbf{交付物}
\begin{itemize}[leftmargin=2.2em]
  \item 统一数据层：同一资产池的日频与分钟频价格、持仓量、成交额；股票需含复权因子与停复牌标记；期货连续合约需含换月跳空守门与前复权处理。
  \item 统一成本层：参与率约束的冲击模型（冲击与参与率平方根成正比，设定参与率上限并对超限标记）；股票端区分佣金、印花税与借券费。
  \item 统一评估层：13 项强制指标（年化收益、年化波动、夏普、最大回撤、Calmar、偏度、超额峰度、年化换手、毛利年化、成本拖累、月胜率、Newey--West $t$ 统计量、成本占毛利比）。
  \item 策略骨架：两套策略共用同一回测引擎，信号函数以符号参数区分方向。
\end{itemize}

\textbf{验证门}
\begin{itemize}[leftmargin=2.2em]
  \item \textbf{G1 一致性}：同一信号、符号翻转后，毛收益序列应严格互为相反数（差异应为零）。这是统一框架的直接检验——若不能通过，说明实现中有非对称处理。
  \item \textbf{G2 成本守恒}：两策略在同一调仓频率下的换手应完全相同；成本拖累应完全相等。若趋势与反转的换手不同，说明信号平滑或阈值处理引入了非对称。
  \item \textbf{G3 无 look-ahead}：在信号计算中注入未来一日数据，回测指标应出现可观测的改善；以此为哨兵检验。
\end{itemize}

G1 与 G2 是本路线图的核心设计。它们在本文自建回测中已经得到验证：期货时序口径下，趋势与反转的换手几乎完全相等（如 $L=5$ 时分别为 52.41 与 52.38），成本拖累一致，净值差异完全来自信号符号——这正是“唯一区别是符号”的实证体现。

\subsection{回测代码框架：五个模块}

Phase 0 的四项交付物落在同一个代码框架里。框架分五层，两套策略共享除信号符号外的全部代码——策略差异被隔离在信号层这一个函数内，G1 与 G2 才能用机器自动校验，而不是靠人工比对。

\begin{quote}\ttfamily\small
数据层：\\
\quad 输入日频行情、复权因子、停复牌标记、期货换月记录\\
\quad 输出统一的收益序列与可交易标记；换月跳空并入收益，停牌日标记为不可交易\\[2pt]
信号层（策略唯一分岔点）：\\
\quad 输入回看窗口 $L$ 与方向符号 $\sigma\in\{+1,-1\}$\\
\quad $s_t \gets \sigma\times\sum_{k=0}^{L-1}\ln(1+r_{t-k})$ \hfill \# 趋势 $\sigma=+1$，反转 $\sigma=-1$\\[2pt]
引擎层（逐日推进，先信号后仓位）：\\
\quad 对每个交易日 $t$：\\
\qquad 用截至 $t$ 日收盘的全部可得信息计算 $s_t$\\
\qquad 目标权重 $w_t \gets$ 缩放与归一化（两策略共用同一函数）\\
\qquad 仓位自 $t+1$ 日起生效，持有期内按收益漂移更新权重\\
\qquad 换手 $\gets\sum_i\bigl|w_{i,t+1}-w^{\mathrm{current}}_{i,t}\bigr|$，仅对换入换出部分计成本\\[2pt]
成本层：\\
\quad 固定基点部分期货单边 2、A 股单边 15；\\
\quad 另设参与率上限 10\% 的冲击项，超限标记该笔不可成交\\[2pt]
评估层：\\
\quad 输出 13 项强制指标；同时输出扣成本前的毛收益序列，供门禁校验\\
\end{quote}

三道哨兵直接挂在引擎层与评估层上，每次跑完自动触发：G1 对 $\sigma=+1$ 与 $\sigma=-1$ 各跑一遍比对毛收益序列，G2 比对两次运行的逐日换手与成本，G3 用信号层的未来数据注入做 look-ahead 探测（定义见 Phase 0）。三道全过，框架才可用于后续实验——任何一道失败，都先修实现再谈结果。

\subsection{Phase 1：尺度扫描与 U 形检验}

\textbf{交付物}：在统一口径下对 $L\in\{5,10,20,60,120,250\}$ 做完整网格，覆盖时序（TS）与截面（XS）两种口径、期货与股票两个市场，输出“$L$—夏普”曲线。

\textbf{验证门}
\begin{itemize}[leftmargin=2.2em]
  \item \textbf{G4 U 形存在性}：至少在一种口径与一个市场上，趋势与反转的相对夏普应随 $L$ 反转（短 $L$ 反转占优、中 $L$ 趋势占优、长 $L$ 反转占优）。若在任何口径上都不出现，则需重新审视“尺度决定归属”的普适性。
  \item \textbf{G5 成本归因}：在同一 $L$ 上，趋势净夏普 $+$ 反转净夏普 应约等于 $-2\times$(成本拖累/波动)。这是成本非对称性的定量检验。
\end{itemize}

\subsection{Phase 2：稳健性与多重检验}

\textbf{交付物}
\begin{itemize}[leftmargin=2.2em]
  \item 样本外检验：滚动 walk-forward，形成期与交易期严格分离。
  \item 多重检验校正：对全部网格组合计算 White Reality Check、Hansen SPA 与紧缩夏普比率；报告校正前后的显著性。
  \item 参数稳健性：$L$ 邻域内的夏普曲面平滑度；是否存在孤立的高夏普点（过拟合征兆）。
\end{itemize}

\textbf{验证门}
\begin{itemize}[leftmargin=2.2em]
  \item \textbf{G6 多重检验后存活}：至少一条“趋势优于反转”或“反转优于趋势”的结论在校正后仍显著。
  \item \textbf{G7 稳健性}：最优 $L$ 的邻域夏普不得出现剧烈跳变；若最优点的夏普显著高于邻域均值，判定为过拟合。
\end{itemize}

\subsection{落地流程：从网格到判决的四道闸门}

把 Phase 2 的第二项交付物展开成可执行的流程。方法学本身（Reality Check、SPA、紧缩夏普）已在第~\ref{sec:audit}~节作为文献标准介绍过，这里只回答操作问题：先算什么、后算什么、哪一步否决哪一步。

\begin{enumerate}[leftmargin=2.2em]
  \item \textbf{锁定试验族}。网格的全部组合（第~\ref{sec:empirical}~节的六窗口 × 两方向 × 四组口径）在实验开始前一次性定义为试验族。族内每个组合都是潜在的“被选中者”，事后追加的组合不得并入——并入即扩大族，先前的校正随之失效。这一步是整个流程的地基，也是最容易被跳过的一步。
  \item \textbf{基准统计量}。对族内每个组合计算净夏普与 Newey--West 调整的 $t$ 统计量（收益自相关下常规标准误失效，见第~\ref{sec:empirical}~节的口径声明），记族内最优者的原始统计量为 $t_{\max}$。到这一步为止，一切与单组合检验无异；$t_{\max}$ 本身不构成证据——它是 48 次搜索的最大值。
  \item \textbf{跨组合校正}。对 $t_{\max}$ 做 block bootstrap 重抽样（保持序列相关的块长），得到 White Reality Check 与 Hansen SPA 的 $p$ 值，回答“最优组合是否只是搜索的产物”；再用 Romano--Wolf 逐步法定位校正后仍显著的组合\emph{集合}——集合可能为空，空集意味着没有任何单一组合在族水平上站得住。
  \item \textbf{稳健性折算与判决}。对通过第 3 步的组合报告紧缩夏普（把估计误差从原始夏普中折掉）与回测过拟合概率（CSCV 口径，0.5 等价于抛硬币）。判决对接 G6：RC 与 SPA 的 $p<0.05$、且过拟合概率明显低于 0.5，判“校正后存活”；任一不满足，该结论降级为方向性描述，只能进比较矩阵的置信度列，不能进结论。
\end{enumerate}

四道闸门的顺序不可调换：先锁定族再算统计量，否则校正对象不确定；先跨组合校正再折算稳健性，否则折算的是一个未经检验的候选。流程跑完输出两张表——校正前后的显著性对照，与被否决组合的清单；后者和前者同样重要，它记录了搜索代价。

\subsection{Phase 3：执行层对照}

\textbf{交付物}
\begin{itemize}[leftmargin=2.2em]
  \item 波动率目标化的不对称检验：对两策略分别施加与不施加波动率目标，比较夏普、偏度与最大回撤的变化。
  \item 资金约束模拟：为反转策略设定保证金缓冲与偏离上限，测试在压力情景（流动性骤降）下的生存性。
  \item 组合融合检验：按尺度分层、按状态加权、按风险预算三种融合方案与简单等权叠加的对比。
\end{itemize}

\textbf{验证门}
\begin{itemize}[leftmargin=2.2em]
  \item \textbf{G8 不对称性确认}：波动率目标化应改善趋势策略的风险调整收益，但不应同等改善（或应恶化）反转策略。若两者同等改善，则第~\ref{sec:execution}~节的核心执行差异判断需修正。
  \item \textbf{G9 融合增益}：至少一种融合方案的夏普或尾部指标优于两策略单独持有与简单叠加。
\end{itemize}

\subsection{Phase 4：沉淀与复现包}

\textbf{交付物}：完整代码仓库（数据层、成本层、引擎、策略、评估），含环境锁定文件、随机种子、以及一键复现脚本；实验记录含每次运行的参数与结果哈希。

\textbf{通过条件}：第三方在无作者协助下，用同一复现包得到同一组指标（允许浮点误差）。

\subsection{与本文实证的关系}

本文第~\ref{sec:empirical}~节的自建回测可以视为本路线图的 Phase 0 与 Phase 1 的初步实现：它通过了 G1、G2（期货时序口径下两策略换手几乎相等），并部分回答了 G4（在期货截面口径上观察到 U 形结构）。Phase 2 至 Phase 4 尚未执行，这是本文实证的未完成部分，也是后续工作的直接入口。

\section{研究现状、遗留课题与未来方向}\label{sec:future}

\subsection{当前研究现状}

把两条脉络的最新进展放在一起看，有四个趋势值得注意。

\textbf{其一，两条脉络正在方法上合流}。最直接的证据是“慢动量加快反转”这类频段融合工作\cite{wood2021}。更广泛的迹象是机器学习方法的介入：Gu 等(2020) 的资产定价实证显示，树模型与神经网络的收益预测能力显著优于线性模型，且动量类变量在其中扮演重要角色\cite{gu2020}。当模型不再预设“动量或反转”，而是从数据中学习响应函数的形状（甚至允许它随状态变化）时，二者作为\emph{两个候选假设}的地位被削弱，变成\emph{同一个响应函数在不同状态下的取值}。

\textbf{其二，反转一侧正在从“选对”转向“构造组合”}。早期配对交易的核心问题是选对（距离法、协整法）；Cuturi 与 d'Aspremont(2013) 的方差阈值方法把问题改写为：直接求解使均值回归程度最大化、同时约束方差的线性组合\cite{cuturi2013}。这一转向的实质是承认“配对”只是构造均值回归载体的一种启发式，更一般的问题是优化式地合成这样的载体。近年 OU 框架的工程化复现（如 arXiv:2412.12458）与深度学习统计套利（arXiv:2106.04028）延续了这一路线\cite{suchato2024,guijarro2021}；强化学习方案的引入（如 arXiv:2403.12180）则试图把仓位管理与信号生成放进同一学习目标\cite{ning2024}。

\textbf{其三，趋势一侧的重心转向容量与拥挤}。信号构造问题在 2010 年代基本定型（形成期/持有期、波动率缩放、信号平滑），此后文献的注意力明显移向容量与拥挤：Stein(2009) 的理论、Lou--Polk(2022) 的 comomentum 度量、Baltas(2019) 的拥挤综述\cite{stein2009,loupolk2022,baltas2019}。这反映了一个现实：趋势策略的信号收益难以继续提升，制约它的是“钱太多”而非“信号太弱”。

\textbf{其四，两策略都进入了“事后审计”时代}。紧缩夏普比率、回测过拟合概率、多重检验校正已从学术规范变成行业惯例\cite{bailey2014,bailey2016,hlz2016,harveysliu2015}。这对两个策略是同等的要求。

\subsection{遗留课题}

以下课题在本文的证据梳理中反复出现，且都未解决。

\textbf{Q1：U 形自相关是否跨市场普适？} 本文的对照实证给出了一个值得警惕的反例：在中国期货的\emph{时序}口径上，趋势策略在所有 $L$ 上都优于反转策略，“短尺度反转”未出现；U 形结构只出现在\emph{截面}口径上（短 $L$ 趋势占优、中长 $L$ 反转占优）。这说明“尺度决定归属”的普适性依赖具体口径与市场。其成因（投资者结构、涨跌停制度、T+1 交易）尚缺系统研究。

\textbf{Q2：趋势收益中“预测性”与“风险缩放”各占多少？} Kim 等(2016) 指出时间序列动量的收益很大部分来自波动率缩放本身\cite{kim2016}；Huang 等(2020) 发现其表现与简单历史均值规则几乎相同\cite{huang2020}。把这两个成分干净地分解，是判定“趋势溢价是否真实”的关键，但尚无公认的分解方案。

\textbf{Q3：拥挤度能否事前度量？} comomentum 属于滞后代理（从已实现相关性推断），无法事前预警；领先代理（资金流、持仓集中度）数据可得性差。一个可事前计算、且与未来收益衰减确有预测关系的拥挤指标，至今没有。

\textbf{Q4：反转策略的容量上限能否给出定量估计？} 与趋势不同，反转的约束来自自身成本，理论上更可测算（带宽与成本的三次方根关系提供了一个起点\cite{martin2011}）。但从理论带宽到实际可管理规模之间缺少完整的映射。

\textbf{Q5：两策略相关性的状态依赖如何建模？} 二者在某些状态下高度负相关（信号相反），在资金冲击下正相关（同为去杠杆受害者）。这一“相关性变号”现象对组合构建是致命的，但目前多被当作经验事实处理，缺少结构模型。

\textbf{Q6：非线性冲击成本如何改变比较结论？} 本文的成本模型是固定基点。反转策略的交易密集发生在流动性差的时点，其真实冲击成本高于平均值；若采用参与率约束的平方根模型，反转的净收益会被进一步压低。这一敏感性尚未被系统量化。

\textbf{Q7：做空约束下反转策略的可实现性？} 反转的多空结构依赖空头端，而多数新兴市场（含 A 股）做空困难。纯多头版本的反转策略（只买输家）在收益与风险上如何，缺少系统研究。

\textbf{Q8：机器学习是否改变了两策略的本质区分？} 若模型可以自由学习响应函数形状，“趋势”与“反转”退化为两个特例。那么二者的区分是否还具有结构意义，取决于响应函数在实际数据中是否呈现稳定的分段结构——这是一个可检验但尚未被检验的问题。

\textbf{Q9：机制之争能否被判决？} 目前的风险解释与行为解释在价格数据上难以区分\cite{cochrane2011}。一个可能的出路是寻找\emph{只在其中一种机制下成立}的辅助预测（如资金流数据、投资者结构变化），但目前尚无成功案例。

\textbf{Q10：极端资金冲击下两策略同时失效，如何对冲？} 2007 年 8 月的教训是：冲击来自资金端，分散化失效\cite{khandanilo2011,brunnermeier2009}。一个能在资金冲击下真正分散的机制（如显式的流动性保险、跨市场分散、或有减仓规则）尚未被提出。

\subsection{未来方向}

\textbf{方向一：频段分解的统一框架}。把收益序列按时间尺度分解（小波、经验模态分解或频域滤波），分别在每个频段上估计自相关符号与幅度，从而把两策略从“二选一”变成“多频段风险预算”。这一方向在文献中已有雏形\cite{wood2021}，但缺少系统化的理论与实证。它也是本文统一框架最自然的延伸。

\textbf{方向二：状态感知的动态融合}。用 regime 切换模型（Hamilton 1989；Ang--Timmermann 2012）估计“趋势占优”与“反转占优”的状态概率，据此动态分配两策略的权重\cite{hamilton1989,angtimmermann2012}。关键在于状态变量必须\emph{可事前观测}，否则退化为事后标定。

\textbf{方向三：信号—成本—资金的联合优化}。Gârleanu 与 Pedersen(2013) 给出了含成本的最优动态组合解\cite{garleanu2013}，但未纳入资金约束与保证金机制。把三者纳入同一优化问题，对反转策略尤其重要——它的成败往往由资金端决定。

\textbf{方向四：中国市场的制度性比较研究}。中国市场的制度特征（T+1、涨跌停、融券受限、投资者结构以散户为主）为检验“两策略差异的制度依赖性”提供了独特样本。本文的对照实证只是初步观察，其成因（尤其是时序口径下短反转的缺失）值得专门研究。

\textbf{方向五：事前拥挤度量}。从持仓集中度、资金流、期权隐含信息等可得数据构造可事前计算的拥挤代理，并检验其与趋势策略未来收益衰减的关系。这是把第~\ref{sec:duality}~节“容量约束”对偶从定性判断推进到定量工具的关键。

\subsection{全文结论}

本文从一句直觉出发——“趋势与反转只是同一信号的两种符号选择”——做了三件事：把两条脉络编年并置，给出统一与对偶的数学框架，并用统一口径的对照实证检验核心预测。

主要的结论可以压缩成四点。

\textbf{第一，统一是真的，但只统一在信号层}。两策略共享同一信号（过去收益的符号或对均值的偏离），共享同一数学容器（带符号回归强度的 OU 族），共享同一统计量（方差比/自相关函数）。凡涉及“用什么信息、赌什么量”的问题，二者是同一个东西。

\textbf{第二，差异是真的，而且集中在执行层}。平稳性假设、收益偏度、时间尺度、风险来源、容量与换手、失效模式、入场门槛——七组对偶中没有一组是修辞性的，它们都可以追溯到具体的工程后果：反转要拓宽带宽、要做因子剥离、要留资金冗余；趋势要平滑信号、要做状态过滤、要监控拥挤度。

\textbf{第三，成本是两策略最不对称的一把尺子}。本文的对照实证给出了一个干净的示范：在期货时序口径、$L=5$ 的设定下，趋势与反转的换手几乎完全相同（52.41 对 52.38），成本拖累完全相等；但由于二者毛收益互为相反数，成本这一“常数扣项”让表现好的一侧被削去一截、表现差的一侧被推向更差——净夏普分别为 $+0.20$ 与 $-0.60$，而非预期的 $+0.4$ 与 $-0.4$。这不是数据噪声，而是成本对符号的必然作用：\emph{成本不区分方向，它只惩罚交易}。谁更频繁地交易，谁就更多地被惩罚。

\textbf{第四，机制叙事应当被降权}。关于“风险还是行为”的争论无法从价格数据中裁决，它会影响我们对策略稳健性的先验，但不应当成为选择策略的依据。可依据的是可观测的结构事实：自相关在哪个尺度上显著、成本与容量的实测约束、资金端的状态。

回到最初的问题——两个策略的本质异同是什么？本文的答案是：\textbf{它们赌的是同一件事（收益的符号），只是押在了符号的两端；它们的分野不在信号里，而在信号之后的每一步——验证、缩放、对冲、交易、融资}。把这句话说清楚的实际价值是：当有人问“该用趋势还是反转”时，正确的问题不是“你相信市场会走远还是会回头”，而是“你要在哪个尺度上、以什么成本、承担什么样的资金约束，去押注哪个符号”。

\appendix

\section{数学工具备查}\label{sec:appendix}

正文多次使用 OU 过程、方差比与最优阈值带这三个工具，但只给出结论。本附录补齐推导要点与使用边界，使读者可以独立复现本文的数学论断。

\subsection{OU 过程的离散估计与半衰期}

连续时间的 OU 过程
\begin{equation}
dX_t=\theta(\mu-X_t)\,dt+\sigma\,dW_t
\end{equation}
在等间隔 $\Delta$ 上的精确解是
\begin{equation}
X_{t+\Delta}=\mu\left(1-e^{-\theta\Delta}\right)+e^{-\theta\Delta}X_t
+\sigma\sqrt{\frac{1-e^{-2\theta\Delta}}{2\theta}}\;\varepsilon_t,\qquad
\varepsilon_t\sim \mathcal N(0,1).
\end{equation}
这是一个 AR(1)：
\begin{equation}
X_{t+\Delta}=c+\phi X_t+u_t,\qquad
\phi=e^{-\theta\Delta},\quad
c=\mu(1-\phi),\quad
\mathrm{Var}(u_t)=\frac{\sigma^2}{2\theta}\left(1-e^{-2\theta\Delta}\right).
\end{equation}
因此参数的估计可以分两步完成：先用最小二乘或极大似然估计 $(\hat c,\hat\phi)$，再由
\begin{equation}
\hat\theta=-\frac{\ln\hat\phi}{\Delta},\qquad
\hat\mu=\frac{\hat c}{1-\hat\phi},\qquad
\hat\sigma^2=\frac{2\hat\theta\,\widehat{\mathrm{Var}}(u_t)}{1-e^{-2\hat\theta\Delta}}
\end{equation}
反解连续时间参数。

\textbf{半衰期}由“偏离被拉回一半所需时间”定义，即解 $e^{-\theta H}=1/2$：
\begin{equation}
H=\frac{\ln 2}{\theta}.
\end{equation}
在离散口径下等价于 $H=\Delta\ln2/\ln(1/\phi)$。这一量在配对交易中是核心参数：它同时决定了持仓的期望时长与最优交易带的宽度量级。

\textbf{使用边界}。三条常见误用值得点明。其一，$\hat\phi$ 在小样本下有向下偏误（Kendall 1954 提出过修正），因此短样本估出的 $\theta$ 往往偏大、半衰期偏短；Poterba 与 Summers(1988)、Fama 与 French(1988) 的长期反转结论之所以受到 Richardson(1993) 的质疑，根本原因正是长周期回归在小样本下的检验功效不足\cite{poterbsummers1988,famafrench1988,richardson1993}。其二，$\theta$ 的估计对是否含时间趋势、是否含结构断点极其敏感，而金融价格序列两者都常见。其三，OU 假设线性回复力与常数波动率；若偏离越大、回复力反而越弱（非线性均值回归），OU 估计会高估回归速度。

\subsection{方差比统计量}

设 $\{r_t\}$ 平稳，$q$ 期累计收益的方差与单期方差之比定义为
\begin{equation}
VR(q)=\frac{\mathrm{Var}\left(\sum_{k=0}^{q-1}r_{t-k}\right)}{q\,\mathrm{Var}(r_t)}.
\end{equation}
把分子展开为 $q$ 个方差项与 $q(q-1)$ 个协方差项，整理后得到
\begin{equation}
VR(q)=1+2\sum_{k=1}^{q-1}\left(1-\frac{k}{q}\right)\rho_k ,
\end{equation}
其中 $\rho_k=\mathrm{Corr}(r_t,r_{t-k})$。这一步是正文“统一框架”的关键：它把“策略归属”转化为“加权自相关之和是否大于 1”。

三个立即的推论：

\begin{itemize}[leftmargin=2.2em]
  \item $VR(q)=1$ 对一切 $q$ 当且仅当所有 $\rho_k=0$（随机游走）。
  \item 若 $\rho_1<0$ 而 $\rho_k\approx0\ (k\ge2)$，则 $VR(q)<1$ 且随 $q$ 单调逼近 $1-2|\rho_1|$ 的极限——$VR(q)$ 随 $q$ 上升，对应短期反转。
  \item 若 $\rho_k>0$ 且衰减缓慢（长记忆），$VR(q)\gg1$ 且随 $q$ 持续增长，对应趋势。
\end{itemize}

\textbf{检验统计量}。在同方差假设下，Lo 与 MacKinlay(1988) 给出
\begin{equation}
z(q)=\frac{VR(q)-1}{\sqrt{\dfrac{2(2q-1)(q-1)}{3q\,n}}}\ \xrightarrow{d}\ \mathcal N(0,1),
\end{equation}
$n$ 为样本量。异方差稳健版本 $z^*(q)$ 用估计的 $\hat\delta_j$ 替换渐近方差（构造见原文）。对本文的启示有两条：其一，$z(q)$ 的分母随 $q$ 增大而增大，\emph{长周期检验的功效天然低于短周期}——这解释了为什么长期反转的证据总是比短期反转脆弱；其二，个股层面的 $|\rho_k|$ 被特质噪声稀释，因此 $z(q)$ 在个股上难以显著，而在组合/指数上容易显著\cite{lomackinlay1988,richardson1993}。

\subsection{最优交易带与成本的三次方根}

反转策略不能连续交易——每次穿越阈值都要付成本。考虑价差 $X_t$ 服从 OU 过程，目标是在线性交易成本下最大化长期收益。Martin 与 Sch\"oneborn(2011) 的结果是：最优策略是一个\emph{无交易带}（no-trade band），当 $|X_t-\mu|$ 落在带内时仓位不变；带宽 $b$ 满足
\begin{equation}
b\;\propto\;\left(\frac{\text{单位交易成本}}{\text{风险厌恶}\times \text{回归强度}}\right)^{1/3}.
\end{equation}
三次方根的来源可以直觉地理解：成本上升时，为避免频繁交易必须加宽带；但加宽带意味着等待更久、承担更多偏离风险，因此带宽随成本上升的弹性小于线性。这一结论有很强的实操含义——它把“多高的成本对应多宽的带宽”变成一个可计算的量，代价是假设线性成本与已知的 OU 动态。

在极端情形下，若加入止损退出（防线不可无限等待），问题变为双停时最优控制，Leung 与 Li(2015) 给出了解析的进出场价格区间\cite{leungli2015}；若允许幂效用，Boguslavsky 与 Boguslavskaya(2004) 给出阈值带的解析解\cite{boguslavsky2004}；若目标是“低买高卖”的直接阈值形式，Zhang 与 Zhang(2008) 给出了最优规则及其性质\cite{zhangzhang2008}。三者的共同结构是：\emph{反转策略的核心决策变量是带宽，而非方向}。

\subsection{GARCH 与波动率均值回归}

价格方向之外还有二阶矩的均值回归。Bollerslev(1986) 的 GARCH(1,1) 模型
\begin{equation}
h_t=\omega+\alpha\,\varepsilon_{t-1}^2+\beta\,h_{t-1},\qquad
\varepsilon_t=\sqrt{h_t}\,z_t,\ z_t\sim \mathcal N(0,1)
\end{equation}
的条件方差期望收敛于
\begin{equation}
\bar h=\frac{\omega}{1-\alpha-\beta}\qquad(\alpha+\beta<1),
\end{equation}
偏离以 $\alpha+\beta$ 的速率衰减。这就是“波动率均值回归”的计量载体。

为何把它单列？因为它与价格方向的均值回归在\emph{载体}上不同：$\theta>0$ 刻画的是“价格回到均值”，$h_t\to\bar h$ 刻画的是“波动幅度回到其长期水平”。二者可以同时成立、也可以互相矛盾（例如价格趋势强、而波动率均值回归正常）。第~\ref{sec:unity}~节的符号统一框架只覆盖前者；把它硬套到波动率上会得到错误推论。

与交易直接相关的是方差风险溢价：Carr 与 Wu(2009) 用指数期权定价出方差风险溢价并验证其与收益的关系\cite{carrwu2009}，Bollerslev 等(2009) 进一步发现它可显著预测短期股票收益\cite{bollerslev2009}。这两项工作构成了“做空波动率均值回归”这一策略族的理论依据，也解释了为什么波动率类策略与方向性策略的相关性通常很低。

\subsection{本文实证对称性的形式化陈述}

第~\ref{sec:empirical}~节的 G1/G2 门禁写成式子只需一行。设某 $L$ 下的目标权重为 $w_t$，趋势与反转的权重分别为 $+w_t$ 与 $-w_t$，则换手 $|w_t-w_{t-1}|$ 相同、毛收益 $w_t^\top r_{t+1}$ 与 $-w_t^\top r_{t+1}$ 互为相反数，而净收益两侧各扣除同一个 $c\cdot\text{换手}_t$（$c$ 为不区分方向的单边成本）。毛利夏普因此严格镜像为 $+S$ 与 $-S$，净夏普则约为 $(+S-\kappa)$ 与 $(-S-\kappa)$，$\kappa$ 为成本拖累对应的夏普损失——\textbf{成本把两侧差距从 $2S$ 放大到约 $2S+2\kappa$}：好的一侧被削去 $\kappa$，差的一侧再多压一个 $\kappa$。第~\ref{sec:empirical}~节已用实测数字验证这一结构（$L=5$ 时为 $+0.20$ 与 $-0.60$，而非 $\pm0.20$）。

这条式子可推广：任何“信号符号翻转”型的策略对偶，在扣除不区分方向的成本后，都会表现出“胜者优势被削弱、败者劣势被放大”的不对称。它与具体市场无关，是成本对策略比较的结构性影响。

\section{证据表（核心来源逐条核验）}\label{sec:evidence}

下表收录本文引用的核心来源。\textbf{核验状态}约定：\textbf{已核验}=标题、作者、年份、期刊卷期与 DOI 经 OpenAlex 或 Crossref 官方接口返回互证；\textbf{部分核验}=专著或无 DOI 期刊，经出版页/馆藏/权威引用交叉确认；\textbf{待复核}=仅能定位来源、具体书目细节需人工核对（正文避免依赖其数值）。全部链接指向原始出处。

\begin{longtable}{@{}>{\raggedright\arraybackslash}p{2.55cm}>{\raggedright\arraybackslash}p{4.7cm}>{\raggedright\arraybackslash}p{2.7cm}>{\raggedright\arraybackslash}p{1.9cm}>{\raggedright\arraybackslash}p{1.45cm}@{}}
\caption{核心来源清单（共 111 条）}\label{tab:evidence}\\ 
\toprule
\textbf{来源（作者，年）} & \textbf{标题} & \textbf{出处} & \textbf{链接} & \textbf{核验} \\
\midrule
\endfirsthead
\toprule \textbf{来源} & \textbf{标题} & \textbf{出处} & \textbf{链接} & \textbf{核验} \\ \midrule \endhead
\bottomrule \endlastfoot
Bachelier, L. (1900) & Théorie de la spéculation & Annales scientifiques de l'École Normale Supérieure 17: 21--86 & \href{https://doi.org/10.24033/asens.476}{链接} & 已核验 \\
Uhlenbeck, G. E. \& Ornstein, L. S. (1930) & On the Theory of the Brownian Motion & Physical Review 36(6): 823--841 & \href{https://doi.org/10.1103/PhysRev.36.823}{链接} & 已核验 \\
Cowles, A. (1933) & Can Stock Market Forecasters Forecast? & Econometrica 1(3): 309--324 & \href{https://doi.org/10.2307/1907042}{链接} & 已核验 \\
Cowles, A. \& Jones, H. E. (1937) & Some A Posteriori Probabilities in Stock Market Action & Econometrica 5(3): 280--294 & \href{https://doi.org/10.2307/1905515}{链接} & 已核验 \\
Kendall, M. G. (1953) & The Analysis of Economic Time-Series---Part I: Prices & J. Royal Statistical Society A 116(1): 11--34 & \href{https://doi.org/10.2307/2980947}{链接} & 已核验 \\
Fama, E. F. (1970) & Efficient Capital Markets: A Review of Theory and Empirical Work & Journal of Finance 25(2): 383--417 & \href{https://doi.org/10.1111/j.1540-6261.1970.tb00518.x}{链接} & 待复核 \\
Shiller, R. J. (1981) & Do Stock Prices Move Too Much to be Justified by Subsequent Changes in Dividends? & American Economic Review 71(3): 421--436 & \href{https://www.nber.org/papers/w0456}{链接} & 待复核 \\
Levy, R. A. (1967) & Relative Strength as a Criterion for Investment Selection & Journal of Finance 22(4): 595--610 & \href{https://doi.org/10.1111/j.1540-6261.1967.tb00295.x}{链接} & 已核验 \\
Vasicek, O. (1977) & An Equilibrium Characterization of the Term Structure & Journal of Financial Economics 5(2): 177--188 & \href{https://doi.org/10.1016/0304-405X(77)90016-2}{链接} & 已核验 \\
Dickey, D. A. \& Fuller, W. A. (1979) & Distribution of the Estimators for Autoregressive Time Series With a Unit Root & JASA 74(366): 427--431 & \href{https://doi.org/10.1080/01621459.1979.10482531}{链接} & 已核验 \\
De Bondt, W. F. M. \& Thaler, R. (1985) & Does the Stock Market Overreact? & Journal of Finance 40(3): 793--805 & \href{https://doi.org/10.1111/j.1540-6261.1985.tb05004.x}{链接} & 已核验 \\
Bollerslev, T. (1986) & Generalized Autoregressive Conditional Heteroskedasticity & Journal of Econometrics 31(3): 307--327 & \href{https://doi.org/10.1016/0304-4076(86)90063-1}{链接} & 已核验 \\
Engle, R. F. (1982) & Autoregressive Conditional Heteroscedasticity with Estimates of the Variance of United Kingdom Inflation & Econometrica 50(4): 987--1007 & \href{https://doi.org/10.2307/1912773}{链接} & 已核验 \\
Engle, R. F. \& Granger, C. W. J. (1987) & Co-Integration and Error Correction: Representation, Estimation, and Testing & Econometrica 55(2): 251--276 & \href{https://doi.org/10.2307/1913236}{链接} & 已核验 \\
Lo, A. W. \& MacKinlay, A. C. (1988) & Stock Market Prices Do Not Follow Random Walks: Evidence from a Simple Specification Test & Review of Financial Studies 1(1): 41--66 & \href{https://doi.org/10.1093/rfs/1.1.41}{链接} & 已核验 \\
Poterba, J. M. \& Summers, L. H. (1988) & Mean Reversion in Stock Prices: Evidence and Implications & Journal of Financial Economics 22(1): 27--59 & \href{https://doi.org/10.1016/0304-405X(88)90021-9}{链接} & 已核验 \\
Fama, E. F. \& French, K. R. (1988) & Permanent and Temporary Components of Stock Prices & Journal of Political Economy 96(2): 246--273 & \href{https://doi.org/10.1086/261535}{链接} & 已核验 \\
Perold, A. F. (1988) & The Implementation Shortfall: Paper versus Reality & Journal of Portfolio Management 14(3): 4--9 & \href{https://doi.org/10.3905/jpm.1988.409150}{链接} & 已核验 \\
Summers, L. H. (1986) & Does the Stock Market Rationally Reflect Fundamental Values? & Journal of Finance 41(3): 591--601 & \href{https://doi.org/10.1111/j.1540-6261.1986.tb04519.x}{链接} & 已核验 \\
Lehmann, B. N. (1990) & Fads, Martingales, and Market Efficiency & Quarterly Journal of Economics 105(1): 1--28 & \href{https://doi.org/10.2307/2937816}{链接} & 已核验 \\
Jegadeesh, N. (1990) & Evidence of Predictable Behavior of Security Returns & Journal of Finance 45(3): 881--898 & \href{https://doi.org/10.1111/j.1540-6261.1990.tb05110.x}{链接} & 已核验 \\
Lo, A. W. \& MacKinlay, A. C. (1990) & When Are Contrarian Profits Due to Stock Market Overreaction? & Review of Financial Studies 3(2): 175--205 & \href{https://doi.org/10.1093/rfs/3.2.175}{链接} & 已核验 \\
De Long, J. B., Shleifer, A., Summers, L. H. \& Waldmann, R. J. (1990) & Noise Trader Risk in Financial Markets & Journal of Political Economy 98(4): 703--738 & \href{https://doi.org/10.1086/261703}{链接} & 已核验 \\
Cutler, D. M., Poterba, J. M. \& Summers, L. H. (1991) & Speculative Dynamics & Review of Economic Studies 58(3): 529--546 & \href{https://doi.org/10.2307/2298010}{链接} & 已核验 \\
Brock, W., Lakonishok, J. \& LeBaron, B. (1992) & Simple Technical Trading Rules and the Stochastic Properties of Stock Returns & Journal of Finance 47(5): 1731--1764 & \href{https://doi.org/10.1111/j.1540-6261.1992.tb04681.x}{链接} & 已核验 \\
Jegadeesh, N. \& Titman, S. (1993) & Returns to Buying Winners and Selling Losers: Implications for Stock Market Efficiency & Journal of Finance 48(1): 65--91 & \href{https://doi.org/10.1111/j.1540-6261.1993.tb04702.x}{链接} & 已核验 \\
Richardson, M. (1993) & Temporary Components of Stock Prices: A Skeptic's View & Journal of Business \& Economic Statistics 11(2): 199--207 & \href{https://doi.org/10.1080/07350015.1993.10509948}{链接} & 已核验 \\
Lakonishok, J., Shleifer, A. \& Vishny, R. W. (1994) & Contrarian Investment, Extrapolation, and Risk & Journal of Finance 49(5): 1541--1578 & \href{https://doi.org/10.1111/j.1540-6261.1994.tb04772.x}{链接} & 已核验 \\
Chan, L. K. C., Jegadeesh, N. \& Lakonishok, J. (1996) & Momentum Strategies & Journal of Finance 51(5): 1681--1713 & \href{https://doi.org/10.1111/j.1540-6261.1996.tb05222.x}{链接} & 已核验 \\
Shleifer, A. \& Vishny, R. W. (1997) & The Limits of Arbitrage & Journal of Finance 52(1): 35--55 & \href{https://doi.org/10.1111/j.1540-6261.1997.tb03807.x}{链接} & 已核验 \\
Fung, W. \& Hsieh, D. A. (1997) & Empirical Characteristics of Dynamic Trading Strategies: The Case of Hedge Funds & Review of Financial Studies 10(2): 275--302 & \href{https://doi.org/10.1093/rfs/10.2.275}{链接} & 已核验 \\
Campbell, J. Y., Lo, A. W. \& MacKinlay, A. C. (1997) & The Econometrics of Financial Markets & Princeton University Press & \href{https://doi.org/10.1515/9781400830213}{链接} & 已核验 \\
Richards, A. J. (1997) & Winner-Loser Reversals in National Stock Market Indices: Can They be Explained? & Journal of Finance 52(5): 2129--2144 & \href{https://doi.org/10.2307/2329478}{链接} & 已核验 \\
Rouwenhorst, K. G. (1998) & International Momentum Strategies & Journal of Finance 53(1): 267--284 & \href{https://doi.org/10.1111/0022-1082.95722}{链接} & 已核验 \\
Barberis, N., Shleifer, A. \& Vishny, R. (1998) & A Model of Investor Sentiment & Journal of Financial Economics 49(3): 307--343 & \href{https://doi.org/10.1016/S0304-405X(98)00027-0}{链接} & 已核验 \\
Daniel, K., Hirshleifer, D. \& Subrahmanyam, A. (1998) & Investor Psychology and Security Market Under- and Overreactions & Journal of Finance 53(6): 1839--1885 & \href{https://doi.org/10.1111/0022-1082.00077}{链接} & 已核验 \\
Hong, H. \& Stein, J. C. (1999) & A Unified Theory of Underreaction, Momentum Trading, and Overreaction in Asset Markets & Journal of Finance 54(6): 2143--2184 & \href{https://doi.org/10.1111/0022-1082.00184}{链接} & 已核验 \\
Moskowitz, T. J. \& Grinblatt, M. (1999) & Do Industries Explain Momentum? & Journal of Finance 54(4): 1249--1290 & \href{https://doi.org/10.1111/0022-1082.00146}{链接} & 已核验 \\
Sullivan, R., Timmermann, A. \& White, H. (1999) & Data-Snooping, Technical Trading Rule Performance, and the Bootstrap & Journal of Finance 54(5): 1647--1691 & \href{https://doi.org/10.1111/0022-1082.00163}{链接} & 已核验 \\
Balvers, R., Wu, Y. \& Gilliland, E. (2000) & Mean Reversion across National Stock Markets and Parametric Contrarian Investment Strategies & Journal of Finance 55(2): 745--772 & \href{https://doi.org/10.1111/0022-1082.00225}{链接} & 已核验 \\
White, H. (2000) & A Reality Check for Data Snooping & Econometrica 68(5): 1097--1126 & \href{https://doi.org/10.1111/1468-0262.00152}{链接} & 已核验 \\
Almgren, R. \& Chriss, N. (2000) & Optimal Execution of Portfolio Transactions & The Journal of Risk 3(2): 5--39 & \href{https://doi.org/10.21314/jor.2001.041}{链接} & 已核验 \\
Harvey, C. R. \& Siddique, A. (2000) & Conditional Skewness in Asset Pricing Tests & Journal of Finance 55(3): 1263--1295 & \href{https://doi.org/10.1111/0022-1082.00247}{链接} & 已核验 \\
Jorion, P. (2000) & Risk Management Lessons from Long-Term Capital Management & European Financial Management 6(3): 277--300 & \href{https://doi.org/10.1111/1468-036X.00125}{链接} & 已核验 \\
Lowenstein, R. (2000) & When Genius Failed: The Rise and Fall of Long-Term Capital Management & Random House (ISBN 0-375-50317-X) & \href{https://www.randomhousebooks.com/}{链接} & 部分核验 \\
Fung, W. \& Hsieh, D. A. (2001) & The Risk in Hedge Fund Strategies: Theory and Evidence from Trend Followers & Review of Financial Studies 14(2): 313--341 & \href{https://doi.org/10.1093/rfs/14.2.313}{链接} & 已核验 \\
Jegadeesh, N. \& Titman, S. (2001) & Profitability of Momentum Strategies: An Evaluation of Alternative Explanations & Journal of Finance 56(2): 699--720 & \href{https://doi.org/10.1111/0022-1082.00342}{链接} & 已核验 \\
Lo, A. W. (2002) & The Statistics of Sharpe Ratios & Financial Analysts Journal 58(4): 36--52 & \href{https://doi.org/10.2469/faj.v58.n4.2453}{链接} & 已核验 \\
Ang, A. \& Bekaert, G. (2002) & International Asset Allocation With Regime Shifts & Review of Financial Studies 15(4): 1137--1187 & \href{https://doi.org/10.1093/rfs/15.4.1137}{链接} & 已核验 \\
Pástor, Ľ. \& Stambaugh, R. F. (2003) & Liquidity Risk and Expected Stock Returns & Journal of Political Economy 111(3): 642--685 & \href{https://doi.org/10.1086/374184}{链接} & 已核验 \\
Vidyamurthy, G. (2004) & Pairs Trading: Quantitative Methods and Analysis & John Wiley \& Sons (ISBN 978-0-471-46067-1) & \href{https://www.wiley.com/}{链接} & 部分核验 \\
Boguslavsky, M. \& Boguslavskaya, E. (2004) & Arbitrage under Power & Risk 17(6): 69--73（无 DOI） & \href{https://www.risk.net/}{链接} & 部分核验 \\
Korajczyk, R. A. \& Sadka, R. (2004) & Are Momentum Profits Robust to Trading Costs? & Journal of Finance 59(3): 1039--1082 & \href{https://doi.org/10.1111/j.1540-6261.2004.00656.x}{链接} & 已核验 \\
Cochrane, J. H. (2005) & Financial Markets and the Real Economy & NBER WP 11193 & \href{https://doi.org/10.3386/w11193}{链接} & 已核验 \\
Hansen, P. R. (2005) & A Test for Superior Predictive Ability & Journal of Business \& Economic Statistics 23(4): 365--380 & \href{https://doi.org/10.1198/073500105000000063}{链接} & 已核验 \\
Romano, J. P. \& Wolf, M. (2005) & Stepwise Multiple Testing as Formalized Data Snooping & Econometrica 73(4): 1237--1282 & \href{https://doi.org/10.1111/j.1468-0262.2005.00615.x}{链接} & 已核验 \\
Elliott, R. J., van der Hoek, J. \& Malcolm, W. P. (2005) & Pairs Trading & Quantitative Finance 5(3): 271--276 & \href{https://doi.org/10.1080/14697680500149370}{链接} & 已核验 \\
Gatev, E., Goetzmann, W. N. \& Rouwenhorst, K. G. (2006) & Pairs Trading: Performance of a Relative-Value Arbitrage Rule & Review of Financial Studies 19(3): 797--827 & \href{https://doi.org/10.1093/rfs/hhj020}{链接} & 已核验 \\
Erb, C. B. \& Harvey, C. R. (2006) & The Strategic and Tactical Value of Commodity Futures & Financial Analysts Journal 62(2): 69--97 & \href{https://doi.org/10.2469/faj.v62.n2.4084}{链接} & 已核验 \\
Zhang, H. \& Zhang, Q. (2008) & Trading a Mean-Reverting Asset: Buy Low and Sell High & Automatica 44(6): 1511--1518 & \href{https://doi.org/10.1016/j.automatica.2007.11.003}{链接} & 已核验 \\
Stein, J. C. (2009) & Presidential Address: Sophisticated Investors and Market Efficiency & Journal of Finance 64(4): 1517--1548 & \href{https://doi.org/10.1111/j.1540-6261.2009.01472.x}{链接} & 已核验 \\
Brunnermeier, M. K. \& Pedersen, L. H. (2009) & Market Liquidity and Funding Liquidity & Review of Financial Studies 22(6): 2201--2238 & \href{https://doi.org/10.1093/rfs/hhn098}{链接} & 已核验 \\
Carr, P. \& Wu, L. (2009) & Variance Risk Premiums & Review of Financial Studies 22(3): 1311--1341 & \href{https://doi.org/10.1093/rfs/hhn038}{链接} & 已核验 \\
Bollerslev, T., Tauchen, G. \& Zhou, H. (2009) & Expected Stock Returns and Variance Risk Premia & Review of Financial Studies 22(11): 4463--4492 & \href{https://doi.org/10.1093/rfs/hhp008}{链接} & 已核验 \\
Avellaneda, M. \& Lee, J.-H. (2010) & Statistical Arbitrage in the US Equities Market & Quantitative Finance 10(7): 761--782 & \href{https://doi.org/10.1080/14697680903124632}{链接} & 已核验 \\
Gromb, D. \& Vayanos, D. (2010) & Limits of Arbitrage & Annual Review of Financial Economics 2: 251--275 & \href{https://doi.org/10.1146/annurev-financial-073009-104107}{链接} & 已核验 \\
Martin, R. \& Schöneborn, T. (2011) & Mean Reversion Pays, but Costs & Risk 24(2): 84--89；arXiv:\allowbreak 1103.4934 & \href{https://arxiv.org/abs/1103.4934}{链接} & 部分核验 \\
Moskowitz, T. J., Ooi, Y. H. \& Pedersen, L. H. (2012) & Time Series Momentum & Journal of Financial Economics 104(2): 228--250 & \href{https://doi.org/10.1016/j.jfineco.2011.11.003}{链接} & 已核验 \\
Asness, C. S., Moskowitz, T. J. \& Pedersen, L. H. (2013) & Value and Momentum Everywhere & Journal of Finance 68(3): 929--985 & \href{https://doi.org/10.1111/jofi.12021}{链接} & 已核验 \\
Do, B. \& Faff, R. (2012) & Are Pairs Trading Profits Robust to Trading Costs? & Journal of Financial Research 35(2): 261--287 & \href{https://doi.org/10.1111/j.1475-6803.2012.01317.x}{链接} & 已核验 \\
Khandani, A. E. \& Lo, A. W. (2011) & What Happened to the Quants in August 2007? Evidence from Factors and Transactions Data & Journal of Financial Markets 14(1): 1--46 & \href{https://doi.org/10.1016/j.finmar.2010.07.005}{链接} & 已核验 \\
Menkhoff, L., Sarno, L., Schmeling, M. \& Schrimpf, A. (2012) & Currency Momentum Strategies & Journal of Financial Economics 106(3): 660--684 & \href{https://doi.org/10.1016/j.jfineco.2012.06.009}{链接} & 已核验 \\
Frazzini, A., Israel, R. \& Moskowitz, T. J. (2018) & Trading Costs of Asset Pricing Anomalies & SSRN / 工作论文 & \href{https://doi.org/10.2139/ssrn.2294498}{链接} & 已核验 \\
Baltas, N. \& Kosowski, R. (2011) & Momentum Strategies in Futures Markets and Trend-Following Funds & SSRN 工作论文 & \href{https://doi.org/10.2139/ssrn.1968996}{链接} & 已核验 \\
Baltas, N. \& Kosowski, R. (2013) & Improving Time-Series Momentum Strategies: The Role of Trading Signals and Volatility Estimators & SSRN 工作论文 & \href{https://doi.org/10.2139/ssrn.2140091}{链接} & 已核验 \\
Cuturi, M. \& d'Aspremont, A. (2013) & Mean Reversion with a Variance Threshold & ICML / PMLR 28(3): 271--279 & \href{https://proceedings.mlr.press/v28/cuturi13.html}{链接} & 已核验 \\
Ilmanen, A. (2011) & Expected Returns: An Investor's Guide to Harvesting Market Rewards & John Wiley \& Sons & \href{https://doi.org/10.1002/9781118467190}{链接} & 已核验 \\
Ang, A. \& Timmermann, A. (2012) & Regime Changes and Financial Markets & Annual Review of Financial Economics 4: 313--337 & \href{https://doi.org/10.1146/annurev-financial-110311-101808}{链接} & 已核验 \\
Hamilton, J. D. (1989) & A New Approach to the Economic Analysis of Nonstationary Time Series and the Business Cycle & Econometrica 57(2): 357--384 & \href{https://doi.org/10.2307/1912559}{链接} & 已核验 \\
Gârleanu, N. \& Pedersen, L. H. (2013) & Dynamic Trading with Predictable Returns and Transaction Costs & Journal of Finance 68(6): 2309--2340 & \href{https://doi.org/10.1111/jofi.12080}{链接} & 已核验 \\
Krauss, C. (2017) & Statistical Arbitrage Pairs Trading Strategies: Review and Outlook & Journal of Economic Surveys 31(2): 513--545 & \href{https://doi.org/10.1111/joes.12153}{链接} & 已核验 \\
Barroso, P. \& Santa-Clara, P. (2015) & Momentum Has Its Moments & Journal of Financial Economics 116(1): 111--120 & \href{https://doi.org/10.1016/j.jfineco.2014.11.010}{链接} & 已核验 \\
Leung, T. \& Li, X. (2015) & Optimal Mean Reversion Trading with Transaction Costs and Stop-Loss Exit & International Journal of Theoretical and Applied Finance 18(3): 1550020 & \href{https://doi.org/10.1142/S021902491550020X}{链接} & 已核验 \\
Harvey, C. R. \& Liu, Y. (2015) & Backtesting & Journal of Portfolio Management 42(1): 13--28 & \href{https://doi.org/10.3905/jpm.2015.42.1.013}{链接} & 已核验 \\
Daniel, K. \& Moskowitz, T. J. (2016) & Momentum Crashes & Journal of Financial Economics 122(2): 221--247 & \href{https://doi.org/10.1016/j.jfineco.2015.12.002}{链接} & 已核验 \\
Harvey, C. R., Liu, Y. \& Zhu, H. (2016) & …and the Cross-Section of Expected Returns & Review of Financial Studies 29(1): 5--68 & \href{https://doi.org/10.1093/rfs/hhv059}{链接} & 已核验 \\
McLean, R. D. \& Pontiff, J. (2016) & Does Academic Research Destroy Stock Return Predictability? & Journal of Finance 71(1): 5--32 & \href{https://doi.org/10.1111/jofi.12365}{链接} & 已核验 \\
Kim, A. Y., Tse, Y. \& Wald, J. K. (2016) & Time Series Momentum and Volatility Scaling & Journal of Financial Markets 30: 103--124 & \href{https://doi.org/10.1016/j.finmar.2016.05.003}{链接} & 已核验 \\
Bailey, D. H., Borwein, J. M., López de Prado, M. \& Zhu, Q. J. (2016) & The Probability of Backtest Overfitting & The Journal of Computational Finance 20(4): 39--69 & \href{https://doi.org/10.21314/jcf.2016.322}{链接} & 已核验 \\
Novy-Marx, R. \& Velikov, M. (2016) & A Taxonomy of Anomalies and Their Trading Costs & Review of Financial Studies 29(1): 104--147 & \href{https://doi.org/10.1093/rfs/hhv063}{链接} & 已核验 \\
Hurst, B., Ooi, Y. H. \& Pedersen, L. H. (2017) & A Century of Evidence on Trend-Following Investing & Journal of Portfolio Management 44(1): 15--29 & \href{https://doi.org/10.3905/jpm.2017.44.1.015}{链接} & 已核验 \\
Georgopoulou, A. \& Wang, J. (2017) & The Trend Is Your Friend: Time-Series Momentum Strategies across Equity and Commodity Markets & Review of Finance 21(4): 1557--1592 & \href{https://doi.org/10.1093/rof/rfw048}{链接} & 已核验 \\
Harvey, C. R., Hoyle, E., Korgaonkar, R., Rattray, S., Sargaison, M. \& van Hemert, O. (2018) & The Impact of Volatility Targeting & Journal of Portfolio Management 45(1): 14--33 & \href{https://doi.org/10.3905/jpm.2018.45.1.014}{链接} & 已核验 \\
Moreira, A. \& Muir, T. (2017) & Volatility-Managed Portfolios & Journal of Finance 72(4): 1611--1644 & \href{https://doi.org/10.1111/jofi.12513}{链接} & 已核验 \\
Bailey, D. H. \& López de Prado, M. (2014) & The Deflated Sharpe Ratio & Journal of Portfolio Management 40(5): 94--107 & \href{https://doi.org/10.3905/jpm.2014.40.5.094}{链接} & 已核验 \\
Huang, D., Li, J., Wang, L. \& Zhou, G. (2020) & Time Series Momentum: Is It There? & Journal of Financial Economics 135(3): 774--794 & \href{https://doi.org/10.1016/j.jfineco.2019.08.004}{链接} & 已核验 \\
Gu, S., Kelly, B. \& Xiu, D. (2020) & Empirical Asset Pricing via Machine Learning & Review of Financial Studies 33(5): 2223--2283 & \href{https://doi.org/10.1093/rfs/hhaa009}{链接} & 已核验 \\
Baltas, N. (2019) & The Impact of Crowding in Alternative Risk Premia Investing & Financial Analysts Journal 75(3): 89--104 & \href{https://doi.org/10.1080/0015198x.2019.1600955}{链接} & 已核验 \\
Baltas, N. \& Kosowski, R. (2020) & Demystifying Time-Series Momentum Strategies & Wiley book chapter & \href{https://doi.org/10.1002/9781119599364.ch3}{链接} & 已核验 \\
Cochrane, J. H. (2011) & Presidential Address: Discount Rates & Journal of Finance 66(4): 1047--1108 & \href{https://doi.org/10.1111/j.1540-6261.2011.01671.x}{链接} & 已核验 \\
Wood, K., Roberts, S. \& Zohren, S. (2021) & Slow Momentum with Fast Reversion: A Trading Strategy Using Deep Learning and Changepoint Detection & Journal of Financial Data Science；arXiv:\allowbreak 2105.13727 & \href{https://arxiv.org/abs/2105.13727}{链接} & 已核验 \\
Ehsani, S. \& Linnainmaa, J. T. (2022) & Factor Momentum and the Momentum Factor & Journal of Finance 77(3): 1877--1919 & \href{https://doi.org/10.1111/jofi.13131}{链接} & 已核验 \\
Lou, D. \& Polk, C. (2022) & Comomentum: Inferring Arbitrage Activity from Return Correlations & Review of Financial Studies 35(2): 735--776 & \href{https://doi.org/10.1093/rfs/hhab117}{链接} & 已核验 \\
Pitkäjärvi, A., Suominen, M. \& Vaittinen, J. (2020) & Cross-Asset Signals and Time Series Momentum & Journal of Financial Economics 136(2): 459--478 & \href{https://doi.org/10.1016/j.jfineco.2019.02.011}{链接} & 已核验 \\
Guijarro-Ordonez, J., Pelger, M. \& Zanotti, G. (2021) & Deep Learning Statistical Arbitrage & arXiv:\allowbreak 2106.04028 & \href{https://arxiv.org/abs/2106.04028}{链接} & 已核验 \\
Ning, B. \& Lee, A. (2024) & Advanced Statistical Arbitrage with Reinforcement Learning & arXiv:\allowbreak 2403.12180 & \href{https://arxiv.org/abs/2403.12180}{链接} & 已核验 \\
Suchato, P., Wiryadi, A., Chen, Y., Zhao, A. \& Yue, X. (2024) & An Application of the Ornstein--Uhlenbeck Process to Pairs Trading & arXiv:\allowbreak 2412.12458 & \href{https://arxiv.org/abs/2412.12458}{链接} & 已核验 \\
Fama, E. F. \& Blume, M. E. (1966) & Filter Rules and Stock-Market Trading & Journal of Business 39(1): 226--241 & \href{https://doi.org/10.1086/294849}{链接} & 已核验 \\
Lo, A. W., Mamaysky, H. \& Wang, J. (2000) & Foundations of Technical Analysis: Computational Algorithms, Statistical Inference, and Empirical Implementation & Journal of Finance 55(4): 1705--1765 & \href{https://doi.org/10.1111/0022-1082.00265}{链接} & 已核验 \\
Novy-Marx, R. (2012) & Is Momentum Really Momentum? & Journal of Financial Economics 104(2): 363--390 & \href{https://doi.org/10.1016/j.jfineco.2011.05.003}{链接} & 已核验 \\
Chabot, B., Ghysels, E. \& Jagannathan, R. (2014) & Momentum Trading, Return Chasing, and Predictable Crashes & NBER WP 20660 & \href{https://doi.org/10.3386/w20660}{链接} & 已核验 \\
\end{longtable}

\section*{参考文献}
\addcontentsline{toc}{section}{参考文献}
\begin{thebibliography}{99}
\small
\bibitem{bachelier1900} Bachelier, L. (1900). \emph{Théorie de la spéculation}. Annales scientifiques de l'École Normale Supérieure 17: 21--86. \href{https://doi.org/10.24033/asens.476}{https://doi.org/10.24033/asens.476}.
\bibitem{uo1930} Uhlenbeck, G. E. \& Ornstein, L. S. (1930). \emph{On the Theory of the Brownian Motion}. Physical Review 36(6): 823--841. \href{https://doi.org/10.1103/PhysRev.36.823}{https://doi.org/10.1103/PhysRev.36.823}.
\bibitem{cowles1933} Cowles, A. (1933). \emph{Can Stock Market Forecasters Forecast?}. Econometrica 1(3): 309--324. \href{https://doi.org/10.2307/1907042}{https://doi.org/10.2307/1907042}.
\bibitem{cowlesjones1937} Cowles, A. \& Jones, H. E. (1937). \emph{Some A Posteriori Probabilities in Stock Market Action}. Econometrica 5(3): 280--294. \href{https://doi.org/10.2307/1905515}{https://doi.org/10.2307/1905515}.
\bibitem{kendall1953} Kendall, M. G. (1953). \emph{The Analysis of Economic Time-Series---Part I: Prices}. J. Royal Statistical Society A 116(1): 11--34. \href{https://doi.org/10.2307/2980947}{https://doi.org/10.2307/2980947}.
\bibitem{fama1970} Fama, E. F. (1970). \emph{Efficient Capital Markets: A Review of Theory and Empirical Work}. Journal of Finance 25(2): 383--417. \href{https://doi.org/10.1111/j.1540-6261.1970.tb00518.x}{https://doi.org/10.1111/j.1540-6261.1970.tb00518.x}.
\bibitem{shiller1981} Shiller, R. J. (1981). \emph{Do Stock Prices Move Too Much to be Justified by Subsequent Changes in Dividends?}. American Economic Review 71(3): 421--436. \href{https://www.nber.org/papers/w0456}{https://www.nber.org/papers/w0456}.
\bibitem{levy1967} Levy, R. A. (1967). \emph{Relative Strength as a Criterion for Investment Selection}. Journal of Finance 22(4): 595--610. \href{https://doi.org/10.1111/j.1540-6261.1967.tb00295.x}{https://doi.org/10.1111/j.1540-6261.1967.tb00295.x}.
\bibitem{vasicek1977} Vasicek, O. (1977). \emph{An Equilibrium Characterization of the Term Structure}. Journal of Financial Economics 5(2): 177--188. \href{https://doi.org/10.1016/0304-405X(77)90016-2}{https://doi.org/10.1016/0304-405X(77)90016-2}.
\bibitem{dickeyfuller1979} Dickey, D. A. \& Fuller, W. A. (1979). \emph{Distribution of the Estimators for Autoregressive Time Series With a Unit Root}. JASA 74(366): 427--431. \href{https://doi.org/10.1080/01621459.1979.10482531}{https://doi.org/10.1080/01621459.1979.10482531}.
\bibitem{debondt1985} De Bondt, W. F. M. \& Thaler, R. (1985). \emph{Does the Stock Market Overreact?}. Journal of Finance 40(3): 793--805. \href{https://doi.org/10.1111/j.1540-6261.1985.tb05004.x}{https://doi.org/10.1111/j.1540-6261.1985.tb05004.x}.
\bibitem{bollerslev1986} Bollerslev, T. (1986). \emph{Generalized Autoregressive Conditional Heteroskedasticity}. Journal of Econometrics 31(3): 307--327. \href{https://doi.org/10.1016/0304-4076(86)90063-1}{https://doi.org/10.1016/0304-4076(86)90063-1}.
\bibitem{engle1982} Engle, R. F. (1982). \emph{Autoregressive Conditional Heteroscedasticity with Estimates of the Variance of United Kingdom Inflation}. Econometrica 50(4): 987--1007. \href{https://doi.org/10.2307/1912773}{https://doi.org/10.2307/1912773}.
\bibitem{englegranger1987} Engle, R. F. \& Granger, C. W. J. (1987). \emph{Co-Integration and Error Correction: Representation, Estimation, and Testing}. Econometrica 55(2): 251--276. \href{https://doi.org/10.2307/1913236}{https://doi.org/10.2307/1913236}.
\bibitem{lomackinlay1988} Lo, A. W. \& MacKinlay, A. C. (1988). \emph{Stock Market Prices Do Not Follow Random Walks: Evidence from a Simple Specification Test}. Review of Financial Studies 1(1): 41--66. \href{https://doi.org/10.1093/rfs/1.1.41}{https://doi.org/10.1093/rfs/1.1.41}.
\bibitem{poterbsummers1988} Poterba, J. M. \& Summers, L. H. (1988). \emph{Mean Reversion in Stock Prices: Evidence and Implications}. Journal of Financial Economics 22(1): 27--59. \href{https://doi.org/10.1016/0304-405X(88)90021-9}{https://doi.org/10.1016/0304-405X(88)90021-9}.
\bibitem{famafrench1988} Fama, E. F. \& French, K. R. (1988). \emph{Permanent and Temporary Components of Stock Prices}. Journal of Political Economy 96(2): 246--273. \href{https://doi.org/10.1086/261535}{https://doi.org/10.1086/261535}.
\bibitem{perold1988} Perold, A. F. (1988). \emph{The Implementation Shortfall: Paper versus Reality}. Journal of Portfolio Management 14(3): 4--9. \href{https://doi.org/10.3905/jpm.1988.409150}{https://doi.org/10.3905/jpm.1988.409150}.
\bibitem{summers1986} Summers, L. H. (1986). \emph{Does the Stock Market Rationally Reflect Fundamental Values?}. Journal of Finance 41(3): 591--601. \href{https://doi.org/10.1111/j.1540-6261.1986.tb04519.x}{https://doi.org/10.1111/j.1540-6261.1986.tb04519.x}.
\bibitem{lehmann1990} Lehmann, B. N. (1990). \emph{Fads, Martingales, and Market Efficiency}. Quarterly Journal of Economics 105(1): 1--28. \href{https://doi.org/10.2307/2937816}{https://doi.org/10.2307/2937816}.
\bibitem{jegadeesh1990} Jegadeesh, N. (1990). \emph{Evidence of Predictable Behavior of Security Returns}. Journal of Finance 45(3): 881--898. \href{https://doi.org/10.1111/j.1540-6261.1990.tb05110.x}{https://doi.org/10.1111/j.1540-6261.1990.tb05110.x}.
\bibitem{lomackinlay1990} Lo, A. W. \& MacKinlay, A. C. (1990). \emph{When Are Contrarian Profits Due to Stock Market Overreaction?}. Review of Financial Studies 3(2): 175--205. \href{https://doi.org/10.1093/rfs/3.2.175}{https://doi.org/10.1093/rfs/3.2.175}.
\bibitem{delong1990} De Long, J. B., Shleifer, A., Summers, L. H. \& Waldmann, R. J. (1990). \emph{Noise Trader Risk in Financial Markets}. Journal of Political Economy 98(4): 703--738. \href{https://doi.org/10.1086/261703}{https://doi.org/10.1086/261703}.
\bibitem{cutler1991} Cutler, D. M., Poterba, J. M. \& Summers, L. H. (1991). \emph{Speculative Dynamics}. Review of Economic Studies 58(3): 529--546. \href{https://doi.org/10.2307/2298010}{https://doi.org/10.2307/2298010}.
\bibitem{brock1992} Brock, W., Lakonishok, J. \& LeBaron, B. (1992). \emph{Simple Technical Trading Rules and the Stochastic Properties of Stock Returns}. Journal of Finance 47(5): 1731--1764. \href{https://doi.org/10.1111/j.1540-6261.1992.tb04681.x}{https://doi.org/10.1111/j.1540-6261.1992.tb04681.x}.
\bibitem{jegadeesh1993} Jegadeesh, N. \& Titman, S. (1993). \emph{Returns to Buying Winners and Selling Losers: Implications for Stock Market Efficiency}. Journal of Finance 48(1): 65--91. \href{https://doi.org/10.1111/j.1540-6261.1993.tb04702.x}{https://doi.org/10.1111/j.1540-6261.1993.tb04702.x}.
\bibitem{richardson1993} Richardson, M. (1993). \emph{Temporary Components of Stock Prices: A Skeptic's View}. Journal of Business \& Economic Statistics 11(2): 199--207. \href{https://doi.org/10.1080/07350015.1993.10509948}{https://doi.org/10.1080/07350015.1993.10509948}.
\bibitem{lakonishok1994} Lakonishok, J., Shleifer, A. \& Vishny, R. W. (1994). \emph{Contrarian Investment, Extrapolation, and Risk}. Journal of Finance 49(5): 1541--1578. \href{https://doi.org/10.1111/j.1540-6261.1994.tb04772.x}{https://doi.org/10.1111/j.1540-6261.1994.tb04772.x}.
\bibitem{chan1996} Chan, L. K. C., Jegadeesh, N. \& Lakonishok, J. (1996). \emph{Momentum Strategies}. Journal of Finance 51(5): 1681--1713. \href{https://doi.org/10.1111/j.1540-6261.1996.tb05222.x}{https://doi.org/10.1111/j.1540-6261.1996.tb05222.x}.
\bibitem{shleifervishny1997} Shleifer, A. \& Vishny, R. W. (1997). \emph{The Limits of Arbitrage}. Journal of Finance 52(1): 35--55. \href{https://doi.org/10.1111/j.1540-6261.1997.tb03807.x}{https://doi.org/10.1111/j.1540-6261.1997.tb03807.x}.
\bibitem{funghsieh1997} Fung, W. \& Hsieh, D. A. (1997). \emph{Empirical Characteristics of Dynamic Trading Strategies: The Case of Hedge Funds}. Review of Financial Studies 10(2): 275--302. \href{https://doi.org/10.1093/rfs/10.2.275}{https://doi.org/10.1093/rfs/10.2.275}.
\bibitem{campbell1997} Campbell, J. Y., Lo, A. W. \& MacKinlay, A. C. (1997). \emph{The Econometrics of Financial Markets}. Princeton University Press. \href{https://doi.org/10.1515/9781400830213}{https://doi.org/10.1515/9781400830213}.
\bibitem{richards1997} Richards, A. J. (1997). \emph{Winner-Loser Reversals in National Stock Market Indices: Can They be Explained?}. Journal of Finance 52(5): 2129--2144. \href{https://doi.org/10.2307/2329478}{https://doi.org/10.2307/2329478}.
\bibitem{rouwenhorst1998} Rouwenhorst, K. G. (1998). \emph{International Momentum Strategies}. Journal of Finance 53(1): 267--284. \href{https://doi.org/10.1111/0022-1082.95722}{https://doi.org/10.1111/0022-1082.95722}.
\bibitem{bsv1998} Barberis, N., Shleifer, A. \& Vishny, R. (1998). \emph{A Model of Investor Sentiment}. Journal of Financial Economics 49(3): 307--343. \href{https://doi.org/10.1016/S0304-405X(98)00027-0}{https://doi.org/10.1016/S0304-405X(98)00027-0}.
\bibitem{dhs1998} Daniel, K., Hirshleifer, D. \& Subrahmanyam, A. (1998). \emph{Investor Psychology and Security Market Under- and Overreactions}. Journal of Finance 53(6): 1839--1885. \href{https://doi.org/10.1111/0022-1082.00077}{https://doi.org/10.1111/0022-1082.00077}.
\bibitem{hongstein1999} Hong, H. \& Stein, J. C. (1999). \emph{A Unified Theory of Underreaction, Momentum Trading, and Overreaction in Asset Markets}. Journal of Finance 54(6): 2143--2184. \href{https://doi.org/10.1111/0022-1082.00184}{https://doi.org/10.1111/0022-1082.00184}.
\bibitem{moskowitzgrinblatt1999} Moskowitz, T. J. \& Grinblatt, M. (1999). \emph{Do Industries Explain Momentum?}. Journal of Finance 54(4): 1249--1290. \href{https://doi.org/10.1111/0022-1082.00146}{https://doi.org/10.1111/0022-1082.00146}.
\bibitem{sullivan1999} Sullivan, R., Timmermann, A. \& White, H. (1999). \emph{Data-Snooping, Technical Trading Rule Performance, and the Bootstrap}. Journal of Finance 54(5): 1647--1691. \href{https://doi.org/10.1111/0022-1082.00163}{https://doi.org/10.1111/0022-1082.00163}.
\bibitem{balvers2000} Balvers, R., Wu, Y. \& Gilliland, E. (2000). \emph{Mean Reversion across National Stock Markets and Parametric Contrarian Investment Strategies}. Journal of Finance 55(2): 745--772. \href{https://doi.org/10.1111/0022-1082.00225}{https://doi.org/10.1111/0022-1082.00225}.
\bibitem{white2000} White, H. (2000). \emph{A Reality Check for Data Snooping}. Econometrica 68(5): 1097--1126. \href{https://doi.org/10.1111/1468-0262.00152}{https://doi.org/10.1111/1468-0262.00152}.
\bibitem{almgrenchriss2000} Almgren, R. \& Chriss, N. (2000). \emph{Optimal Execution of Portfolio Transactions}. The Journal of Risk 3(2): 5--39. \href{https://doi.org/10.21314/jor.2001.041}{https://doi.org/10.21314/jor.2001.041}.
\bibitem{harveysiddique2000} Harvey, C. R. \& Siddique, A. (2000). \emph{Conditional Skewness in Asset Pricing Tests}. Journal of Finance 55(3): 1263--1295. \href{https://doi.org/10.1111/0022-1082.00247}{https://doi.org/10.1111/0022-1082.00247}.
\bibitem{jorion2000} Jorion, P. (2000). \emph{Risk Management Lessons from Long-Term Capital Management}. European Financial Management 6(3): 277--300. \href{https://doi.org/10.1111/1468-036X.00125}{https://doi.org/10.1111/1468-036X.00125}.
\bibitem{lowenstein2000} Lowenstein, R. (2000). \emph{When Genius Failed: The Rise and Fall of Long-Term Capital Management}. Random House (ISBN 0-375-50317-X). \href{https://www.randomhousebooks.com/}{https://www.randomhousebooks.com/}.
\bibitem{funghsieh2001} Fung, W. \& Hsieh, D. A. (2001). \emph{The Risk in Hedge Fund Strategies: Theory and Evidence from Trend Followers}. Review of Financial Studies 14(2): 313--341. \href{https://doi.org/10.1093/rfs/14.2.313}{https://doi.org/10.1093/rfs/14.2.313}.
\bibitem{jegadeesh2001} Jegadeesh, N. \& Titman, S. (2001). \emph{Profitability of Momentum Strategies: An Evaluation of Alternative Explanations}. Journal of Finance 56(2): 699--720. \href{https://doi.org/10.1111/0022-1082.00342}{https://doi.org/10.1111/0022-1082.00342}.
\bibitem{lo2002} Lo, A. W. (2002). \emph{The Statistics of Sharpe Ratios}. Financial Analysts Journal 58(4): 36--52. \href{https://doi.org/10.2469/faj.v58.n4.2453}{https://doi.org/10.2469/faj.v58.n4.2453}.
\bibitem{angbekaert2002} Ang, A. \& Bekaert, G. (2002). \emph{International Asset Allocation With Regime Shifts}. Review of Financial Studies 15(4): 1137--1187. \href{https://doi.org/10.1093/rfs/15.4.1137}{https://doi.org/10.1093/rfs/15.4.1137}.
\bibitem{pastorstambaugh2003} Pástor, Ľ. \& Stambaugh, R. F. (2003). \emph{Liquidity Risk and Expected Stock Returns}. Journal of Political Economy 111(3): 642--685. \href{https://doi.org/10.1086/374184}{https://doi.org/10.1086/374184}.
\bibitem{vidyamurthy2004} Vidyamurthy, G. (2004). \emph{Pairs Trading: Quantitative Methods and Analysis}. John Wiley \& Sons (ISBN 978-0-471-46067-1). \href{https://www.wiley.com/}{https://www.wiley.com/}.
\bibitem{boguslavsky2004} Boguslavsky, M. \& Boguslavskaya, E. (2004). \emph{Arbitrage under Power}. Risk 17(6): 69--73（无 DOI）. \href{https://www.risk.net/}{https://www.risk.net/}.
\bibitem{korajczyksadka2004} Korajczyk, R. A. \& Sadka, R. (2004). \emph{Are Momentum Profits Robust to Trading Costs?}. Journal of Finance 59(3): 1039--1082. \href{https://doi.org/10.1111/j.1540-6261.2004.00656.x}{https://doi.org/10.1111/j.1540-6261.2004.00656.x}.
\bibitem{cochrane2005} Cochrane, J. H. (2005). \emph{Financial Markets and the Real Economy}. NBER WP 11193. \href{https://doi.org/10.3386/w11193}{https://doi.org/10.3386/w11193}.
\bibitem{hansen2005} Hansen, P. R. (2005). \emph{A Test for Superior Predictive Ability}. Journal of Business \& Economic Statistics 23(4): 365--380. \href{https://doi.org/10.1198/073500105000000063}{https://doi.org/10.1198/073500105000000063}.
\bibitem{romanowolf2005} Romano, J. P. \& Wolf, M. (2005). \emph{Stepwise Multiple Testing as Formalized Data Snooping}. Econometrica 73(4): 1237--1282. \href{https://doi.org/10.1111/j.1468-0262.2005.00615.x}{https://doi.org/10.1111/j.1468-0262.2005.00615.x}.
\bibitem{elliott2005} Elliott, R. J., van der Hoek, J. \& Malcolm, W. P. (2005). \emph{Pairs Trading}. Quantitative Finance 5(3): 271--276. \href{https://doi.org/10.1080/14697680500149370}{https://doi.org/10.1080/14697680500149370}.
\bibitem{gatev2006} Gatev, E., Goetzmann, W. N. \& Rouwenhorst, K. G. (2006). \emph{Pairs Trading: Performance of a Relative-Value Arbitrage Rule}. Review of Financial Studies 19(3): 797--827. \href{https://doi.org/10.1093/rfs/hhj020}{https://doi.org/10.1093/rfs/hhj020}.
\bibitem{erbharvey2006} Erb, C. B. \& Harvey, C. R. (2006). \emph{The Strategic and Tactical Value of Commodity Futures}. Financial Analysts Journal 62(2): 69--97. \href{https://doi.org/10.2469/faj.v62.n2.4084}{https://doi.org/10.2469/faj.v62.n2.4084}.
\bibitem{zhangzhang2008} Zhang, H. \& Zhang, Q. (2008). \emph{Trading a Mean-Reverting Asset: Buy Low and Sell High}. Automatica 44(6): 1511--1518. \href{https://doi.org/10.1016/j.automatica.2007.11.003}{https://doi.org/10.1016/j.automatica.2007.11.003}.
\bibitem{stein2009} Stein, J. C. (2009). \emph{Presidential Address: Sophisticated Investors and Market Efficiency}. Journal of Finance 64(4): 1517--1548. \href{https://doi.org/10.1111/j.1540-6261.2009.01472.x}{https://doi.org/10.1111/j.1540-6261.2009.01472.x}.
\bibitem{brunnermeier2009} Brunnermeier, M. K. \& Pedersen, L. H. (2009). \emph{Market Liquidity and Funding Liquidity}. Review of Financial Studies 22(6): 2201--2238. \href{https://doi.org/10.1093/rfs/hhn098}{https://doi.org/10.1093/rfs/hhn098}.
\bibitem{carrwu2009} Carr, P. \& Wu, L. (2009). \emph{Variance Risk Premiums}. Review of Financial Studies 22(3): 1311--1341. \href{https://doi.org/10.1093/rfs/hhn038}{https://doi.org/10.1093/rfs/hhn038}.
\bibitem{bollerslev2009} Bollerslev, T., Tauchen, G. \& Zhou, H. (2009). \emph{Expected Stock Returns and Variance Risk Premia}. Review of Financial Studies 22(11): 4463--4492. \href{https://doi.org/10.1093/rfs/hhp008}{https://doi.org/10.1093/rfs/hhp008}.
\bibitem{avellaneda2010} Avellaneda, M. \& Lee, J.-H. (2010). \emph{Statistical Arbitrage in the US Equities Market}. Quantitative Finance 10(7): 761--782. \href{https://doi.org/10.1080/14697680903124632}{https://doi.org/10.1080/14697680903124632}.
\bibitem{grombvayanos2010} Gromb, D. \& Vayanos, D. (2010). \emph{Limits of Arbitrage}. Annual Review of Financial Economics 2: 251--275. \href{https://doi.org/10.1146/annurev-financial-073009-104107}{https://doi.org/10.1146/annurev-financial-073009-104107}.
\bibitem{martin2011} Martin, R. \& Schöneborn, T. (2011). \emph{Mean Reversion Pays, but Costs}. Risk 24(2): 84--89；arXiv:\allowbreak 1103.4934. \href{https://arxiv.org/abs/1103.4934}{https://arxiv.org/abs/1103.4934}.
\bibitem{mop2012} Moskowitz, T. J., Ooi, Y. H. \& Pedersen, L. H. (2012). \emph{Time Series Momentum}. Journal of Financial Economics 104(2): 228--250. \href{https://doi.org/10.1016/j.jfineco.2011.11.003}{https://doi.org/10.1016/j.jfineco.2011.11.003}.
\bibitem{asness2013} Asness, C. S., Moskowitz, T. J. \& Pedersen, L. H. (2013). \emph{Value and Momentum Everywhere}. Journal of Finance 68(3): 929--985. \href{https://doi.org/10.1111/jofi.12021}{https://doi.org/10.1111/jofi.12021}.
\bibitem{dofaff2012} Do, B. \& Faff, R. (2012). \emph{Are Pairs Trading Profits Robust to Trading Costs?}. Journal of Financial Research 35(2): 261--287. \href{https://doi.org/10.1111/j.1475-6803.2012.01317.x}{https://doi.org/10.1111/j.1475-6803.2012.01317.x}.
\bibitem{khandanilo2011} Khandani, A. E. \& Lo, A. W. (2011). \emph{What Happened to the Quants in August 2007? Evidence from Factors and Transactions Data}. Journal of Financial Markets 14(1): 1--46. \href{https://doi.org/10.1016/j.finmar.2010.07.005}{https://doi.org/10.1016/j.finmar.2010.07.005}.
\bibitem{menkhoff2012} Menkhoff, L., Sarno, L., Schmeling, M. \& Schrimpf, A. (2012). \emph{Currency Momentum Strategies}. Journal of Financial Economics 106(3): 660--684. \href{https://doi.org/10.1016/j.jfineco.2012.06.009}{https://doi.org/10.1016/j.jfineco.2012.06.009}.
\bibitem{frazzini2018} Frazzini, A., Israel, R. \& Moskowitz, T. J. (2018). \emph{Trading Costs of Asset Pricing Anomalies}. SSRN / 工作论文. \href{https://doi.org/10.2139/ssrn.2294498}{https://doi.org/10.2139/ssrn.2294498}.
\bibitem{baltaskosowski2011} Baltas, N. \& Kosowski, R. (2011). \emph{Momentum Strategies in Futures Markets and Trend-Following Funds}. SSRN 工作论文. \href{https://doi.org/10.2139/ssrn.1968996}{https://doi.org/10.2139/ssrn.1968996}.
\bibitem{baltaskosowski2013imp} Baltas, N. \& Kosowski, R. (2013). \emph{Improving Time-Series Momentum Strategies: The Role of Trading Signals and Volatility Estimators}. SSRN 工作论文. \href{https://doi.org/10.2139/ssrn.2140091}{https://doi.org/10.2139/ssrn.2140091}.
\bibitem{cuturi2013} Cuturi, M. \& d'Aspremont, A. (2013). \emph{Mean Reversion with a Variance Threshold}. ICML / PMLR 28(3): 271--279. \href{https://proceedings.mlr.press/v28/cuturi13.html}{https://proceedings.mlr.press/v28/cuturi13.html}.
\bibitem{ilmanen2011} Ilmanen, A. (2011). \emph{Expected Returns: An Investor's Guide to Harvesting Market Rewards}. John Wiley \& Sons. \href{https://doi.org/10.1002/9781118467190}{https://doi.org/10.1002/9781118467190}.
\bibitem{angtimmermann2012} Ang, A. \& Timmermann, A. (2012). \emph{Regime Changes and Financial Markets}. Annual Review of Financial Economics 4: 313--337. \href{https://doi.org/10.1146/annurev-financial-110311-101808}{https://doi.org/10.1146/annurev-financial-110311-101808}.
\bibitem{hamilton1989} Hamilton, J. D. (1989). \emph{A New Approach to the Economic Analysis of Nonstationary Time Series and the Business Cycle}. Econometrica 57(2): 357--384. \href{https://doi.org/10.2307/1912559}{https://doi.org/10.2307/1912559}.
\bibitem{garleanu2013} Gârleanu, N. \& Pedersen, L. H. (2013). \emph{Dynamic Trading with Predictable Returns and Transaction Costs}. Journal of Finance 68(6): 2309--2340. \href{https://doi.org/10.1111/jofi.12080}{https://doi.org/10.1111/jofi.12080}.
\bibitem{krauss2017} Krauss, C. (2017). \emph{Statistical Arbitrage Pairs Trading Strategies: Review and Outlook}. Journal of Economic Surveys 31(2): 513--545. \href{https://doi.org/10.1111/joes.12153}{https://doi.org/10.1111/joes.12153}.
\bibitem{barroso2015} Barroso, P. \& Santa-Clara, P. (2015). \emph{Momentum Has Its Moments}. Journal of Financial Economics 116(1): 111--120. \href{https://doi.org/10.1016/j.jfineco.2014.11.010}{https://doi.org/10.1016/j.jfineco.2014.11.010}.
\bibitem{leungli2015} Leung, T. \& Li, X. (2015). \emph{Optimal Mean Reversion Trading with Transaction Costs and Stop-Loss Exit}. International Journal of Theoretical and Applied Finance 18(3): 1550020. \href{https://doi.org/10.1142/S021902491550020X}{https://doi.org/10.1142/S021902491550020X}.
\bibitem{harveysliu2015} Harvey, C. R. \& Liu, Y. (2015). \emph{Backtesting}. Journal of Portfolio Management 42(1): 13--28. \href{https://doi.org/10.3905/jpm.2015.42.1.013}{https://doi.org/10.3905/jpm.2015.42.1.013}.
\bibitem{danielmoskowitz2016} Daniel, K. \& Moskowitz, T. J. (2016). \emph{Momentum Crashes}. Journal of Financial Economics 122(2): 221--247. \href{https://doi.org/10.1016/j.jfineco.2015.12.002}{https://doi.org/10.1016/j.jfineco.2015.12.002}.
\bibitem{hlz2016} Harvey, C. R., Liu, Y. \& Zhu, H. (2016). \emph{…and the Cross-Section of Expected Returns}. Review of Financial Studies 29(1): 5--68. \href{https://doi.org/10.1093/rfs/hhv059}{https://doi.org/10.1093/rfs/hhv059}.
\bibitem{mcleanpontiff2016} McLean, R. D. \& Pontiff, J. (2016). \emph{Does Academic Research Destroy Stock Return Predictability?}. Journal of Finance 71(1): 5--32. \href{https://doi.org/10.1111/jofi.12365}{https://doi.org/10.1111/jofi.12365}.
\bibitem{kim2016} Kim, A. Y., Tse, Y. \& Wald, J. K. (2016). \emph{Time Series Momentum and Volatility Scaling}. Journal of Financial Markets 30: 103--124. \href{https://doi.org/10.1016/j.finmar.2016.05.003}{https://doi.org/10.1016/j.finmar.2016.05.003}.
\bibitem{bailey2016} Bailey, D. H., Borwein, J. M., López de Prado, M. \& Zhu, Q. J. (2016). \emph{The Probability of Backtest Overfitting}. The Journal of Computational Finance 20(4): 39--69. \href{https://doi.org/10.21314/jcf.2016.322}{https://doi.org/10.21314/jcf.2016.322}.
\bibitem{novymarxvelikov2016} Novy-Marx, R. \& Velikov, M. (2016). \emph{A Taxonomy of Anomalies and Their Trading Costs}. Review of Financial Studies 29(1): 104--147. \href{https://doi.org/10.1093/rfs/hhv063}{https://doi.org/10.1093/rfs/hhv063}.
\bibitem{hop2017} Hurst, B., Ooi, Y. H. \& Pedersen, L. H. (2017). \emph{A Century of Evidence on Trend-Following Investing}. Journal of Portfolio Management 44(1): 15--29. \href{https://doi.org/10.3905/jpm.2017.44.1.015}{https://doi.org/10.3905/jpm.2017.44.1.015}.
\bibitem{georgopoulou2017} Georgopoulou, A. \& Wang, J. (2017). \emph{The Trend Is Your Friend: Time-Series Momentum Strategies across Equity and Commodity Markets}. Review of Finance 21(4): 1557--1592. \href{https://doi.org/10.1093/rof/rfw048}{https://doi.org/10.1093/rof/rfw048}.
\bibitem{harvey2018} Harvey, C. R., Hoyle, E., Korgaonkar, R., Rattray, S., Sargaison, M. \& van Hemert, O. (2018). \emph{The Impact of Volatility Targeting}. Journal of Portfolio Management 45(1): 14--33. \href{https://doi.org/10.3905/jpm.2018.45.1.014}{https://doi.org/10.3905/jpm.2018.45.1.014}.
\bibitem{moreiramuir2017} Moreira, A. \& Muir, T. (2017). \emph{Volatility-Managed Portfolios}. Journal of Finance 72(4): 1611--1644. \href{https://doi.org/10.1111/jofi.12513}{https://doi.org/10.1111/jofi.12513}.
\bibitem{bailey2014} Bailey, D. H. \& López de Prado, M. (2014). \emph{The Deflated Sharpe Ratio}. Journal of Portfolio Management 40(5): 94--107. \href{https://doi.org/10.3905/jpm.2014.40.5.094}{https://doi.org/10.3905/jpm.2014.40.5.094}.
\bibitem{huang2020} Huang, D., Li, J., Wang, L. \& Zhou, G. (2020). \emph{Time Series Momentum: Is It There?}. Journal of Financial Economics 135(3): 774--794. \href{https://doi.org/10.1016/j.jfineco.2019.08.004}{https://doi.org/10.1016/j.jfineco.2019.08.004}.
\bibitem{gu2020} Gu, S., Kelly, B. \& Xiu, D. (2020). \emph{Empirical Asset Pricing via Machine Learning}. Review of Financial Studies 33(5): 2223--2283. \href{https://doi.org/10.1093/rfs/hhaa009}{https://doi.org/10.1093/rfs/hhaa009}.
\bibitem{baltas2019} Baltas, N. (2019). \emph{The Impact of Crowding in Alternative Risk Premia Investing}. Financial Analysts Journal 75(3): 89--104. \href{https://doi.org/10.1080/0015198x.2019.1600955}{https://doi.org/10.1080/0015198x.2019.1600955}.
\bibitem{baltaskosowski2020} Baltas, N. \& Kosowski, R. (2020). \emph{Demystifying Time-Series Momentum Strategies}. Wiley book chapter. \href{https://doi.org/10.1002/9781119599364.ch3}{https://doi.org/10.1002/9781119599364.ch3}.
\bibitem{cochrane2011} Cochrane, J. H. (2011). \emph{Presidential Address: Discount Rates}. Journal of Finance 66(4): 1047--1108. \href{https://doi.org/10.1111/j.1540-6261.2011.01671.x}{https://doi.org/10.1111/j.1540-6261.2011.01671.x}.
\bibitem{wood2021} Wood, K., Roberts, S. \& Zohren, S. (2021). \emph{Slow Momentum with Fast Reversion: A Trading Strategy Using Deep Learning and Changepoint Detection}. Journal of Financial Data Science；arXiv:\allowbreak 2105.13727. \href{https://arxiv.org/abs/2105.13727}{https://arxiv.org/abs/2105.13727}.
\bibitem{ehsani2022} Ehsani, S. \& Linnainmaa, J. T. (2022). \emph{Factor Momentum and the Momentum Factor}. Journal of Finance 77(3): 1877--1919. \href{https://doi.org/10.1111/jofi.13131}{https://doi.org/10.1111/jofi.13131}.
\bibitem{loupolk2022} Lou, D. \& Polk, C. (2022). \emph{Comomentum: Inferring Arbitrage Activity from Return Correlations}. Review of Financial Studies 35(2): 735--776. \href{https://doi.org/10.1093/rfs/hhab117}{https://doi.org/10.1093/rfs/hhab117}.
\bibitem{pitkajarvi2020} Pitkäjärvi, A., Suominen, M. \& Vaittinen, J. (2020). \emph{Cross-Asset Signals and Time Series Momentum}. Journal of Financial Economics 136(2): 459--478. \href{https://doi.org/10.1016/j.jfineco.2019.02.011}{https://doi.org/10.1016/j.jfineco.2019.02.011}.
\bibitem{guijarro2021} Guijarro-Ordonez, J., Pelger, M. \& Zanotti, G. (2021). \emph{Deep Learning Statistical Arbitrage}. arXiv:\allowbreak 2106.04028. \href{https://arxiv.org/abs/2106.04028}{https://arxiv.org/abs/2106.04028}.
\bibitem{ning2024} Ning, B. \& Lee, A. (2024). \emph{Advanced Statistical Arbitrage with Reinforcement Learning}. arXiv:\allowbreak 2403.12180. \href{https://arxiv.org/abs/2403.12180}{https://arxiv.org/abs/2403.12180}.
\bibitem{suchato2024} Suchato, P., Wiryadi, A., Chen, Y., Zhao, A. \& Yue, X. (2024). \emph{An Application of the Ornstein--Uhlenbeck Process to Pairs Trading}. arXiv:\allowbreak 2412.12458. \href{https://arxiv.org/abs/2412.12458}{https://arxiv.org/abs/2412.12458}.
\bibitem{famablume1966} Fama, E. F. \& Blume, M. E. (1966). \emph{Filter Rules and Stock-Market Trading}. Journal of Business 39(1): 226--241. \href{https://doi.org/10.1086/294849}{https://doi.org/10.1086/294849}.
\bibitem{lo2000} Lo, A. W., Mamaysky, H. \& Wang, J. (2000). \emph{Foundations of Technical Analysis: Computational Algorithms, Statistical Inference, and Empirical Implementation}. Journal of Finance 55(4): 1705--1765. \href{https://doi.org/10.1111/0022-1082.00265}{https://doi.org/10.1111/0022-1082.00265}.
\bibitem{novymarx2012} Novy-Marx, R. (2012). \emph{Is Momentum Really Momentum?}. Journal of Financial Economics 104(2): 363--390. \href{https://doi.org/10.1016/j.jfineco.2011.05.003}{https://doi.org/10.1016/j.jfineco.2011.05.003}.
\bibitem{chabot2014} Chabot, B., Ghysels, E. \& Jagannathan, R. (2014). \emph{Momentum Trading, Return Chasing, and Predictable Crashes}. NBER WP 20660. \href{https://doi.org/10.3386/w20660}{https://doi.org/10.3386/w20660}.
\end{thebibliography}


\end{document}
